% DER_II-a: Grothendieck's text with the commentary set in, then the overview, inside a group that applies the chunk's own definitions
\dgrchapter{II}{Chapter II}{Cofinalité (à droite et à gauche) (préliminaire à la cofinalité cohomologique)}
\dgrpart{0}{II-a}{Part 1 of 2, manuscript pages 1--53}{II-a\enspace\textperiodcentered\enspace Chapter II, part 1 of 2\enspace\textperiodcentered\enspace manuscript pages 1--53}{Chapter II \textperiodcentered\ \textup{II-a}}
\dgrchunkbegin
% DER_II-a: the definitions of the chunk's preamble, applied inside its group
\let\printedpage\relax\newcommand{\printedpage}[1]{\begingroup\def\dcpagevalue{#1}\def\dcblankpage{--}\ifx\dcpagevalue\dcblankpage\else\marginpar{\scriptsize\raggedright p.~#1}\fi\endgroup}
\def\dcmsnumber page #1\relax{#1}
\newif\ifdcskipmsp
\let\dcpairedpage\relax\newcommand{\dcpairedpage}[2]{\marginpar{\scriptsize\raggedright p.~#1\\\mbox{[#2]}}\global\dcskipmsptrue}
\let\msp\relax\newcommand{\msp}[1]{\ifdcskipmsp\global\dcskipmspfalse\else\marginpar{\scriptsize\raggedright[\dcmsnumber#1\relax]}\fi}
% printed notation: the category C' enclosed in a V-shaped wedge with the arrow name under it
\let\Vcat\relax\newcommand{\Vcat}[2]{\underset{#2}{\underset{\vee}{#1}}}
\let\Homint\relax\newcommand{\Homint}{\operatorname{Hom\,int}}
\let\sX\relax\newcommand{\sX}{\mathcal{X}} % resolved against PDF: CMSY calligraphic X
\let\X\relax\newcommand{\X}{\mathcal{X}}
\let\Hom\relax\newcommand{\Hom}{\operatorname{Hom}}
\let\Ker\relax\newcommand{\Ker}{\operatorname{Ker}}
\let\Ob\relax\newcommand{\Ob}{\operatorname{Ob}}
\let\mspin\relax\newcommand{\mspin}[1]{\texttt{[#1]}}
\let\HOT\relax\newcommand{\HOT}{\mathrm{HOT}}
\let\Cat\relax\newcommand{\Cat}{\mathrm{Cat}}
\let\Fl\relax\newcommand{\Fl}{\mathrm{Fl}}
\let\Ens\relax\newcommand{\Ens}{\mathrm{Ens}}
\let\pl\relax\newcommand{\pl}{\texttt{[plutôt}}
\let\uHom\relax\newcommand{\uHom}{\underline{\operatorname{Hom}}}
\let\bs\relax\newcommand{\bs}{\backslash}
\let\cM\relax\newcommand{\cM}{\mathcal{M}}
\let\dbl\relax\newcommand{\dbl}[2]{\mathrel{\overset{#1}{\underset{#2}{\rightrightarrows}}}}
\let\defeq\relax\newcommand{\defeq}{\overset{\text{déf}}{=}}

% ---- from Der-ChII-234_p001-005.tex
% ===================== PDF page 1 =====================
\printedpage{--}
% no page number printed (title page)
% head: (none)

\begin{center}
\vspace*{4cm}
{\Large\textsf{LES DÉRIVATEURS}}

\bigskip\bigskip

\textbf{ALEXANDRE GROTHENDIECK}

\bigskip\bigskip\bigskip

\textbf{Chapitre II}

\textbf{Cofinalité (à droite et à gauche)}

\textbf{(préliminaire à la cofinalité cohomologique)}
\end{center}

\clearpage

% ===================== PDF page 2 =====================
\printedpage{--}
% no page number printed (editorial note)
% head: (none)

\vspace*{6cm}

\noindent Ce texte a été déchiffré et transcrit en \LaTeXe{} par M. Künzer. Il a été édité par M. Künzer, J. Malgoire, et G. Maltsiniotis. La transcription est aussi fidèle que possible au manuscrit. Pour les quelques corrections évidentes, ou rares commentaires des éditeurs, ainsi que pour la numérotation originale des pages du manuscrit, les caractères de machine à écrire \texttt{[typewriter]} entre crochets sont utilisés. Un point d'interrogation entre crochets signifie que l'on n'est pas sûr du mot qui précède.

\noindent Cette édition est provisoire. Les remarques, commentaires et corrections sont bienvenus. Envoyer un message à :

\noindent\textbf{[e-mail address removed]}

\hfill G. Maltsiniotis\hspace*{2cm}

\clearpage

% ===================== PDF page 3 =====================
\printedpage{1}
% head: A. Grothendieck : Dérivateurs, Ch. II, Édition des Universités de Montpellier II et Paris VII
% foot: Version 8 juillet 2001

\msp{page 1}

\section*{1\quad Cofinalité cohomologique}

Soit
\[ \alpha : C' \longrightarrow C \]
un foncteur de (petites) catégories, $\mathcal{M}$ une catégorie ayant les $\varprojlim$ pertinentes, $F$ un préfaisceau (donc contravariant en $i$) sur $C$ à valeurs dans $\mathcal{M}$. On a un homomorphisme canonique dans $\mathcal{M}$
\[ \varprojlim_{C} F \longrightarrow \varprojlim_{C'} \underbrace{\alpha^*(F)}_{F\circ\alpha} \tag*{(1)} \]
et on veut des conditions sur $\alpha$ pour que ce soit toujours un isomorphisme (\footnote{On dit alors que $\alpha$ est \textit{cofinal,} cf. SGA 4, I 8.1, ou 3.3, p. 45.}). On est ramené aussitôt au cas où $\mathcal{M} = (\mathrm{Ens})$. Disons qu'une catégorie $I$ est \textit{faiblement filtrante à gauche} si $I \neq \emptyset$ et si $\forall\, i', i'' \in \operatorname{Ob} I$, $\exists\, i \in \operatorname{Ob} I$ et des flèches
\[
\begin{tikzcd}[row sep=small, column sep=small]
 & i' \\
i \arrow[ur] \arrow[dr] & \\
 & i''.
\end{tikzcd}
\]
Cela signifie aussi que $\operatorname{Ob} I$ en tant qu'ensemble ordonné est filtrant à gauche.

\dgrblock{dgrb78}{The presheaf represented by a change of index}
\begin{dcomment}{exposition}{p1:On est ramené aussitôt au cas où}{The presheaf represented by a change of index}
Grothendieck begins again with ordinary limits because their elementary descent argument already contains the object which the cohomological theory must derive. We keep his variance throughout: a presheaf on $C$ is a functor $C^{\mathrm{op}}\to\mathrm{Set}$, and $c/\alpha$, also written $c\backslash C'$, has objects $(a,f:c\to\alpha a)$. The arrow in the first proposition is printed in the opposite direction; its intended direction is fixed both by this notation and by the proof on manuscript page~6. This is an edition slip, not a change of convention.

Put $K_\alpha(c)=\pi_0(c/\alpha)$. Precomposition with $c_0\to c_1$ defines $K_\alpha(c_1)\to K_\alpha(c_0)$, so $K_\alpha$ is a presheaf. There is a natural bijection
\[
 \lim_{C'^{\mathrm{op}}}\alpha^*F
 \;\cong\;\operatorname{Nat}(K_\alpha,F).
\]
Indeed a compatible family $(s_a)$ determines the map $[a,f]\mapsto F(f)(s_a)$. A morphism $(a,f)\to(b,g)$ in $c/\alpha$ gives $g=\alpha(h)f$, and compatibility gives $s_a=F(\alpha h)(s_b)$; hence the value is constant along every zigzag. Naturality in $c$ follows by composition. Conversely, from $t:K_\alpha\to F$, put $s_a=t_{\alpha a}[a,1_{\alpha a}]$. For $h:a\to b$, the arrow $(a,1)\to(b,\alpha h)$ and naturality of $t$ prove compatibility. These constructions are inverse.

Equivalently, $K_\alpha=\alpha_!1_{C'}$. Thus the restriction map on limits is induced by the augmentation $K_\alpha\to1_C$. Yoneda, applied in the presheaf category, proves the exact criterion:
\[
 \text{restriction preserves every presheaf limit}
 \quad\Longleftrightarrow\quad
 \pi_0(c/\alpha)=\{*\}\quad(c\in C).
\]
The nonempty, connected comma categories are what matter. A common predecessor for each pair of objects is a stronger sufficient condition; it has no equalization requirement. For coefficients in a locally small category $\mathcal M$, apply $\operatorname{Hom}_{\mathcal M}(M,-)$ for every $M$. These functors preserve existing limits and jointly reflect isomorphisms, which justifies the reduction to sets without imposing extra completeness on $\mathcal M$.

In current covariant terminology, $\alpha$ is a final functor, while $\alpha^{\mathrm{op}}$ is initial; the latter is the functor acting on the diagrams whose limits are written here. SGA~4 uses the cofinal terminology followed in the manuscript \citep[Expos\'e~I, \S8.1]{SGA4I1972}. Proposition~3.2 below returns to precisely this criterion. We shall use the representing presheaf again there, rather than repeat the elementwise argument.
\end{dcomment}


\textbf{Proposition.} \textit{Pour que} (1) \textit{soit toujours un isomorphisme} ($\alpha$ \textit{`anodin pour} $\varprojlim$ \textit{'), il suffit que pour tout} $x_0 \in \operatorname{Ob} C$\textit{, la catégorie} $x_0\backslash C'$ \textit{des couples} $(x', u)$\textit{, avec} $x' \in \operatorname{Ob} C'$\textit{,} $u : \alpha(x') \longrightarrow x$\textit{, soit faiblement filtrante à gauche} (\footnote{Cf. amélioration des remarques p. 5.}, \footnote{Cf. 3.2, p. 43-44, pour conditions nécessaires et suffisantes de cofinalité.}). % [sic] u : \alpha(x') \longrightarrow x -> x_0 in place of x, the category is x_0\backslash C'

Ceci n'implique pas la conclusion plus forte
\[ F \underbrace{\xrightarrow{\;\sim\;}}_{\text{adj.}} \alpha_*\alpha^* F \]
que le morphisme d'adjonction canonique soit un isomorphisme. Pour qu'il soit ainsi, il suffit que $\forall\, x_0 \in C$, le foncteur $\alpha/x_0 : C'/x_0 \longrightarrow C/x_0$ induit par $\alpha$ soit `anodin pour les $\varprojlim$ ', ce qui signifie aussi que pour toute flèche $u : x_1 \longrightarrow x_0$ dans $C$, la catégorie $\Vcat{C'}{u}$ formée des triples $(x', \xi, \eta)$ avec $x' \in \operatorname{Ob} C'$ et $\xi : x_1 \longrightarrow \alpha(x')$, $\eta : \alpha(x') \longrightarrow x_0$ rendant commutatif
% printed: C' inside a V-shaped wedge, with u printed below the wedge
\msp{page 2}
le diagramme
\printedpage{2}
% head: A. Grothendieck : Dérivateurs, Ch. II, Édition des Universités de Montpellier II et Paris VII
% foot: Version 8 juillet 2001
\[
\begin{tikzcd}
x_1 \arrow[rr, "u"] \arrow[dr, "\xi"'] & & x_0 \\
 & \alpha(x'), \arrow[ur, "\eta"'] &
\end{tikzcd}
\]
est faiblement filtrante à gauche.

Plus généralement, considérons
\[
\begin{tikzcd}
X' = C' \arrow[rr, "\alpha"] \arrow[dr, "f'"'] & & C = X \arrow[dl, "f"] \\
 & S &
\end{tikzcd}
\]
et posons nous la question quand l'homomorphisme fonctoriel, par rapport à un préfaisceau $F$ variable sur $X$,
\[ f_*(F) \longrightarrow f'_*(\alpha^*(F)) \]
est toujours un isomorphisme. (Le cas précédent est celui où $S = X$, $f = \mathrm{id}_X$.) On trouve qu'il suffit que $\forall\, s_0 \in \operatorname{Ob} S$, le foncteur
\[ \alpha/s_0 : X'/s_0 \longrightarrow X/s_0 \]
soit `anodin pour les $\varprojlim$ ', et pour ceci, il suffit que pour \texttt{[tout]} objet $(x_0, u_0)$ de $X/s_0$ (où
\[ f(x_0) \xrightarrow{u_0} s_0) \]
la catégorie $(x_0, u_0)\backslash(X'/s_0)$ soit faiblement filtrante à gauche.

\msp{page 3}

Cette catégorie s'interprète comme celle des triples $(x', \xi, u)$ avec
\[
\left\{
\begin{aligned}
& x' \in \operatorname{Ob} C' \\
& u : g(x') = f(\alpha(x')) \longrightarrow s_0 \\
& \xi : x_0 \longrightarrow \alpha(x')
\end{aligned}
\right.
\]
\texttt{[plutôt} $f'(x')$ \texttt{au lieu de} $g(x')$\texttt{]} de façon à rendre commutatif le diagramme
\[
\begin{tikzcd}
f(x_0) \arrow[rr, "u_0"] \arrow[dr, "f(\xi)"'] & & s_0 \\
 & f(\alpha(x')). \arrow[ur, "u"'] &
\end{tikzcd}
\]
Je voudrais dégager des variantes de ces énoncés dans le cadre cohomologique, les foncteurs $\varprojlim$ étant remplacés par des $R\Gamma$ (de préfaisceaux abéliens p.ex. - mais on pense aussi aux `cas non commutatifs', définis par des dérivateurs notamment.)

\clearpage

% ===================== PDF page 5 =====================
\printedpage{3}
% head: A. Grothendieck : Dérivateurs, Ch. II, Édition des Universités de Montpellier II et Paris VII
% foot: Version 8 juillet 2001

Si $I$ est une catégorie \textit{ordonnée finie}, la condition \texttt{[faiblement]} filtrante à gauche équivaut à l'existence d'un objet initial, ce qui implique que $I$ est contractile. Dans les énoncés précedents, il se pourrait que la bonne philosophie soit de remplacer partout `faiblement filtrante à gauche' par `contractile'. Si p.ex. $C$ est la catégorie finale, alors remplacer la condition que $C'$ soit filtrante à gauche, par l'acyclicité. C'est en tous cas en accord avec l'esprit de l'axiome $\mathrm{Der}_5$ (cas où $J$ est la catégorie finale). En effet, dans le cas où $C = e$, la condition que $\alpha$ soit anodin pour les $\varprojlim$ équivaut visiblement à dire que le foncteur $\alpha^* : \mathcal{M} \longrightarrow \underline{\mathrm{Hom}}(C'^{o}, \mathcal{M})$, associant à tout $F \in \mathcal{M}$ le préfaisceau constant de valeur $F$ % [sic] précedents -> précédents, missing accent on the second e
\msp{page 4}
sur $C'$, donne un morphisme d'adjonction
\[ \varprojlim_{e} F = F \longrightarrow \alpha_*(\alpha^*(F)) = \varprojlim_{C'} \alpha^* F \]
qui est un isomorphisme. Et l'axiome $\mathrm{Der}_5$ assure qu'il doit en être ainsi si $C'$ est contractile.

\dgrblock{dgrb79}{Why the relative problem has two endpoints}
\begin{dcomment}{exposition}{p3}{Why the relative problem has two endpoints}
The passage from global sections to a direct image introduces a second endpoint. For $f:X\to S$ and a presheaf $F$, the pointwise right Kan extension formula is
\[
 (f_*F)(s)=\lim_{(f/s)^{\mathrm{op}}}F(x),
 \qquad (f/s)=\{(x,f(x)\to s)\}.
\]
This is the ordinary comma-category formula for a Kan extension \citep[Theorem~6.2.1 and Corollary~6.2.9]{Riehl2016}. If $f'=f\alpha$, the comparison at $s$ is restriction along $\alpha/s$. Its comma category at $(x_0,u_0)$ consists exactly of factorizations
\[
 f(x_0)\xrightarrow{\,f(\xi)\,}f(\alpha a)
 \xrightarrow{\,u\,}s,
 \qquad u f(\xi)=u_0.
\]
Connectedness and nonemptiness of these categories therefore suffice. The weak cofiltering condition written here also suffices, but the phrase ``ce qui signifie aussi'' on manuscript page~1 should not turn that sufficient condition into a necessary one.

When $S=X$ and $f=1_X$, the test concerns factorizations of each arrow $x_0\to x_1$ through $\alpha a$. It is stronger than the test involving only $x_0/\alpha$. This is the same distinction already isolated in the commentary on Chapter~I, manuscript p.~49 (``The factorization category controls base change''); Proposition~4.1 will identify its ordinary categorical form. Keeping the two endpoints visible prevents an erroneous inference from invariance of global sections to full faithfulness of restriction.
\end{dcomment}


La condition filtrante dans la proposition est d'ailleurs une condition \textit{suffisante} d'anodinité, mais nullement nécessaire. Si p.ex. $C$ est la catégorie finale, alors la condition d'anodinité pour $C' \longrightarrow C$ est simplement que $C'$ soit connexe non vide, i.e. $\pi_0(C') = \{0\}$. (Condition qui visiblement néglige les $\pi_i$ supérieurs.)

Soit \fbox{$\alpha : X' \longrightarrow X$} un foncteur de petites catégories. La condition suggérée par la proposition, que les $x_0\backslash X'$ soient contractiles, signifie que $\alpha$ est une hot-équivalence co-stricte, i.e. que $\alpha^{o} : X'^{o} \longrightarrow X^{o}$ est une hot-équivalence `localement sur $X^{o}$', donc aussi globalement. Si $\alpha$
\msp{page 5}
lui-même était une hot-équivalence stricte, au lieu de $\alpha^{o}$, alors on aurait par $\mathrm{Der}_5$, que pour $F \in \mathbf{D}(X)$, $F \xrightarrow{\;\sim\;} \alpha_*\alpha^* F$ ($\mathbf{D}$ dérivateur à gauche), ce qui est bien plus fort que $\Gamma_{\mathbf{D}}(X, F) \xrightarrow[\sim]{\alpha^*} \Gamma_{\mathbf{D}}(X', \alpha^* F)$. La condition duale implique (si $\mathbf{D}$ \texttt{[est un]} dérivateur à droite) $\alpha_!\alpha^* F \xrightarrow{\;\sim\;} F$. Mais est-il raisonnable de s'attendre qu'elle implique aussi la condition `du côté cohomologie, au lieu de l'homologie',
\[ \Gamma_{\mathbf{D}}(X, F) \xrightarrow{\;\sim\;} \Gamma_{\mathbf{D}}(X', \alpha^*(F)) \ \ ? \]

Le premier cas, qui devrait emporter conviction, est celui du dérivateur DerAb, i.e. des catégories dérivées des faisceaux abéliens :

\textbf{Question.} \textit{Soit} $\alpha : X' \longrightarrow X$ \textit{une} hot-\textit{équivalence co-stricte,} $F$ \textit{un préfaisceau abélien sur} $X$\textit{. A-t-on alors} $H^i(X, F) \xrightarrow[\sim]{\alpha^*} H^i(X', \alpha^*(F))$ \textit{pour tout} $i \in \mathbf{Z}$ \textit{?}

Si $X$ admet un objet initial $\emptyset$, alors $\emptyset\backslash X' \simeq X'$, donc $X'$ contractile, donc c'est OK pour $F$ constant.

\dgrblock{dgrb80}{Cohomological cofinality and the derived weight}
\begin{dcomment}{later}{p5:donc c'est OK pour $F$ constant.}{Cohomological cofinality and the derived weight}
Grothendieck's question about abelian coefficients has an affirmative answer. Its proof can follow the argument he has just discovered, provided that the presheaf $K_\alpha$ is replaced by its chain-level predecessor. We give the argument for a commutative ring $k$; $k=\mathbb Z$ answers the stated question. Write $D(C;k)$ for the derived category of $k$-module presheaves and use $\mathbf R\operatorname{Hom}_{C;k}$ for the \emph{global} derived Hom, valued in $D(k)$.

Derived adjunction gives, naturally in $F$,
\[
 \mathbf R\Gamma(C',\alpha^*F)
 \simeq
 \mathbf R\operatorname{Hom}_{C;k}
       (\mathbf L\alpha_!k_{C'},F).
\]
For clarity, this can be constructed without assuming the disputed internal formula on manuscript page~10. Resolve $k_{C'}$ by the augmented simplicial presheaf of free modules whose degree $n$ is
\[
 P_n=\bigoplus_{a_0\to\cdots\to a_n} k\operatorname{Hom}_{C'}(-,a_0).
\]
At $a$, its augmented chain complex is the nerve complex of $a/C'$, contracted by its initial object $1_a$. Each $P_n$ is projective. Its normalized complex, in nonpositive cohomological degrees, is therefore a projective resolution. Since restriction is exact, $\alpha_!$ takes these projectives to projectives. Evaluating $\alpha_!P_\bullet$ at $c$ gives the free chain complex of $c/\alpha$. In particular
\[
 (\mathbf L\alpha_!k_{C'})(c)
 \simeq C_*(N(c/\alpha);k),
\]
where a chain degree $n$ is placed in cohomological degree $-n$. The usual formulation using the opposite comma category gives the same complex up to signed reversal.

If every $N(c/\alpha)$ is weakly contractible, its augmented complex is quasi-isomorphic to $k$. Consequently $\mathbf L\alpha_!k_{C'}\to k_C$ is an isomorphism. Applying derived Hom yields the desired comparison in every degree, and indeed for every complex $F$. This proof needs only $k$-acyclicity of these nerves; weak contractibility is the uniform condition which works for homotopy types as well. The two conditions must not be confused with mere connectedness. For example, $BG\to e$ is cofinal for ordinary presheaf limits whenever $G$ is a group, but for $G=\mathbb Z$ and the constant coefficient $\mathbb Z$ one has $H^1(BG;\mathbb Z)=\mathbb Z$.

For spaces, the corresponding weight is $c\mapsto B(c/\alpha)$ in $\operatorname{Fun}(C^{\mathrm{op}},\mathcal S)$. The same mapping-object adjunction proves cofinality; the general higher categorical criterion is \citet[Theorem~4.1.3.1, together with Proposition~4.1.1.8]{Lurie2009HTT}, applied with the appropriate opposites. Its corollary~4.1.3.3 is Quillen's Theorem~A. These results settle the question on page~5, without requiring the stronger unit or counit assertions in its preceding paragraph. Those assertions are conditional there on the manuscript's invocation of $\mathrm{Der}_5$; interpreted as consequences of comma-category asphericity under present derivator axioms, they are false. The endpoint calculation in the next block makes the distinction explicit, continuing the caution about the changing axioms in the commentary on Chapter~I, manuscript p.~43 (``What the localization axioms do and do not say'').
\end{dcomment}


\textbf{Remarque.} Je m'aperçois que dans la proposition 1, au lieu de la condition que les $x\backslash C'$ soient faiblement filtrantes à gauche, il suffirait qu'elles soient non vides et connexes, i.e.

% ---- from Der-ChII-234_p006-010.tex
\printedpage{4}
% head: A. Grothendieck : Dérivateurs, Ch. II, Édition des Universités de Montpellier II et Paris VII
% foot: Version 8 juillet 2001
0-connexes (\footnote[4]{Cela signifie aussi que $\alpha$ est $W_0$-coasphérique. Pour un traitement ab ovo, cf. prop. 3.2, page 43.}). Le fait que les $x\backslash C'$ non vide implique que $\varprojlim\limits_{C} \longrightarrow \varprojlim\limits_{C'}$ soit \textit{injectif.} Pour la sujectivité, si $\xi' \in \varprojlim\limits_{C'}$ (cas ensembliste), on définit $\xi \in \varprojlim\limits_{C}$ par les $\xi(x) \in F(x)$ à l'aide de % [sic] sujectivité -> surjectivité, lettre r omise
\msp{page 6}
\[ x \xrightarrow{u} \alpha(x') \]
en prenant $u^*(\xi(x'))$ (où $\xi(x') \in (\alpha^*(F))(x') = F(\alpha(x')))$, le point essentiel est de montrer que les $u^*(\xi(x'))$ ne dépendent que de $x$, pas du choix de $(x',u) \in x\backslash C'$. Or pour $(x',u)$ variable dans $x\backslash C'$, l'application
\[
\begin{aligned}
\varphi : \operatorname{Ob}(x\backslash C') &\quad\longrightarrow\quad F(x)\\
(x',u) &\quad\longmapsto\quad u^*(\xi(x'))
\end{aligned}
\]
est telle que s'il existe une flèche
\[ (x',u) \xrightarrow{u'} (\bar{x}',\bar{u}) \qquad \text{où} \qquad u' : x' \longrightarrow \bar{x}' \]
(donc le diagramme
\[
\begin{tikzcd}
 & \alpha(x') \arrow[dd, "\alpha(u')"] \\
x \arrow[ur, "u"] \arrow[dr, "\bar{u}"'] & \\
 & \alpha(\bar{x}')
\end{tikzcd}
\]
est commutatif), alors $\varphi(x',u) = \varphi(\bar{x}',\bar{u})$. Donc $\varphi$ est constant sur les composantes connexes de $x\backslash C'$, donc constant si $x\backslash C'$ est 0-connexe, d'où $\xi(x)$. On prouve alors aussitôt que pour $x$ variable,
\[ \xi \overset{\text{déf}}{=} (\xi(x))_{x\in\operatorname{Ob} C} \simeq \varprojlim F, \qquad \text{\texttt{[plutôt}}\ \ \in \varprojlim F, \text{\texttt{]}} \]
et que $\alpha^*(\xi) = \xi'$, cqfd.

On peut dire que pour $x_0 \in \operatorname{Ob} C_0$ fixé, si on regarde le préfaisceau \textit{constant} sur $x_0\backslash C'$ de valeur $F(x_0)$, soit $\Phi_{x_0}$, l'application
\[ (x',u) \longmapsto \xi^*(u') : x_0\backslash C' \longrightarrow F(x_0) \]
peut être regardée comme une section de ce préfaisceau sur $x\backslash C'$, défini par $\xi' \in \varprojlim\limits_{C'} \alpha^*F$, donc on a
\[ \varprojlim_{C'} \alpha^*F \xrightarrow{c_{x_0}} \Gamma(x_0\backslash C', \Phi_{x_0}) \qquad (\xleftarrow{\ \sim\ } F(x_0) \text{ si } x_0\backslash C \text{ 0-connexe}). \]
Si on a \fbox{$x_0 \xrightarrow{v} x_1$} dans $C'$, on en déduit
\[ x_1\backslash C' \xrightarrow{V} x_0\backslash C' \]
(les $x\backslash C'$ sont les fibres d'une \textit{catégorie fibrée} $C\backslash C'$ sur $C$ \dots) et on a $F(x_1) \longrightarrow F(x_0)$, d'où \fbox{$\Phi_{x_1} \longrightarrow V^*\Phi_{x_0}$} pour les faisceaux constants correspondants sur $x_0\backslash C'$ resp. $x_1\backslash C'$
\printedpage{5}
% head: A. Grothendieck : Dérivateurs, Ch. II, Édition des Universités de Montpellier II et Paris VII
% foot: Version 8 juillet 2001
\msp{page 7}
et on a commutativité dans
\[
\begin{tikzcd}[column sep=large]
\varprojlim\limits_{C'} \alpha^*F \arrow[r, "c_{x_1}"] \arrow[d, "c_{x_0}"'] & \Gamma(x_0\backslash C', F(x_0)) \simeq \operatorname{Hom}(\pi_0(x_0\backslash C'), F(x_0)) \arrow[d, "V^*"] & \\
\operatorname{Hom}(\pi_0(x_1\backslash C'), F(x_1)) \simeq \Gamma(x_1\backslash C', F(x_1)) \arrow[r, "\substack{v^*\\ (F(x_1) \longrightarrow F(x_0))}"'] & \Gamma(x_1\backslash C', F(x_0)) \simeq \operatorname{Hom}(\pi_0(x_1\backslash C'), F(x_0)) & \\
\pi_0(x_1\backslash C') \arrow[r, "\pi_0(V)"] \arrow[d, dotted] & \pi_0(x_0\backslash C') \arrow[d, dotted] & \text{commutatif.} \\
F(x_1) \arrow[r, "F(v)"] & F(x_0) &
\end{tikzcd}
\]
On considère le préfaisceau d'ensembles
\[ x \longmapsto \pi_0(x\backslash C') \qquad \text{sur } C, \]
appelons-le $\pi_0(C'/C)$, on aura un morphisme canonique
\[ \boxed{\varprojlim_{C'} \alpha^*(F) \xrightarrow{\ \sim\ } \operatorname{Hom}_{\substack{\text{préfaisceaux}\\ \text{d'ensembles sur $C$}}} (\pi_0(C'/C), F),} \]
sans doute toujours \texttt{[un]} isomorphisme.

Si on travaille dans la catégorie dérivée des préfaisceaux des $k$-modules ($k$ un anneau), la formule précédente suggère conjecturalement
\[ \boxed{\mathbf{R}\Gamma(C', \alpha^*F) \simeq \mathbf{R}\underbrace{\operatorname{Hom}}_{\text{Hom internes}}\ (\underbrace{\mathbf{L}\alpha_!(k_{C'})}_{\substack{\text{homologie}\\ \text{relative de}\\ \text{$C'$ sur $C$}}}, F).} \]

\msp{page 8}
L'augmentation canonique
\[ \mathbf{L}\alpha_!(k_{C'}) \longrightarrow k_C \]
redonne, en prenant les $\mathbf{R}\mathrm{Hom}$ dans $F$,
\[ \mathbf{R}\Gamma(C,F) \longrightarrow \mathbf{R}\operatorname{Hom}(\mathbf{L}\alpha_!(k_{C'}), F) \xleftarrow{\ \sim\ } \mathbf{R}\Gamma(C', \alpha^*F) \]
l'homomorphisme canonique de `contravariance' de la cohomologie.

Si par contre on travaille avec le dérivateur non additif HOT, donc $F$ un `type d'homotopie relatif' sur $C$, on doit avoir
\[ \boxed{\mathbf{R}\Gamma_{\mathrm{HOT}}(C', \alpha^*F) \xrightarrow{\ \sim\ } \mathbf{R}\Gamma(C\,,\ \underbrace{\operatorname{Hom}}_{\substack{\text{Hom interne}\\ \text{de la catégorie des}\\ \text{types d'homotopie}\\ \text{relatifs sur $C$}}}\ (\ \alpha_!\ (\ \underbrace{e_{C'}}_{\substack{\text{objet final}\\ \text{de HOT($C'$),}\\ \text{type d'homotopie}\\ \text{relatif final}\\ \text{sur $C'$}}}\ ),\ F)).} \]

\dgrblock{dgrb81}{Global adjunction does not identify an internal Hom}
\begin{dcomment}{examples}{p8:Si par contre on travaille avec le dérivateur non additif HOT}{Global adjunction does not identify an internal Hom}
The proposed formula is useful because its global version is right. The difficulty begins only when its two sides are compared as objects over $C$. For presheaves of spaces, the correct equivalence is
\[
 \operatorname{Map}_{\mathcal P(C')}(1,\alpha^*F)
 \simeq
 \operatorname{Map}_{\mathcal P(C)}(\alpha_!1,F),
 \qquad \mathcal P(C)=\operatorname{Fun}(C^{\mathrm{op}},\mathcal S).
\]
Since $\mathcal P(C)$ is cartesian closed, the right side is also
$\mathbf R\Gamma(C,\underline{\operatorname{Hom}}(\alpha_!1,F))$.
This proves exactly the globally sectioned formula on page~8. The analogous abelian formula on page~7 takes values in $D(k)$ when ``Hom internes'' means the $k$-enriched global Hom. A sheaf-valued internal Hom would require an additional $\mathbf R\Gamma(C,-)$.

To see what global sections forget, take $C=[1]=(0\to1)$ and $i:e\to C$ with $i(*)=1$. A presheaf is an arrow $F_1\to F_0$. Direct calculation of all three adjoints gives
\[
 i^*F=F_1,\qquad i_!A=(A\xrightarrow{1}A),
 \qquad i_*A=(A\to *).
\]
Here the displayed order is always the restriction map from value at $1$ to value at $0$. Every $c/i$ is terminal, so $i$ passes even the contractible-comma criterion; $\Gamma(C,F)=F_1$ and the comparison with $\Gamma(e,i^*F)$ is an identity. Nevertheless the unit has component $F_0\to *$ at $0$, and the counit $i_!i^*F\to F$ has component $F_1\to F_0$ there. Neither is generally invertible.

Moreover, $i_!1=1_C$, hence
\[
 \underline{\operatorname{Hom}}_C(i_!1,F)=F,
 \qquad i_*i^*F=(F_1\to *).
\]
Taking $F=(\{*\}\to\{0,1\})$ disproves the internal identification. Applying global sections to both sides gives $\{*\}$, as it should. With $j(*)=0$ instead, every $j/c$ is terminal, whereas $j_*j^*F=(F_0\xrightarrow{1}F_0)$; thus the opposite comma orientation also fails to imply the universal unit assertion. The same examples work in $D(k)$, replacing $*$ by the zero module and choosing an arrow which is not an isomorphism.

The shift argument on page~9 establishes the expected Hom groups once the derived adjunction is available. Recovering an equivalence of enriched Hom objects requires the enriched adjunction itself, or a constructed comparison and a conservativity argument. It is not supplied merely by a bijection on $\pi_0$ in the space-valued case. No such extra step is needed for the explicit bar-resolution proof above.
\end{dcomment}


\printedpage{6}
% head: A. Grothendieck : Dérivateurs, Ch. II, Édition des Universités de Montpellier II et Paris VII
% foot: Version 8 juillet 2001
En fait, ces formules sont quasi-tautologiques, si on admet l'existence des foncteurs $\alpha_!$ comme adjoints \textit{à gauche} des $\alpha^*$. Dans le cas abélien, disons,
\msp{page 9}
l'isomorphie des $H^0$ des deux membres est la formule d'adjonction
\[ \operatorname{Hom}_{C'}(k_{C'}, \alpha^*F) \simeq \operatorname{Hom}_C(\alpha_!(k_{C'}), F), \]
d'où l'isomorphie pour les $H^i$ des deux membres, par décalage. De même dans le cas de HOT, l'isomorphie canonique des $\pi_0$ des deux membres est la formule d'adjonction. Pour avoir l'isomorphisme plus fin dans $\mathbf{D}(e)$, la catégorie fondamentale du dérivateur envisagé, en termes de Hom internes, on doit admettre l'existence d'objets unité (les $k_X$ dans un cas, les $e_X$ \texttt{[plutôt} $\varepsilon_X$\texttt{]} dans l'autre), compatibles avec images inverses, de telle façon \texttt{[que]} $\Homint(\varepsilon_X, \underbrace{F}_{\in\mathbf{D}_X}) \simeq F$. On aura
\[ \Gamma(C', \alpha^*F) \overset{\text{déf}}{=} f'_*(\alpha^*(F)) \simeq f_*(\alpha_*\alpha^*(F)), \]
et la question est donc de donner une interprétation de $\alpha_*\alpha^*F$ comme
\[ \alpha_*\alpha^*F \simeq \Homint(\alpha_!(\varepsilon_{C'}), F) \qquad \text{dans } \mathbf{D}(C). \]
Pour ceci, on va écrire
\[ \alpha^*(F) \simeq \Homint(\varepsilon_{C'}, F), \]
et la formule doit résulter d'une `formule d'adjonction interne'
\msp{page 10}
\[ ? \qquad\qquad \boxed{\alpha_*(\Homint_{C'}(G', \alpha^*F)) \simeq \Homint_C(\alpha_!G', F).} \]
(\textit{Formule d'adjonction `interne',} pour $\alpha_!$ et $\alpha^*$, valable en principe pour \textit{tout} $\alpha : C' \longrightarrow C$ dans Diag, tout $G'$ dans $\mathbf{D}(C')$ et tout $F$ dans $\mathbf{D}(C)$.)

\dgrblock{dgrb82}{The projection formula is an extra property}
\begin{dcomment}{exposition}{p9:et la formule doit résulter}{The projection formula is an extra property}
Grothendieck correctly reduces the proposed internal adjunction to a projection formula. What fails is the expectation that every change of indexing category has this property. In a closed symmetric monoidal setting, write $L=\alpha_!$, $R=\alpha^*$ and $U=\alpha_*$. The canonical map supplied by the adjunction $L\dashv R$ has direction
\[
 L(RG\otimes G')\longrightarrow G\otimes LG'.
\]
It is adjoint to $RG\otimes G'\to RG\otimes RLG'$, using the unit and the monoidal structure of $R$. The arrow printed on page~10 points in the inverse direction; such an inverse is available only after invertibility has been established.

For every $H$, the adjunctions give
\[
\begin{aligned}
 \operatorname{Hom}(H,U\underline{\operatorname{Hom}}(G',RF))
 &\cong\operatorname{Hom}(L(RH\otimes G'),F),\\
 \operatorname{Hom}(H,\underline{\operatorname{Hom}}(LG',F))
 &\cong\operatorname{Hom}(H\otimes LG',F).
\end{aligned}
\]
Thus the canonical internal comparison
$\underline{\operatorname{Hom}}(LG',F)\to
U\underline{\operatorname{Hom}}(G',RF)$
is invertible for all $F,G'$ exactly when the projection maps are invertible for all $H,G'$. This proves the reduction, and also locates the additional hypothesis precisely.

For cartesian products this property is called Frobenius reciprocity. An essential geometric morphism need not satisfy it. Moreover, Frobenius reciprocity itself is weaker than its validity in every slice, the latter being the condition appropriate to a locally connected geometric morphism. \citet[\S1]{GarnerStreicher2021} explicitly distinguish these conditions; they also recall the equivalence between Frobenius reciprocity and preservation of exponentials by the right adjoint of an adjunction between cartesian closed categories. These later distinctions explain why existence of both Kan extensions does not settle the question. They should not be replaced by an unrestricted ``projection formula for derivators''.
\end{dcomment}


Comment prouver cette formule ? En prenant l'opération produit $\otimes$ associé à l'opération Hom\,int, on est ramené (en prenant les foncteurs contravariants associés) à une formule
\[ G \otimes \alpha_!G' \xrightarrow{\ \sim\ } \alpha_!(\alpha^*G \otimes G'). \]
On trouve par adjonction un homomorphisme canonique, il faut voir s'il est toujours \texttt{[un]} isomorphisme.

\dgrblock{dgrb83}{The exact obstruction for a point inclusion}
\begin{dcomment}{development}{p10:On trouve par adjonction un homomorphisme canonique}{The exact obstruction for a point inclusion}
The walking-arrow example can be sharpened to measure exactly which restriction maps obstruct the projection formula. This makes the question on page~10 concrete even before a general theory of Frobenius adjunctions is introduced.

\textbf{Point-inclusion criterion.}\ours\
Let $d$ be an object of a small category $C$, and let $i:e\to C$ select $d$. For a set-valued presheaf $G$, the projection map
\[
 i_!(i^*G\times A)\longrightarrow G\times i_!A
\]
is invertible for every set $A$ if and only if $G(f):G(d)\to G(c)$ is bijective for every arrow $f:c\to d$. Consequently the formula holds for all $G,A$ if and only if every arrow with codomain $d$ is an isomorphism.

\textit{Proof.}
The left adjoint to evaluation at $d$ is $i_!A=A\times h_d$, where $h_d(c)=C(c,d)$. At $c$ the projection map is
\[
 A\times G(d)\times C(c,d)\longrightarrow
 A\times G(c)\times C(c,d),\qquad
 (a,s,f)\longmapsto(a,G(f)s,f).
\]
For $A=1$ it is a disjoint union, indexed by $f$, of the maps $G(f)$; it is bijective precisely when each of them is. This also suffices for every $A$. If every $f$ is invertible, every $G(f)$ is bijective. Conversely, suppose the formula holds for every presheaf. Taking $G=h_z$ gives a bijection $C(d,z)\to C(c,z)$ by precomposition with $f$, for every $z$. With $z=c$, obtain $g:d\to c$ such that $gf=1_c$. With $z=d$, injectivity and $(fg)f=f=1_df$ imply $fg=1_d$. Thus $f$ is invertible.\hfill$\square$

This includes the object $1$ of $[1]$, where the noninvertible incoming arrow is exactly the obstruction. It also explains why the analogous formula for a group action behaves differently: all indexing arrows are invertible. For a fixed $G$, local constancy along the arrows entering $d$ is the precise repair; no conclusion about all coefficients follows from connectedness of a comma category.

We checked the pointwise Kan-extension treatment of \citet[\S6.2]{Riehl2016} and the Frobenius discussion of \citet[\S1]{GarnerStreicher2021}, and searched for the point-evaluation formulation. We could not find this elementary criterion stated there in this form. It is recorded as our worked continuation of the calculation, with no claim of priority.
\end{dcomment}


\printedpage{7}
% head: A. Grothendieck : Dérivateurs, Ch. II, Édition des Universités de Montpellier II et Paris VII
% foot: Version 8 juillet 2001
\msp{page 11}

\dgrblock{dgrb84}{Why the argument passes through toposes}
\begin{dcomment}{context}{p10}{Why the argument passes through toposes}
The change of subject at the beginning of \S2 is part of the cofinality problem. Grothendieck wants conditions on indexing categories which survive passage to categories of coefficients and which force the relevant homotopy types to disappear. A formulation in terms of the presheaf topos can detect all objects over the category, whereas the underlying preorder remembers only whether an arrow exists.

This is a small-scale instance of the program described in \citet[\S69, especially the discussion of diagram categories on pp.~1--3 of the excerpt]{Grothendieck1983}: the derived category of coefficients is accompanied by derived categories of diagrams of every shape. The concern there comes from operations on coherent diagrams which a single triangulated category does not retain. The present calculation asks what changing the shape does to sections and their derived functors. In parallel, \citet[Chapter~II, \S1; Chapter~III, \S2]{Heller1988} axiomatizes homotopy theories and develops the computation of homotopy Kan extensions. These references explain the setting, not a historical claim that one author's argument was taken from the other.

The additional tensor language on pages~7--10 needs separate structure. Heller's Chapter~IV treats tensor and cotensor operations; the preceding projection counterexample shows why that structure cannot be inferred from the mere existence of left and right Kan extensions. We now return to a question for which the ordinary presheaf topos is already sufficient: which part of the filtering conditions is visible in the behavior of its connected objects?
\end{dcomment}


\section*{2\quad Les catégories localement filtrantes sont les catégories localement totalement $0$-connexes}

Soit $\mathcal{T}$ un topos. On dit que $\mathcal{T}$ est \textit{irréductible} si l'ensemble ordonné $\operatorname{Ouv}(\mathcal{T})$ des sous-objets de $e_{\mathcal{T}}$ a la propriété suivante : a) $\operatorname{Ouv}(\mathcal{T})$ n'est pas réduit à un élément, i.e. $\emptyset_{\mathcal{T}} \longrightarrow e_{\mathcal{T}}$ n'est pas un isomorphisme ($\mathcal{T}$ n'est pas identique à son sous-topos vide). b) Si $U'$, $U''$ sont deux objets `non vides', i.e. non initiaux, de $\operatorname{Ouv}(\mathcal{T})$, leur intersection est `non vide'. Il revient au même pour b) de dire que si $\mathcal{T}'$, $\mathcal{T}''$ sont deux sous-topos fermés distinct de $\mathcal{T}$, alors leur supremum $\mathcal{T}' \cup \mathcal{T}''$ est distinct de $\mathcal{T}$, ou encore que tout ouvert de $\mathcal{T}$ est connexe. On dit que $\mathcal{T}$ est \textit{localement irréductible} si a) $\mathcal{T}$ est localement connexe (i.e. chaque faisceau sur $\mathcal{T}$ est somme de faisceaux 0-connexes), et si b) pour tout faisceau 0-connexe $U$ sur $\mathcal{T}$, le topos induit $\mathcal{T}/U$ est irréductible (i.e. l'ensemble ordonné $\operatorname{Ouv}(U)$ des sous-objets de $U$ a la propriété b) ci-dessus ( a) signifie simplement que $U \not\simeq \emptyset_{\mathcal{T}}$, i.e. que $U$ n'est pas objet initial dans $\mathcal{T}^\wedge$ ; ou plus simplement) : tout sous-faisceau d'un faisceau connexe est connexe (\footnote[5]{\textbf{NB} Irréductible n'implique pas localement irréductible, pas plus que connexe n'implique pas localement connexe etc.}) (\footnote[6]{$\mathcal{T}$ est localement irréductible si et seulement s'il a un ensemble des générateurs $(U_\alpha)$, avec les $U_\alpha$ irréductibles (donc 0-connexes, donc $\mathcal{T}$ est localement connexe).}).

\dgrblock{dgrb85}{Irreducibility of an object and of its support}
\begin{dcomment}{exposition}{p11:tout sous-faisceau d'un faisceau connexe est connexe}{Irreducibility of an object and of its support}
Grothendieck is distinguishing two ways to forget morphisms. The space $\operatorname{Esp}(X)$ remembers downward-closed sets of objects; the topos $\widehat X=\operatorname{Fun}(X^{\mathrm{op}},\mathrm{Set})$ also remembers every category of elements. The first sees a common predecessor of two objects. The second must see common predecessors compatible with prescribed arrows.

Here an object is \emph{connected} if it has no decomposition into two noninitial disjoint summands; ``$0$-connected'' adds noninitiality. Thus the assertion that every subobject of a connected object is connected allows the initial subobject. An object $A$ is irreducible when it is noninitial and any two noninitial subobjects of $A$ meet. Equivalently, every subobject is connected: a disconnected subobject gives two disjoint nonempty subobjects, while disjoint subobjects have a disconnected union. Applied to $A=1$, this gives the definition of an irreducible topos. Applied to every nonempty connected $A$, after assuming local connectedness, it gives the local definition.

For $\widehat X$, subterminal objects are precisely sieves of objects: if $x$ belongs and $y\to x$ exists, then $y$ belongs. Any two nonempty such sieves meet exactly when each pair of objects has a common predecessor. This proves
\[
 \widehat X\text{ irreducible}
 \quad\Longleftrightarrow\quad
 X\ne\varnothing\text{ and }
 (\forall x,y)(\exists z)(z\to x, z\to y).
\]
It is a condition on the associated preorder, not on equalities of arrows. A least class in that preorder need not give an initial object: uniqueness of arrows has been discarded. If the preorder quotient is artinian, a minimal class exists, and common lower bounds force it below every other class; this proves the assertion on page~13. Finite principal lower sets suffice as well. ``Locally finite'' here must have that meaning, as indicated by the finite slices $X/x$; finiteness of intervals alone would fail, for example for the ordered set $\mathbb Z$.

The list of equivalences on pages~12--13 is provisional. Its claim about components being weakly cofiltered is too weak, and its subsequent inclusion of PS~2 is too strong. The annotations and Proposition~2.1 correct these in opposite directions. The surviving characterization is PS~1 alone. We prove it below, retaining the successive formulations as part of the text's development.
\end{dcomment}


Quand $\mathcal{T}$ est de la forme $\operatorname{Top}(T)$, où $T$ est un espace topologique, alors $\mathcal{T}$ irréductible signifie que $T$ est irréductible (car alors $\operatorname{Ouv}(\mathcal{T}) \simeq \operatorname{Ouv}(T)$). Cela implique déjà trivialement que $T$ est localement connexe (tout ouvert est somme disjoint d'ouverts connexes), donc $\operatorname{Top}(\mathcal{T})$ est localement connexe. De plus, cela implique aussi que $\operatorname{Top}(T)$ est localement irréductible (\footnote[7]{expliciter la démonstration.}). (\textbf{NB} Les objets 0-connexes % [sic] \operatorname{Top}(\mathcal{T}) -> \operatorname{Top}(T), Top s'applique à l'espace T et non au topos \mathcal{T}
\msp{page 12}
de $\operatorname{Top}(T)^\wedge$, i.e. les faisceaux d'ensembles 0-connexes sur $T$, sont ceux dont la `fibre générique' a exactement un point \dots) Aussi $\mathcal{T} = \operatorname{Top}(T)$ est localement irréductible si et seulement si $T$ \texttt{[est]} localement connexe, et les composantes connexes de $T$ sont irréductibles, i.e. si et seulement si $T = \coprod_i T_i$, les $T_i$ irréductibles.

\dgrblock{dgrb86}{The generic fibre without a generic point}
\begin{dcomment}{examples}{p12:les composantes connexes de $T$ sont irréductibles}{The generic fibre without a generic point}
The topological case supplies the intuition for the definition, but it also calls for a precise interpretation of ``generic fibre''. An irreducible space need not possess a generic point. For a sheaf $F$ on a nonempty irreducible space $T$, define instead
\[
 \gamma(F)=\mathop{\operatorname{colim}}_{\varnothing\ne U\subseteq T}F(U),
\]
where transition maps are restrictions as $U$ shrinks. Nonempty opens form a filtering system under intersection. If a generic point exists, this is its stalk; otherwise it is the stalk at the corresponding point of the soberification. The following direct argument avoids requiring that point.

For every nonempty open $U$, a section of $F(U)$ has a well-defined generic germ. Because $U$ is connected, the constant sheaf with value $\gamma(F)$ has sections $\gamma(F)$ on $U$. We obtain a map $F\to\underline{\gamma(F)}$ and a decomposition
\[
 F\cong\coprod_{g\in\gamma(F)}F_g
\]
into its fibres. Every $F_g$ is noninitial because $g$ has a representative section. If $A,B$ are two noninitial subobjects of $F_g$, choose sections $a\in A(U)$ and $b\in B(V)$ on nonempty opens. They have the same generic germ, so on some nonempty $W\subseteq U\cap V$ their restrictions agree. This supplies a section of $A\cap B$, proving that $F_g$ is irreducible.

Consequently these fibres are exactly the connected components of $F$, and $F$ is $0$-connected if and only if $\gamma(F)$ is a singleton. This proves both the assertion about the generic fibre and the missing implication in footnote~7: the sheaf topos of an irreducible space is locally irreducible. More generally, if $T$ is a disjoint union of open irreducible spaces, the conclusion follows component by component. Conversely, a locally irreducible sheaf topos is locally connected; its terminal object decomposes into connected subterminal objects, which are disjoint open subsets of $T$, and each is irreducible. This proves the topological characterization stated here.

For a general presheaf topos, no generic point of the underlying preorder can replace this argument. Parallel arrows produce different objects in slices, even when they have the same source and target. The next example shows exactly where the analogy with topological spaces ceases to suffice.
\end{dcomment}


Si $\mathcal{T}$ est de la forme $\operatorname{Top}(X)$, $X$ une catégorie, on sait que $\mathcal{T}$ est localement connexe (et même localement $\infty$-connexe). On voit d'autre part que $\mathcal{T}$ est irréductible si et seulement si $\operatorname{Esp}(X)$ (l'espace topologique associé, dont l'ensemble sous-jacent est $\operatorname{Ob} X$) est irréductible (\footnote[8]{car les ouverts de $\operatorname{Top}(X)$, ceux de $\operatorname{Esp}(X)$ et ceux de $X$ se correspondent.}) ; et on dit alors que $X$ est irréductible (\footnote[9]{Cela signifie aussi que pour $x$, $y$ dans $X$, $\exists\, z$ avec \begin{tikzcd}[ampersand replacement=\&, row sep=0.2em, column sep=1.2em, baseline=(\tikzcdmatrixname-2-1.base)] \& x \\ z \arrow[ur] \arrow[dr] \& \\ \& y \end{tikzcd}, i.e. $\operatorname{Ord}(X)$ est filtrant à gauche.}). Comme les composantes connexes de $X$ et celles de $\operatorname{Esp}(X)$ ou $\operatorname{Top}(X) = \mathcal{T}$ se correspondent, on voit donc que les conditions suivantes sont équivalentes : a) le topos $\operatorname{Top}(X)$ est localement irréductible, b)
\printedpage{8}
% head: A. Grothendieck : Dérivateurs, Ch. II, Édition des Universités de Montpellier II et Paris VII
% foot: Version 8 juillet 2001
$\operatorname{Esp}(X)$ est localement irréductible, c) les composantes connexes de $X$ sont irréductibles, d) les $X/x$ \texttt{[sont]} irréductibles. On dit alors que $X$ est localement irréductible (\footnote[10]{Faux, on a (a $\Longleftrightarrow$ d) $\Longrightarrow$ (b $\Longleftrightarrow$ c) $X$ irréductible n'implique \textit{pas} que les $X/x$ le soient, exemple $X = (0 \overset{\alpha}{\underset{\beta}{\rightrightarrows}} 1)$ $(\alpha \neq \beta)$, on a $X/1 =$ \begin{tikzcd}[ampersand replacement=\&, row sep=0.2em, column sep=1.2em, baseline=(\tikzcdmatrixname-2-2.base)] 0' \arrow[dr] \& \\ \& 1 \\ 0'' \arrow[ur] \& \end{tikzcd}, qui n'est \textit{pas} irréductible, ni localement irréductible puisqu'il est connexe. Pourtant il est étale sur $X$.}).

Cela implique donc que pour tout $X'$ étale sur $X$, avec $X'$ 0-connexe, $X'$ est irréductible.

\textit{Dire conditions équivalentes.}
\begin{enumerate}
\item[a)] $\operatorname{Top}(X)$ est localement irréductible.
\item[b)] Les $X/x$ \texttt{[sont]} irréductibles.
\item[c)] Tout objet 0-connexe de $X^\wedge$ est irréductible.
\item[d)] Les composantes connexes de $X$ sont faiblement cofiltrantes (i.e. faiblement filtrantes à gauche).
\item[e)] Condition PS $1_{\mathrm{g}}$ de SGA 4, I 2.
\end{enumerate}
Si $X$ est \texttt{[une]} catégorie, $X$ irréductible signifie aussi que $X \neq \emptyset$ et $\forall\, x, y \in \operatorname{Ob} X$, \texttt{[il]} existe \texttt{[un]} $s$ avec \begin{tikzcd}[row sep=0.2em, column sep=1.2em, baseline=(\tikzcdmatrixname-2-1.base)] & x \\ s \arrow[ur] \arrow[dr] & \\ & y \end{tikzcd}. (Cela signifie aussi que $X$ est \textit{faiblement filtrante} (\footnote[11]{Il est mieux l'appeler \textit{pseudo-filtrante,} et réservons faiblement filtrante à : localement faiblement filtrante + connexe.}) à gauche, ou faiblement cofiltrante.) $X$ localement irreductible signifie que pour $x$, $y$ \textit{appartenant à une même composante connexe,} \texttt{[il]} existe \texttt{[un]} $s$ avec \begin{tikzcd}[row sep=0.2em, column sep=1.2em, baseline=(\tikzcdmatrixname-2-1.base)] & x \\ s \arrow[ur] \arrow[dr] & \\ & y \end{tikzcd}, % [sic] irreductible -> irréductible, accent omis
\msp{page 13}
i.e. $X/[x] \cap X/[y] \neq \emptyset$ (où $[x]$ est considéré comme élément de l'ensemble ordonné $\operatorname{Ob} X$). Il revient encore au même de dire que les catégories induites $X/x$ sont irréductibles (\footnote[12]{faux, cf. annotation marginale plus haut.}).

\textbf{Exemple.} Une catégorie \textit{ayant un objet initiale,} ou seulement un `plus petit objet', est irréductible. Inversement, si l'ensemble \textit{ordonné} $\operatorname{Ord}(X)$ associé à $X$ (quotient de $\operatorname{Ob}(X)$ par la relation d'équivalence associé à son préordre) est fini, p.ex. $X$ fini, ou seulement localement fini, p.ex. $X$ localement fini (i.e. les catégories $X/x$ finies), alors si $X$ est irréductible, $\operatorname{Ob} X$ a un plus petit élément, i.e. $X$ a un plus petit objet. C'est vrai plus généralement encore, si $\operatorname{Ord}(X)$ est artinien, i.e. toute suite décroissante dans $\operatorname{Ord}(X)$ est stationnaire. % [sic] initiale -> initial, accord avec « objet »

Mais la condition intéressante est la condition topossique de \textit{locale} irréductibilité, qui équivaut : $I$ pseudo-filtrante à gauche, i.e. les composantes connexes de $I$ filtrantes, i.e. les conditions PS 1, PS 2

PS 1\enspace Tout diagramme \begin{tikzcd}[row sep=0.2em, column sep=1.2em, baseline=(\tikzcdmatrixname-2-2.base)] x \arrow[dr] & \\ & s \\ y \arrow[ur] & \end{tikzcd} se complète en \begin{tikzcd}[row sep=0.2em, column sep=1.2em, baseline=(\tikzcdmatrixname-2-1.base)] & x \arrow[dr] & \\ z \arrow[ur] \arrow[dr] & & s \\ & y \arrow[ur] & \end{tikzcd} \textit{commutatif.}

PS 2\enspace Pour toute double flèche $x \rightrightarrows y$, \texttt{[il]} existe un égalisateur $z \longrightarrow x \rightrightarrows y$.

% ---- from Der-ChII-234_p011-015.tex
\printedpage{9}
% head: A. Grothendieck : Dérivateurs, Ch. II, Édition des Universités de Montpellier II et Paris VII
% foot: Version 8 juillet 2001

\msp{page 14}

En résumé, on a la

\noindent\textbf{Proposition 2.1.} \textit{Soit $X$ un objet de} Cat. \textit{Conditions équivalentes :}
\begin{enumerate}[label={},leftmargin=3em]
\item[\textit{a)}] \textit{Le topos associé} $\operatorname{Top}(X)$ \textit{est localement irréductible, i.e. (étant localement connexe) pour tout objet $F$ de $X^{\wedge}$ et tout sous-objet $F'$ de $F$, si $F$ est connexe, $F'$ l'est.}
\item[\textit{a$'$)}] \textit{Pour tout $X'$ étale sur $X$, si $X'$ est connexe, tout ouvert de $X'$ est connexe.}
\item[\textit{a$''$)}] \textit{Tout ouvert d'une catégorie localisée $X/a$ ($a \in \operatorname{Ob} X$) est connexe (i.e. $X/a$ est une catégorie irréductible, i.e. le topos} $\operatorname{Top}(X/a)$ \textit{l'est).}
\item[\textit{b)}] \textit{Pour tout objet $a$ de $X$,} $\operatorname{Esp}(X/a)$ \textit{est un espace topologique irréductible.}
\item[\textit{c)}] \textit{Les catégories $X/a$ sont faiblement filtrantes à gauche (ou faiblement cofiltrantes), i.e. pour deux objets $\xi$, $\eta$ d'une telle catégorie,} \texttt{[il]} \textit{existe un objet $\zeta$ et deux flèches}
\begin{tikzcd}[ampersand replacement=\&, row sep=tiny, column sep=small, baseline=(\tikzcdmatrixname.center)]
 \& \xi \\
\zeta \arrow[ur] \arrow[dr] \& \\
 \& \eta
\end{tikzcd}
\textit{(i.e. l'ensemble ordonné Ord($X/a$) associé à $X/a$ est filtrant à gauche)}
\item[\textit{c$'$)}] \textit{Pour tout diagramme en traits pleins dans $X$,}
\begin{tikzcd}[ampersand replacement=\&, row sep=tiny, column sep=small, baseline=(\tikzcdmatrixname.center)]
 \& x \arrow[dr] \& \\
z \arrow[ur, dotted] \arrow[dr, dotted] \& \& s \\
 \& y \arrow[ur] \&
\end{tikzcd}
\textit{, on peut le compléter en un carré} commutatif \textit{comme indiqué. (Condition} PS 1 \textit{des catégories localement filtrantes.)}
\end{enumerate}

\noindent Ces conditions impliquent que toute composante connexe $X_{\alpha}$ de $X$ est irréductible, ou ce qui revient au même, que $\operatorname{Esp}(X_{\alpha})$ est irréductible,

\msp{page 15}

\noindent mais ne sont pas impliquées par celle-ci.

Contrexemple, la catégorie 0-connexe $0 \underset{\beta}{\overset{\alpha}{\rightrightarrows}} 1$, avec $\alpha \neq \beta$. On a $X/1 \simeq$
\begin{tikzcd}[ampersand replacement=\&, row sep=tiny, column sep=small, baseline=(\tikzcdmatrixname.center)]
0 \arrow[dr] \& \\
 \& 1 \\
0' \arrow[ur] \&
\end{tikzcd}, qui n'est \textit{pas} irréductible, puisque $\Phi \backslash \{1\} \simeq \varepsilon$ est un ouvert disconnexe \texttt{[}$\Phi = X/1$\texttt{,} $\varepsilon = \bullet\ \bullet$\texttt{]}.

J'aimerais une condition `topossique' sur $X$, équivalente à $X$ `pseudo-filtrante à gauche', i.e. les composantes connexes filtrantes, i.e. PS 1\textsuperscript{opp} (condition c$'$) ci-dessus) + PS 2\textsuperscript{opp} (condition de la double flèche), et qui de plus implique que tout objet 0-connexe de $X^{\wedge}$ (i.e. tout objet $X'$ de Cat étale sur $X$ et 0-connexe) soit $W_{\infty}$-asphérique. Et je songe à la notion de `totale asphéricité' de Pursuing Stacks.

\dgrblock{dgrb87}{Completing Proposition 2.1}
\begin{dcomment}{exposition}{p15:J'aimerais une condition}{Completing Proposition 2.1}
Proposition~2.1 isolates the part of filtering which can be expressed without equalizing parallel arrows. We give its proof because it explains both the corrected statement and the mechanism which will later be strengthened.

Recall that PS~1 says that every cospan $a\to x\leftarrow b$ admits a commuting square with some upper-left object $c$; no universal property is requested. Suppose PS~1 holds. In any category with this property, the relation ``$a$ and $b$ have a common predecessor'' is reflexive and symmetric. It is transitive: given $a\leftarrow r\to b$ and $b\leftarrow s\to d$, apply PS~1 to $r\to b\leftarrow s$ and compose the resulting arrows. Every arrow connects objects related in this sense. Hence any two objects in the same connected component have a common predecessor.

For a presheaf $F$, its category of elements $E=\int_XF\to X$ is a discrete fibration, the meaning of ``\'etale'' here. Every arrow into the image of an object has a unique lift into that object. Thus a PS~1 square in $X$ lifts to a PS~1 square in $E$: the two lifts into the common target agree by uniqueness. Therefore $E$ also satisfies PS~1. Components of $E$ correspond to coproduct components of $F$. If $F$ is $0$-connected, $E$ is nonempty and connected, so the previous paragraph gives common predecessors for every pair of its objects. Subobjects of $F$ correspond to sieves in $E$; any two nonempty such sieves meet. Thus $F$ is irreducible, and $\widehat X$ is locally irreducible.

Conversely, suppose every $0$-connected presheaf is irreducible. The representable $h_x$ is $0$-connected, since its category of elements $X/x$ has a terminal object. Given arrows $a\to x$ and $b\to x$, their generated subobjects of $h_x$ are noninitial. Their intersection is noninitial, so it contains some element represented by arrows $c\to a$ and $c\to b$ with equal composite to $x$. This is PS~1. The descriptions using $X/x$, its associated space, and open subcategories are equivalent by the preceding sieve argument. This proves all the equivalences of Proposition~2.1.

The counterexample is the category with two objects and distinct arrows $\alpha,\beta:0\rightrightarrows1$. Its objects have the common predecessor $0$, so its underlying preorder is the irreducible two-point chain. But in $X/1$, the objects $(0,\alpha)$ and $(0,\beta)$ have no common predecessor: the only endomorphism of $0$ is its identity. Their two singleton lower sets are disjoint nonempty subobjects of $h_1$. Thus $\widehat X$ is not locally irreducible. The mistaken converse on page~16 is precisely the assertion that the preorder remembers this equality of composites; the footnotes already retract it.
\end{dcomment}


\noindent\textbf{Définition 2.2.} (\footnote[13]{Cf. tableau récapitulatif pages 97, 98 et surtout 99.})\ \ Soit $X$ une catégorie.
\begin{enumerate}[label={},leftmargin=3em]
\item[a)] On dit que $X$ est \textit{faiblement filtrante} (\footnote[14]{dire plutôt \textit{pseudo-filtrante}, réserver `faiblement filtrante' à `localement faiblement filtrante' (ou \textit{irréductible}) + connexe, condition PS $1'_{\mathrm{g}}$.}) si $X \neq \emptyset$, et si pour $a, b \in \operatorname{Ob} X$, $\exists\, c$ et des flèches $c$
\begin{tikzcd}[ampersand replacement=\&, row sep=tiny, column sep=small, baseline=(\tikzcdmatrixname.center)]
 \& a \\
{} \arrow[ur, phantom] \& \\
 \& b
\end{tikzcd}
\hspace{-2.6em}% layout aid only
\begin{tikzcd}[ampersand replacement=\&, row sep=tiny, column sep=small, baseline=(\tikzcdmatrixname.center)]
 \& {} \\
{} \arrow[ur] \arrow[dr] \& \\
 \& {}
\end{tikzcd}
, i.e. $X/[a] \cap X/[b] \neq \emptyset$. (\textbf{NB} $X$ est faiblement filtrante (\footnote[15]{pseudo-filtrante.}) à gauche si et seulement si $\operatorname{Ord}(X)$ l'est, i.e. est filtrante à gauche. \textbf{ReNB} $X$ faiblement filtrante $\Longrightarrow$ $X$ 0-connexe.)

\item[b)] \printedpage{10}
% head: A. Grothendieck : Dérivateurs, Ch. II, Édition des Universités de Montpellier II et Paris VII
% foot: Version 8 juillet 2001
On dit que $X$ est \textit{localement faiblement filtrante à gauche} ou plus simplement préfiltrante à gauche si les catégories $X/x$ ($x \in \operatorname{Ob} X$) sont faiblement filtrantes à gauche, i.e. si $X$ satisfait la condition : $\forall$ diagrammes
\begin{tikzcd}[ampersand replacement=\&, row sep=tiny, column sep=small, baseline=(\tikzcdmatrixname.center)]
a \arrow[dr] \& \\
 \& x \\
b \arrow[ur] \&
\end{tikzcd}, $\exists\, c$ dans $X$ et des flèches de $c$ dans $a$, $b$ rendant \textit{commutatif} le carré
\[
\begin{tikzcd}[ampersand replacement=\&, row sep=tiny, column sep=small]
 \& a \arrow[dr] \& \\
c \arrow[ur, dotted] \arrow[dr, dotted] \& \& x \\
 \& b \arrow[ur] \&
\end{tikzcd}
\]
(\footnote[16]{Condition PS $1_{\mathrm{g}}$.}). (\textbf{NB} Cette condition implique que $\operatorname{Ord}(X)$ est localement faiblement filtrant, mais n'équivaut pas à cette dernière. Exemple, $X = 0 \underset{\beta}{\overset{\alpha}{\rightrightarrows}} 1$, $X/1 =$
\begin{tikzcd}[ampersand replacement=\&, row sep=tiny, column sep=small, baseline=(\tikzcdmatrixname.center)]
0' \arrow[dr] \& \\
 \& 1 \\
0'' \arrow[ur] \&
\end{tikzcd}
n'est pas faiblement filtrante à gauche.)

\msp{page 16}

\textbf{NB} Si $X$ localement faiblement filtrante à gauche, tout $X'$ étale sur $X$ itou.

\textbf{NB} (\footnote[17]{\textbf{Faux.}}) $X$ localement faiblement filtrante à gauche $\Longleftrightarrow$ les composantes connexes sont faiblement filtrantes à gauche, ou encore si $X$ est 0-connexe, $X$ localement faiblement filtrante à gauche $\Longleftrightarrow$ $X$ est faiblement filtrante (\footnote[18]{pseudo-filtrante.}) à gauche. On a en effet $\Longleftarrow$ tautologiquement (\footnote[19]{\textbf{faux,} ainsi $0 \underset{\beta}{\overset{\alpha}{\rightrightarrows}} 1$ est 0-connexe et faiblement filtrante à gauche, mais non pas localement faiblement filtrante à gauche car
\begin{tikzcd}[ampersand replacement=\&, row sep=tiny, column sep=small, baseline=(\tikzcdmatrixname.center)]
0 \arrow[dr, "\alpha"] \& \\
 \& 1 \\
0 \arrow[ur, "\beta"'] \&
\end{tikzcd}
ne se complète pas en un carré \textit{commutatif.}}), et $\Longrightarrow$ en montrant que si PS 1 \texttt{[est]} satisfait, alors $PS(a,b)$ défini par $X/[a] \cap X/[b] \neq \emptyset$ est une relation d'équivalence, impliquée par la relation $C(a,b)$ : $\{\operatorname{Hom}(a,b) \neq \emptyset$ \textit{ou} $\operatorname{Hom}(b,a) \neq \emptyset\}$, donc la relation d'équivalence (homotopie) engendrée par $C(a,b)$ est plus fine que $PS(a,b)$.

\item[c)] On dit que $X$ est \textit{filtrante à gauche}, ou \textit{cofiltrante,} si elle est faiblement filtrante (\footnote[20]{pseudo-filtrante.}) à gauche (donc \textit{non vide}), et satisfait (en plus de la condition PS $1'_{\mathrm{g}}$ de a)) la condition PS $2_{\mathrm{g}}$ suivante :

Pour toute double flèche $a \underset{v}{\overset{u}{\rightrightarrows}} b$, \texttt{[il]} existe \texttt{[une]} flèche $c \xrightarrow{w} a$ telle que $uw = vw$ (\footnote[21]{i.e. trois conditions a) $X \neq \emptyset$ b) PS $1'$ habituelle c) PS 2 (double flèche).}).

Il revient au même de dire que $X$ satisfait les trois conditions a$'$) $X$ est 0-connexe (plus forte que $X \neq \emptyset$), b$'$) PS 1, i.e. $X$ localement faiblement filtrante à gauche et c$'$) PS 2 comme avant.

\item[d)] On dit que $X$ est \textit{localement filtrante à gauche} si les composantes connexes sont filtrantes à gauche, i.e. si $X$ satisfait les conditions PS 1 et PS 2 (`pseudo-cofiltrante' dans la terminologie de SGA 4). \textbf{NB} Si $X$ est localement filtrante à gauche, il en est de même de tout $X'$ étale sur $X$, car c'est évident pour PS 1 qui concerne les $X/x$, et pour PS 2 cela résulte du fait que $X' \longrightarrow X$ est fidèle. Donc si $X$ est localement cofiltrante, les $X/a$ sont filtrantes, mais la réciproque est fausse.
\end{enumerate}

\dgrblock{dgrb88}{The equalization condition survives beyond every slice}
\begin{dcomment}{examples}{p16:mais la réciproque est fausse.}{The equalization condition survives beyond every slice}
The definitions now separate PS~1 from the additional condition PS~2. In the latter, ``\'egalisateur'' means an \emph{equalizing arrow}: for $u,v:a\rightrightarrows b$ there is $w:c\to a$ with $uw=vw$. It does not mean that a categorical equalizer exists. This is the usual cofiltered condition, dual to SGA~4's PS~2 \citep[Expos\'e~I, \S2.7]{SGA4I1972}.

To see why all slices can be cofiltered without the original category being locally cofiltered, take the one-object category $BG$ for a nontrivial group $G$. It satisfies PS~1: a pair $f,g$ can be completed by $1$ and $g^{-1}f$. For its unique object $*$, the slice $BG/*$ has objects $g\in G$ and exactly one arrow from $g$ to $h$, namely $h^{-1}g$. It is equivalent to the terminal category and is cofiltered. Yet distinct $u,v\in G$ cannot be equalized on the right, since $uw=vw$ and invertibility of $w$ imply $u=v$. Hence $BG$ fails PS~2.

This example also makes the topos distinction visible. Every connected $G$-set is transitive and therefore irreducible, so the presheaf topos is locally irreducible. Products of connected objects need not remain connected: for the regular right $G$-set, the diagonal right action on $G\times G$ has one orbit for each value of $hg^{-1}$. The product condition introduced next detects the missing equalization information. No defect in Proposition~2.1 is involved; its purpose was exactly to isolate PS~1.
\end{dcomment}


\printedpage{11}
% head: A. Grothendieck : Dérivateurs, Ch. II, Édition des Universités de Montpellier II et Paris VII
% foot: Version 8 juillet 2001

\noindent On a les notions duales : \textit{faiblement filtrante à droite} (ou simplement faiblement filtrante), et \textit{filtrante à droite} ou simplement filtrante. (Notions duales de a), c) ci-dessus.) Pour la notion duale de b), savoir que les $x\backslash X$ soient faiblement filtrantes, on ne peut plus l'appeler `localement faiblement filtrante', mais ce serait \textit{co}localement faiblement filtrante.

\noindent\textbf{Exemple.} $X = \Bigl(0 \xrightarrow{\gamma} 1 \underset{\beta'}{\overset{\alpha'}{\rightrightarrows}} 2 \underset{\beta}{\overset{\alpha}{\rightrightarrows}} 3\Bigr)$ avec $(\alpha\alpha' = \beta\beta') \neq (\alpha\beta' = \beta\alpha')$.

\dgrblock{dgrb89}{The finite example and its homotopy type}
\begin{dcomment}{development}{p16:avec $(\alpha\alpha' = \beta\beta')}{The finite example and its homotopy type}
The example between Definition 2.2 and Proposition 2.3 makes the loss of information in $\operatorname{Ord}(X)$ concrete. We read the displayed graph and its two relations as a presentation of $X$; in particular, no additional relation is imposed on composites with $\gamma$. Put
\[
 r=\alpha\alpha'=\beta\beta',\qquad
 s=\alpha\beta'=\beta\alpha',\qquad r\ne s.
\]
Then $\operatorname{Ord}(X)$ is the chain $0<1<2<3$. Every pair of objects has both a common predecessor and a common successor. Nevertheless, $0$ is not initial: the arrows $\alpha'\gamma,\beta'\gamma:0\rightrightarrows2$ are distinct. Likewise, $3$ is not terminal, since $r,s:1\rightrightarrows3$ are distinct.

The same calculation locates the exact failure of PS 1. A commuting square for the cospan $1\xrightarrow{\alpha'}2\xleftarrow{\beta'}1$ would require two arrows from an object into $1$. Their common source can only be $1$ or $0$, and both arrows would be respectively $1_1$ or $\gamma$. Neither choice makes the square commute. Thus the left PS 1 fails, and the pair $\alpha',\beta'$ also has no equalizing arrow on the left. Dually, a commuting cocone over $3\xleftarrow{\alpha}2\xrightarrow{\beta}3$ can only use identity arrows out of $3$, so the right PS 1 fails; the pair $\alpha,\beta$ cannot be equalized on the right either. The category is consequently neither locally cofiltered nor locally filtered, although its associated poset satisfies both conditions. This is the point of the paragraph immediately following the example.

The crossed composition relations also admit a useful homotopical interpretation which the manuscript does not develop.

\textbf{The homotopy type of the finite example.}\ours
\textit{For the category presented above, $|N X|\simeq\mathbb{R}P^2$. In particular, its contractible reachability poset does not determine even the fundamental group of its nerve.}

\textit{Proof.} Let $C$ be the full subcategory on $1,2,3$. The functor $q:X\to C$ sends $0$ to $1$, sends $\gamma$ to $1_1$, and fixes the remaining generators. It respects the defining relations. For the inclusion $j:C\hookrightarrow X$, one has $qj=1_C$ and a natural transformation $1_X\to jq$ whose component at $0$ is $\gamma$ and whose other components are identities. The resulting homotopy of nerves fixes $N C$, so $|N X|\simeq|N C|$.

Identify each set $C(i,j)$, for $1\le i<j\le3$, with $\mathbf Z/2$: label $\alpha',\alpha,r$ by $0$ and $\beta',\beta,s$ by $1$. The two displayed relations say exactly that composition adds these labels. Let $P$ be the poset with objects $(i,\varepsilon)$, where $i\in\{1,2,3\}$ and $\varepsilon\in\mathbf Z/2$, with $(i,\varepsilon)<(j,\delta)$ precisely when $i<j$. Its involution changes every $\varepsilon$ to $\varepsilon+1$. There is an invariant functor $P\to C$ sending $(i,\varepsilon)$ to $i$ and an arrow to the label $\delta-\varepsilon$. A chain of arrows in $C$ lifts uniquely after choosing the lift of its first vertex, and the two choices are exchanged by the involution. This proves, degree by degree and compatibly with faces and degeneracies,
\[
 N P/(\mathbf Z/2)\simeq N C.
\]
The order complex of $P$ is the join of three two-point sets. Sending their vertices to $\{\pm e_1,\pm e_2,\pm e_3\}$ identifies it with the boundary of the octahedron. Radial projection identifies its realization with $S^2$, and the involution with the antipodal map. Realization commutes with the quotient, giving
\[
 |N C|\simeq |N P|/(\mathbf Z/2)\cong S^2/(x\sim-x)=\mathbb{R}P^2.
\]
Since the sphere is simply connected, this double covering also gives $\pi _1(|N X|)\cong\mathbf Z/2$.\hfill$\square$

The join and antipodal descriptions are classical \citep[p.~9 and Example 1.43]{Hatcher2002}. Our addition is their explicit identification with this particular presentation. The bounded literature searches did not locate that identification; the result is marked as our calculation, without a claim of priority.
\end{dcomment}


\msp{page 17}

\noindent Seule la notion a) (faiblement filtrante (\footnote[22]{pseudo-filtrant.}) à gauche) ne dépend que de $\operatorname{Ord} X$ - mais bien sûr si $X$ est \textit{pré}ordonnée, toutes trois ne dépendent que de $\operatorname{Ord} X$. Notons que dire que ça ne dépend que de $\operatorname{Ord} X$, ou que de $\mathrm{Esp}\,\mathrm{sep}X$, revient au même, puisque les deux se déterminent mutuellement. La condition a) signifie que $\operatorname{Esp} X$, ou aussi $\operatorname{Top}(X)$, est \textit{irréductible.} On dira que $X$ est \textit{irréductible} (au lieu de pseudo-filtrante). Les trois autres notions ne peuvent s'exprimer en termes de $\operatorname{Esp} X$ ou de $\operatorname{Ord} X$ seul, mais elles s'expriment en termes du $\mathrm{Topos}\, X$. La condition b) ($X$ localement faiblement filtrante) signifie que $\operatorname{Top}(X)$ est \textit{localement irréductible,} et on dira aussi que $X$ est \textit{localement irréductible} (au lieu de localement faiblement filtrante). Ça signifie que pour tout $X'$ étale sur $X$ et 0-connexe, $X'$ est faiblement filtrante à gauche, i.e. est irréductible.

Reste à exprimer la condition filtrante elle-même en termes de $\operatorname{Top} X$. Je rappelle à ce propos :

\dgrblock{dgrb90}{Products of connected objects}
\begin{dcomment}{exposition}{p17:Reste à exprimer la condition filtrante elle-même}{Products of connected objects}
Grothendieck now asks for a condition on the topos which remembers the equalization of parallel arrows. Irreducibility alone cannot do this: it records whether two subobjects meet, whereas equalization concerns two maps with the same source. Proposition 2.3 first isolates the intermediate condition that the component functor preserve products. Throughout the following topos discussion, ``connected object'' means noninitial and connected, that is, $0$-connected; the topos is locally connected. Thus every object is the coproduct of its connected components, and $\pi _0$ is left adjoint to the constant-object functor.

Write $\mathcal E$ for the category of objects of the topos and let $S$ be the full generating subcategory of connected objects used in the text. The precise statement of 2.3(b) is
\[
 \pi _0(V\times U)\longrightarrow\pi _0(V)
 \quad\hbox{is bijective whenever $U$ is connected.}
\]
The edition prints $\pi _0(U)$ on the right, although the displayed projection has target $V$. This is a slip of the edition, retained in the source. Condition (a) implies (c) by decomposing $U$ and $V$ into connected components: products distribute over coproducts in a topos, and each resulting product is connected. Conversely, (c) applied to connected objects gives (a); the same decomposition proves the equivalence with (b).

The reduction to generators deserves the proof which the footnote requests. For $F\in\mathcal E$, let $S/F$ have objects $s\to F$, with $s\in S$. The canonical presentation
\[
 F\simeq\mathop{\rm colim}_{(s\to F)\in S/F}s
\]
holds in $\mathcal E$; it is the density presentation for the canonical site of generators. Applying $\pi _0$ gives $\pi _0(F)\simeq\pi _0(S/F)$, since every $s$ is connected. If $s\times t$ is connected for every $s,t\in S$, separate preservation of colimits by products gives
\[
 \pi _0(F\times G)
 \simeq\mathop{\rm colim}_{S/F\times S/G}\{*\}
 \simeq\pi _0(S/F)\times\pi _0(S/G).
\]
This proves (d)$\Rightarrow$(c). Equivalently, one may use the clopen-partition argument in our block on Lemma 2.9 below; this avoids any choice of a presentation.

Applying (c) to the terminal object shows that the diagonal of $\pi _0(1)$ is a bijection onto its square. Consequently $\pi _0(1)$ has at most one element. The nondegenerate case has exactly one, which explains the separate nonemptiness clause in Definition 2.4. Nothing in this argument annihilates higher homotopy groups. The inference to $W_\infty$-asphericity immediately following 2.3 is explicitly withdrawn in the manuscript's annotation; the stronger hereditary hypothesis introduced in 2.6 is what will support the cohomological argument.
\end{dcomment}


\noindent\textbf{Proposition 2.3.} \textit{Les conditions suivantes sur un topos $\sX$ sont équivalentes (quand $S$ (comme `site') désigne une partie génératrice de} $\mathrm{Fais}(\sX) = S^{\wedge}$, formée d'objets 0-connexes (\footnote[23]{donc $\sX$ est localement 0-connexe.}).\textit{)}
\begin{enumerate}[label={},leftmargin=3em]
\item[\textit{a)}] \textit{Si $U$, $V$ sont deux objets $0$-connexes de $\operatorname{Ob} \sX$, $U \times V$ est $0$-connexe.}
\item[\textit{b)}] \textit{Si $U$ est un objet $0$-connexe de $\sX^{\wedge}$, alors $U \longrightarrow e$ est `$W_0$-asphérique', i.e. pour tout $V$, la projection $V \times U \longrightarrow V$ induit $\pi_0(U \times V) \xrightarrow{\sim} \pi_0(U)$.}
\item[\textit{c)}] \textit{Si $U$, $V$} \texttt{[sont]} \textit{dans $\sX^{\wedge}$, alors}
\[
\pi_0(U \times V) \xrightarrow{\sim} \pi_0(U) \times \pi_0(V).
\]
\item[\textit{d)}] \textit{Comme a), avec $U$, $V$ dans $S$. (Mais on suppose les objets de $S$ $0$-connexes.)}
\end{enumerate}
\noindent(\footnote[24]{Il faudrait expliciter la démonstation avec celle de 2.9 ci-dessous (p.~25).}). % [sic] démonstation -> démonstration, letter r missing

\msp{page 18}

\noindent De plus, ces conditions impliquent que $\sX$ est connexe, donc 0-connexe si $\sX \neq \emptyset$. Cela implique de plus que $\sX$ est $W_{\infty}$-asphérique, cf. Pursuing Stacks §35 (Corollary p.~52) (\footnote[25]{\textbf{Faux,} loc. cit. ne dit rien de tel.}).

C'est essentiellement Proposition 1 de Pursuing Stacks §65, cf. loc. cit. p.~170-171 (\footnote[26]{Mais je m'y borne aux topos de la forme $\operatorname{Top}(A)$, avec $A$ dans Cat.}). Cf. aussi Pursuing Stacks, §35, pages 50-55.

\printedpage{12}
% head: A. Grothendieck : Dérivateurs, Ch. II, Édition des Universités de Montpellier II et Paris VII
% foot: Version 8 juillet 2001

\noindent\textbf{Définition 2.4.} On dit que $\sX$ est \textit{totalement $0$-connexe} s'il est localement connexe, non vide et satisfait les conditions équivalentes de 2.3. (NB c'est un cas particulier de la notion `totalement $W$-asphérique' de loc. cit. On a ainsi la notion de totalement $n$-connexe.)

Notons que pour un objet $X$ de Cat, dire que $\operatorname{Top}(X)$ est totalement 0-connexe revient à dire que $X \neq \emptyset$, et que pour $x, y \in \operatorname{Ob} X$, $x \times y$ dans $X^{\wedge}$ est 0-connexe, i.e. $X/x \times y$ est 0-connexe (\footnote[27]{Cela signifie aussi que $\mathrm{diag}_X : X \longrightarrow X \times X$ est co-cofinale, i.e. $X^o \longrightarrow X^o \times X^o$ est cofinale, et implique qu'il en est de même de $X \longrightarrow X^E$ pour tout ensemble \textit{fini} $E$.}). C'est une condition très forte. On dira aussi que $X$ est \textit{totalement $0$-connexe.} Notons ici

\dgrblock{dgrb91}{The ordinary sifted condition and its variance}
\begin{dcomment}{later}{p18:Notons que pour un objet}{The ordinary sifted condition and its variance}
The product condition has acquired a useful name, provided that we keep track of the opposite category. For $\widehat X=\operatorname{Fun}(X^{\mathrm{op}},\mathbf {Set})$, the left adjoint to the constant-presheaf functor is
\[
 \pi _0(F)=\mathop{\rm colim}_{X^{\mathrm{op}}}F
          =\pi _0(X/F).
\]
Thus $\widehat X$ is totally $0$-connected precisely when $X^{\mathrm{op}}$ is \emph{sifted}: its colimits in sets commute with finite products. In the ordinary categorical characterization, the indexing category is nonempty and its diagonal is final. For the conventions of these pages, this says that $X\to X\times X$ is co-cofinal. The comma category at $(x,y)$ consists of common cones $z\to x$, $z\to y$, and is exactly $X/(h_x\times h_y)$.

This characterization also explains why binary products suffice. Repeated use of finality of the diagonal gives finality of every finite nonempty diagonal; nonemptiness handles the nullary product. The terminology and its relation to algebraic theories are developed by \citet{AdamekRosickyVitale2010}. The distinction from filteredness is structural: filtered colimits preserve finite limits, while sifted colimits are required to preserve finite products. Equalizers are precisely the operation absent from the latter requirement.

There is a corresponding distinction in higher category theory, but one must not transfer a statement by changing its vocabulary alone. In the higher version the relevant comma objects must be weakly contractible, rather than merely nonempty and connected. The characterization of filtered colimits by preservation of finite limits in spaces is \citet[Proposition 5.3.3.3]{Lurie2009HTT}; sifted colimits and geometric realizations are treated in his \S5.5.8. Our present identification concerns \emph{ordinary} set-valued colimits. The higher homotopy obstruction which survives it is exhibited later in the commentary on 2.19--2.20.
\end{dcomment}


\noindent\textbf{Proposition 2.5.} \textit{Soit $\sX$ un topos. Si $\sX$ est totalement $0$-connexe, alors $\sX$ est irréductible. (En fait la conclusion vaut, tautologiquement, en supposant seulement que le produit de deux objets $0$-connexes de $\sX^{\wedge}$ est non vide, et il suffit même de le supposer pour des sous-objets de l'objet final $e_{\sX}$.) En particulier, si $\sX = \mathit{Top}(X)$, où $X$ est dans} Cat\textit{, alors $X$ est faiblement filtrante à gauche.}

Mais bien sûr, cela n'implique \textit{pas} que $\sX$ soit localement irréductible.

\msp{page 19}

\noindent Par exemple la catégorie $\boldsymbol{\Delta}$, ou toute petite catégorie stable par produits binaires, est totalement 0-connexe, mais $\boldsymbol{\Delta}$ n'est pas localement irréductible (les catégories induites $\boldsymbol{\Delta}/\boldsymbol{\Delta}^n$, pour $n > 0$, ne sont visiblement \textit{pas} irréductibles). (Mais rappelons que si $\sX = \operatorname{Top}(T)$, $T$ un espace topologique, alors $T$ irréductible implique que $\sX$ est aussi \textit{localement} irréductible. Donc $T$ totalement 0-connexe implique $\operatorname{Top}(T)$ globalement \textit{et} localement irréductible, i.e. pour tout $T'$ étale sur $T$, et 0-connexe, $T'$ est irréductible.)

\dgrblock{dgrb92}{The simplex category separates the two conditions}
\begin{dcomment}{examples}{p19:les catégories induites}{The simplex category separates the two conditions}
Grothendieck's example shows exactly where a global product condition loses the information needed locally. We spell it out because the word ``visibly'' hides two different computations. Let $\Delta$ have the nonempty finite ordinals $[n]=\{0,\ldots,n\}$ and monotone maps. For any small category $X$ with binary products, the product of two representables is representable; its category of elements has a terminal object. Such a nonempty $X$ is therefore totally $0$-connected. The separate assertion about $\Delta$ does not say that $\Delta$ has binary products: the ordered square $[1]\times[1]$ is not an ordinal.

For $\Delta$, the product $h_{[m]}\times h_{[n]}$ is the simplicial set $\Delta^m\times\Delta^n$. Its vertices are the pairs $(i,j)$. An edge joins $(i,j)$ to $(i+1,j)$ whenever the latter exists, and similarly joins $(i,j)$ to $(i,j+1)$. Hence its vertex graph is connected. Every simplex is joined in its category of elements to each of its vertices by a face map. The category of elements is consequently nonempty and connected, which verifies the generator test of 2.3 for all $m,n$.

Local irreducibility fails already over $[1]$. In $\Delta/[1]$, consider the two objects $[0]\xrightarrow{0}[1]$ and $[0]\xrightarrow{1}[1]$. A common predecessor would be a nonempty ordinal $[r]$ mapping to both copies of $[0]$, with equal composites to $[1]$. Those composites are respectively constant at $0$ and $1$, so equality is impossible. The two principal sieves have empty intersection. In particular, $\Delta/[1]$ is not irreducible, and $\Delta$ is not cofiltered. Equivalently, the parallel vertex maps $[0]\rightrightarrows[1]$ admit no equalizing map into $[0]$.

This example gives an explicit positive answer to the request in footnote 31: total $0$-connectedness does not imply local irreducibility. Its opposite also supplies an ordinary sifted category which is not filtered. Neither conclusion requires an argument about derived limits.
\end{dcomment}


\noindent\textbf{Théorème 2.6.} \textit{Soit $\sX$ un topos localement connexe, $S \subseteq \sX^{\wedge}$ une petite partie génératrice formée d'objets $0$-connexes. Conditions équivalentes (}sauf \textit{c, c$'$, qui semblent} (\footnote[28]{\textbf{sont.}}) \textit{plus faibles, cf. démonstration).)}
\begin{enumerate}[label={},leftmargin=3em]
\item[\textit{a)}] \textit{Tout topos induit $\sX/U$ sur $U$ $0$-connexe dans $\sX^{\wedge}$ est totalement $0$-connexe.}
\item[\textit{a$'$)}] \textit{Tout produit fibré $U \times_W V$ d'objets $0$-connexes dans $\sX^{\wedge}$ est $0$-connexe.}
\item[\textit{b)}] \textit{$\sX$ est localement irréductible, et ses composantes connexes sont totalement $0$-connexes.}
\item[\textit{c)}] \textit{Comme a), avec $U$ dans $S$, i.e. les $\sX/U$ ($U \in \operatorname{Ob} S$) sont totalement $0$-connexes.}
\item[\textit{c$'$)}] \textit{Comme a$'$), avec $U$, $V$, $W$ dans $S$.}
\item[\textit{d)}] \textit{$S$ est localement filtrante à gauche, i.e. $S$ satisfait les deux conditions}

PS$^*$ 1\ \ \textit{Tout diagramme}
\begin{tikzcd}[ampersand replacement=\&, row sep=tiny, column sep=small, baseline=(\tikzcdmatrixname.center)]
a \arrow[dr] \& \\
 \& x \\
b \arrow[ur] \&
\end{tikzcd}
\textit{dans $S$ s'insère dans}
\begin{tikzcd}[ampersand replacement=\&, row sep=tiny, column sep=small, baseline=(\tikzcdmatrixname.center)]
 \& a \arrow[dr] \& \\
c \arrow[ur, dotted] \arrow[dr, dotted] \& \& x \\
 \& b \arrow[ur] \&
\end{tikzcd}
\textit{commutatif.}

PS$^*$ 2\ \ \textit{Tout diagramme $a \rightrightarrows b$ s'insère dans un diagramme commutatif $c \longrightarrow a \rightrightarrows b$.}
\end{enumerate}
\noindent(\footnote[29]{NB On a équivalence de \tikz[baseline=(n.base)]\node[draw, rounded corners, inner sep=1.5pt](n){a), a$'$), b), d),}; qui impliquent c), c$'$), lesquels sont équivalents. Mais il semble (\footnotemark[30]) que c), c$'$) soient strictement plus faible. Cf. cependant les variantes e), e$'$) dans 2.12 (p.~28, 29).}). % [sic] plus faible -> plus faibles, adjective does not agree with the plural subject c), c')
\footnotetext[30]{Il s'avère. C'est bien vrai, c) et c$'$) n'impliquent \textit{pas} les autres conditions, cf. annotation marginale
% footnote 30 continues at the foot of printed page 13:
p.~28 (cas de $X = B_G$).}

\printedpage{13}
% head: A. Grothendieck : Dérivateurs, Ch. II, Édition des Universités de Montpellier II et Paris VII
% foot: Version 8 juillet 2001

\msp{page 20}

\noindent Ces conditions impliquent que les objets 0-connexes de $\sX^{\wedge}$ (en particulier les éléments de $S$) sont $\infty$-connexes.

\noindent\textbf{Corollaire 2.7} (de l'équivalence b $\Longleftrightarrow$ d). \textit{Soit $X$ un objet de} Cat. \textit{Pour que $X$ soit filtrant à gauche il faut et il suffit que $\mathit{Top}(X)$ soit totalement $0$-connexe et localement (et donc aussi globalement) irréductible (ou encore que $X$ soit totalement $0$-connexe et que pour tout $x \in \operatorname{Ob} X$, la catégorie localisée $X/x$ soit irréductible (elle est même totalement $0$-connexe), ou encore que pour deux objets $a$, $b$ dans $X$, $a \times b \in X^{\wedge}$ soit $0$-connexe, et pour}
\begin{tikzcd}[ampersand replacement=\&, row sep=tiny, column sep=small, baseline=(\tikzcdmatrixname.center)]
a \arrow[dr] \& \\
 \& x \\
b \arrow[ur] \&
\end{tikzcd}
\textit{dans $X$, $a \times_x b$ (dans $X^{\wedge}$) soit $0$-connexe)} (\footnote[31]{Exemple de $X$ totalement 0-connexe, mais non cofiltrant, i.e. non localement irréductible ? La catégorie $\boldsymbol{\Delta}$ des simplexes type est totalement 0-connexe, mais n'est pas cofiltrante car $\boldsymbol{\Delta}^0 = e \underset{\partial_1}{\overset{\partial_0}{\rightrightarrows}} \boldsymbol{\Delta}^1$ ne peut être égalisé ! Donc \texttt{[la catégorie opposée]} $\boldsymbol{\Delta}^o = I$ est telle que $I \longrightarrow I^E$ soit cofinal pour tout ensemble fini $E$, mais $I$ non filtrante.}).

\noindent\textsc{Démonstration} de 2.6. On a le diagramme d'implications $\pm$ tautologiques
\[
\begin{tikzcd}[ampersand replacement=\&, column sep=large, row sep=large]
 \& \mathrm{a} \arrow[r, Leftrightarrow] \arrow[dl, Rightarrow, "{\substack{\text{par}\\ \text{2.5}}}"'] \arrow[d, Rightarrow] \& \mathrm{a}' \arrow[d, Rightarrow] \\
\mathrm{b} \arrow[d, Rightarrow] \& \mathrm{c} \arrow[r, Leftrightarrow, "{(\Leftarrow \text{ par 2.3.})}"] \& \mathrm{c}' \arrow[dll, Rightarrow] \\
\text{PS 1.} \& \&
\end{tikzcd}
\]
Montrons quand même les deux implications c$'$) $\Longrightarrow$ PS 1 et b $\Longrightarrow$ PS 1.

\dgrblock{dgrb93}{What the hereditary condition adds}
\begin{dcomment}{exposition}{p20:Montrons quand même les deux implications}{What the hereditary condition adds}
Theorem 2.6 makes the product condition stable under passage to a connected object as base. Its surviving equivalence is
\[
 (a)\ \Longleftrightarrow\ (a')\ \Longleftrightarrow\ (b)
 \ \Longleftrightarrow\ (d),
 \qquad (c)\Longleftrightarrow(c'),
\]
with an implication from the first group to the second. Here the three objects $U,V,W$ in (a') are all connected. Connectedness of $W$ is essential: even in sets, two singleton objects mapping to different elements of a two-element set have empty pullback.

The equivalence of (a) and (a') is the same product test inside $\mathcal E/W$. Condition (a) implies local irreducibility by 2.5 and gives the assertion about components in (b). To see the decisive step (b)$\Rightarrow$(d), first use 2.8 for PS 1. Given $u,v:a\rightrightarrows b$ in $S$, both objects lie over the same connected component of $1$. By (b), $b\times b$ is connected. The equalizer is the pullback
\[
 E=a\times_{b\times b}b,
 \qquad a\longrightarrow b\times b=(u,v),\quad
 b\longrightarrow b\times b=\operatorname{diag}.
\]
Local irreducibility, applied over the connected object $b\times b$, makes $E$ noninitial. A generator mapping to $E$ supplies an equalizing arrow, which is PS 2. The text's ``c) est vérifié'' at manuscript p.~24 means (b); its reference to $\mathcal X\times\mathcal X$ should be read as an application of 2.8 in $\mathcal X$ itself.

There is a concise exactness interpretation after restricting to any connected component of the topos. Under (a'), decomposition into connected components gives, for arbitrary $U,V,W$ in that component,
\[
 \pi _0(U\times_WV)
 \simeq\pi _0(U)\times_{\pi _0(W)}\pi _0(V).
\]
Indeed, the summands of the pullback are precisely the products $U_i\times_{W_k}V_j$ for which both factors map to the same component $W_k$; each is connected by (a'). Since $\pi _0(1)=\{*\}$, $\pi _0$ preserves finite limits on this connected topos. We shall use this stronger fact in the cohomological argument. For a disconnected topos one applies it component by component; its global $\pi _0$ does not preserve the terminal object.
\end{dcomment}


\msp{page 21}

\noindent C'est contenu dans le

\noindent\textbf{Corollaire 2.8.} \textit{Soient $\sX$, $S$ comme dans 2.6. Conditions équivalentes :}
\begin{enumerate}[label={},leftmargin=3em]
\item[\textit{a)}] \textit{$\sX$ est localement irréductible, i.e. les topos induits $\sX/U$, où $U$ est $0$-connexe, sont irréductibles.}
\item[\textit{b)}] \textit{Tout produit fibré $U \times_W V$ d'objets $0$-connexes $U$, $V$ est `non vide'.}
\item[\textit{b$'$)}] \textit{Tout produit fibré $U \times_W V$, avec $U$, $V$ `non vide', $W$ $0$-connexe, est non vide.}
\item[\textit{c)}] \textit{Les $\sX/x$ avec $x \in S$ sont irréductibles.}
\item[\textit{d)}] \textit{Comme b), avec $U$, $V$, $W$ dans $S$.}
\item[\textit{e)}] \textit{$S$ est localement faiblement filtrante à gauche, i.e. satisfait PS$^*$ 1.}
\end{enumerate}
\noindent Ici les implications $\pm$ évidentes sont

% ---- from Der-ChII-234_p016-020.tex
\printedpage{14}
% head: A. Grothendieck : Dérivateurs, Ch. II, Édition des Universités de Montpellier II et Paris VII
% foot: Version 8 juillet 2001

\[
\begin{tikzcd}
\mathrm{b}' \arrow[r, Rightarrow] & \mathrm{b} \arrow[rr, Rightarrow] \arrow[d, Rightarrow] & & \mathrm{a} \arrow[d, Rightarrow] \\
 & \mathrm{d} \arrow[r, Leftrightarrow] & \mathrm{e} \arrow[r, Leftrightarrow] & \mathrm{c}.
\end{tikzcd}
\]

Prouvons les implications de la ligne du bas. Considérons d). Si $\begin{tikzcd}[cramped, row sep=0.1em, column sep=0.8em, baseline=(current bounding box.center)] a \arrow[dr] & \\ & x \\ b \arrow[ur] & \end{tikzcd}$ \texttt{[est]} dans $S$, dire que le produit fibré dans $\X$ est `non vide' signifie, puisque $S$ est un système de générateurs et en particulier \textit{non vide,} qu'il existe $c$ dans $S$ et $c\longrightarrow a\times_x b$, ou ce qui revient au même, $\begin{tikzcd}[cramped, row sep=0.1em, column sep=0.8em, baseline=(current bounding box.center)] & a \\ c \arrow[ur, dotted] \arrow[dr, dotted] & \\ & b \end{tikzcd}$, rendant commutatif le diagramme $\begin{tikzcd}[cramped, row sep=0.1em, column sep=0.8em, baseline=(current bounding box.center)] & a \arrow[dr] & \\ c \arrow[ur, dotted] \arrow[dr, dotted] & & x \\ & b \arrow[ur] & \end{tikzcd}$. Donc d) équivaut à PS 1, i.e. à e) ($\X$ localement faiblement filtrant à gauche). Considérons la condition c) : les $\X/a=(S/a)^\sim$ sont irréductibles. Cela signifie que pour deux sous-objets \textit{non vides} $U,V\subseteq a$ dans $\X^\wedge$, on a $U\cap V\neq\emptyset$, ou encore que pour deux cribles

\msp{page 22}$\underline{U}$, $\underline{V}$ de $S/x$, si $\underline{U}$, $\underline{V}$ non vides (ce qui signifie que les sous-objets de $x$ dans $\X^\wedge$ qu'ils définissent sont `non vides', ou que les objets de $S$ sont `non vides' dans $\X^\wedge$), alors $\underline{U}\cap\underline{V}$ est non vide. On peut supposer que $\underline{U}$, $\underline{V}$ est engendré chacun par une seule flèche $a\xrightarrow{u}x$, $b\xrightarrow{v}x$. Le crible intersection est formé des objets $c\longrightarrow x$ de $S/x$ tels que $c\longrightarrow x$ se factorise par un $\begin{tikzcd}[cramped, row sep=0.1em, column sep=0.8em, baseline=(current bounding box.center)] & a \\ c \arrow[ur, dotted] \arrow[dr, dotted] & \\ & b \end{tikzcd}$, donc dire que ce crible n'est pas vide, signifie que PS 1 est satisfait pour les deux flèches données $u$, $v$. Donc c) équivaut à la condition PS 1\texttt{[, i.e. à]} e.

Ainsi on a établi
\[
\mathrm{b}'\Longrightarrow\mathrm{b}\Longrightarrow\mathrm{a}\Longrightarrow(\mathrm{c}\Longleftrightarrow\mathrm{d}\Longleftrightarrow\mathrm{e}),
\]
reste à prouver que c $\Longrightarrow$ b$'$. Soit donc $\begin{tikzcd}[cramped, row sep=0.1em, column sep=0.8em, baseline=(current bounding box.center)] U \arrow[dr] & & V \arrow[dl] \\ & W & \end{tikzcd}$ un diagramme d'objets `non vides', avec $W$ 0-connexe, il faut prouver que le produit fibré est non vide. Considérant $a\longrightarrow U$ et $b\longrightarrow V$ avec $a,b\in S$, on est ramené au cas où $U=a$, $V=b$ dans $S$. Prouvons que si $a\times_W b$ est `vide', alors $W$ est disconnexe. Pour ceci, considérons la sous-catégorie $C$ de $S/W$ formée des $x\longrightarrow W$ tels que $a\times_W x$ soit `non vide'. Désignant par $w(a)$, $w(x)$ l'image de $a$, $x$ dans $W$, cela signifie que $w(a)\cap w(x)\neq\emptyset$.

\msp{page 23}Notons que tout quotient dans $\X^\wedge$ d'un objet irréductible est irréductible, donc $w(a)$, $w(x)$ sont des objets irréductibles. Cela montre que pour $x$ dans $C$ et $x'$ dans $S/W$, on a $x'$ dans $C$, i.e. $w(a)\cap w(x')\neq\emptyset$ si et seulement si $w(x)\cap w(x')\neq\emptyset$. (Lemme direct : dans un objet $W$ d'un topos $\X$, si on a \texttt{[des]} sous-objets $W_1$, $W_2$, $W_3$ irréductibles, si l'un des deux \texttt{[plutôt trois]} rencontre les deux autres, ceux-ci se rencontrent.) Cela montre que $C$ est un crible, puisque $x'\longrightarrow x$ implique $w(x')\subseteq w(x)$, donc $w(x')$ rencontre $w(x)$, donc aussi $w(a)$. Mais alors le co-crible $C'=(S/W)\backslash C$ est également un crible, car si $c$
\printedpage{15}
% head: A. Grothendieck : Dérivateurs, Ch. II, Édition des Universités de Montpellier II et Paris VII
% foot: Version 8 juillet 2001
est dans $C'$, i.e. \texttt{[si]} $w(c)$ ne rencontre pas $w(x)$, a fortiori pas tout $x'$ au-dessus de $x$, pour lequel on aura $w(x')\subseteq w(x)$. Ainsi, $S/W$ est décomposé en somme de deux ouverts $C$, $C'$, non vides puisque l'un contient $a$, l'autre $b$. Par suite, $W$ est décomposé en somme de deux sous-objets non vides, donc $W$ est non 0-connexe.

On a prouvé 2.8, donc aussi (revenant à 2.6) les implications vers PS 1 dans le diagramme page 20. Il reste à prouver
\[
\left\{
\begin{aligned}
&\mathrm{b}\Longrightarrow\mathrm{PS}\ 2,\ \mathrm{c}'\Longrightarrow\mathrm{PS}\ 2\\
&\mathrm{PS}\ 1+\mathrm{PS}\ 2\ (\text{i.e. d})\Longrightarrow\mathrm{a}'
\end{aligned}
\right.
\]
(\footnote[32]{NB finalement on ne prouve pas c$'\Longrightarrow$ PS 2.}).

\dgrblock{dgrb94}{How irreducible generators control an arbitrary slice}
\begin{dcomment}{exposition}{p23:On a prouvé 2.8, donc aussi}{How irreducible generators control an arbitrary slice}
The proof of 2.8 explains why irreducibility, although weaker than filteredness, is genuinely local in the topos. The missing hypothesis in its printed condition (b) is that $W$ be connected. The assertion used in the proof is (b'): noninitial $U$ and $V$ over a connected $W$ have noninitial pullback.

Here is the passage from the generating objects to such a $W$. Suppose every $s\in S$ is irreducible, and consider the images $W_s\subseteq W$ of all maps $s\to W$. These images cover $W$ and remain irreducible: inverse images of two nonzero subobjects under an epimorphism are nonzero, and their nonzero intersection has nonzero image. On this covering family, declare two images related when they meet. The relation is reflexive and symmetric. It is transitive because, if $A$ and $C$ both meet an irreducible $B$, the two nonzero subobjects $A\cap B$ and $B\cap C$ of $B$ meet. Thus the equivalence classes give pairwise disjoint unions of subobjects covering $W$.

These unions are a coproduct decomposition in the topos. Since $W$ is connected, there is only one class. Every pair of generating images therefore meets. Choose a generator mapping into each of the noninitial objects $U$ and $V$. Their images in $W$ meet, and epimorphisms remain epimorphisms after pullback; consequently $U\times_WV$ is noninitial. This proves the step (c)$\Rightarrow$(b') without assuming that every object of a site is representable as a presheaf on that site.

Finally, the condition that $s\times_x t$ be noninitial is equivalent to the existence of a generator $r$ mapping to that pullback. Such a map is exactly a commutative square completing $s\to x\leftarrow t$, namely PS 1. This proves the equivalence with (d) and (e). The implication (b')$\Rightarrow$(a) follows by taking two nonzero subobjects of a connected object. The remaining implications are restrictions of this same test. Grothendieck's argument is therefore sound after restoring the condition on the base; the diagram of implications is not concealing an additional homotopical assumption.
\end{dcomment}


\msp{page 24}Nous nous plaçons donc dans l'hypothèse où PS 1 est satisfait (puisque b), c$'$) et d) impliquent PS 1), donc où on est sous les conditions équivalentes de 2.8 - $\X$ est localement irréductible. Considérons alors la condition PS 2, et prouvons qu'elle est impliquée tant par b), que par c$'$), et qu'elle implique a$'$). Elle signifie que pour une double flèche $a\overset{u}{\underset{v}{\rightrightarrows}}b$ dans $S$, le conoyau \texttt{[plutôt noyau]} dans $S^\wedge$, ou, ce qui revient au même, dans $\X^\wedge=S^\sim\subseteq S^\wedge$, est `non vide'. On peut l'écrire sous la forme d'un produit fibré
\[
\begin{tikzcd}
a \arrow[r, "{\scriptstyle (u,v)}"] & b\times b \\
\operatorname{Ker}(u,v) \arrow[u] \arrow[r] & b. \arrow[u, "\mathrm{diag}"']
\end{tikzcd}
\]
Ici par hypothèse $a$ et $b$ sont des objets 0-connexes de $\X^\wedge$ (ceux de $S$ l'étant), c) est vérifié (($\X$ localement irréductible et) ses composantes connexes \texttt{[sont]} totalement 0-connexes), alors $b\times b$ est 0-connexe comme produit de deux objets 0-connexes sur une \textit{même composante connexe,} laquelle est totalement 0-connexe. Comme $\X\times\X$ est localement irréductible, il s'ensuit que le produit fibré est non vide (2.8), donc b) $\Longrightarrow$ d). Prouvons d) $\Longrightarrow$ a), i.e. si

\msp{page 25}$S$ est localement cofiltrante, alors les topos induits $\X/W$, où $W$ est 0-connexe, sont totalement 0-connexes. Ce topos induit est décrit par le site $S/W$, i.e. $S/W$ est un système de générateurs. Or si $S$ est localement filtrante, $S/W$ (étale sur $S$) l'est aussi. Donc on peut supposer $W=e$, et à prouver que si $S$ est localement cofiltrante et $\X$ 0-connexe, alors $\X$ est totalement 0-connexe. Notons que $\X$ 0-connexe équivaut à $S$ 0-connexe - il faudrait quand même l'expliciter :

\textbf{Lemme 2.9.} \textit{Soit $\X$ un topos, $S\subseteq\X^\wedge$ une sous-catégorie pleine génératrice formée d'objets} 0-connexes. \textit{Alors $\X$ est $0$-connexe si et seulement si $S$ l'est - plus généralement, l'ensemble ordonné des sous-objets à la fois ouverts et fermés de $S$ dans} Cat \textit{(ou encore des
\printedpage{16}
% head: A. Grothendieck : Dérivateurs, Ch. II, Édition des Universités de Montpellier II et Paris VII
% foot: Version 8 juillet 2001
sous-objets ouverts $=$ cribles de $S$ qui sont aussi fermés, i.e. des cocribles) est isomorphe à celui des sommandes directs de $\X$.}

Je l'ai utilisé tacitement dans la démonstration de 2.8, pour le topos induit $\X/W$. Je n'explicite pas la démonstration.

\dgrblock{dgrb95}{Connectedness in a category of generators}
\begin{dcomment}{exposition}{p25:Je n'explicite pas la démonstration}{Connectedness in a category of generators}
Lemma 2.9 is what permits the categorical and topos-theoretic arguments to meet. Its content is stronger than the statement that a generating family covers the terminal object: a partition of the generators must descend to an actual coproduct decomposition of that object.

Let $S\subseteq\mathcal E$ be full, generating, and formed of connected objects. A decomposition $1=A\amalg B$ assigns every $s\in S$ entirely to one summand, since the two pullbacks decompose $s$. Every arrow preserves this assignment. It thus gives a partition $S=S_A\amalg S_B$ into full subcategories with no arrows between them. Conversely, suppose such a partition is given. Let $A$ and $B$ be the unions in $1$ of the images of the objects in the respective parts. They cover $1$, because $S$ generates. If $A\cap B$ were noninitial, some generator would map into it. Pulling back the epimorphic families defining $A$ and $B$, and choosing a generator in a noninitial member, would produce an object $t\in S$ with maps to an object of $S_A$ and one of $S_B$. This contradicts the partition. Hence $A\cap B=0$, and $A\amalg B\simeq1$.

The two constructions are inverse: a connected generator whose image lies in $A$ cannot belong to the other part, and the images of all generators assigned to $A$ jointly cover $A$. The same argument works with any set of parts and proves
\[
 \pi _0(S)\simeq\pi _0(1_{\mathcal E}).
\]
It also applies to $S/F$, which is a full generating subcategory of $\mathcal E/F$. An object $s\to F$ is connected in the slice exactly when its domain $s$ is connected, since coproducts in a slice are computed on domains. We obtain the repeatedly used identity
\[
 \pi _0(S/F)\simeq\pi _0(F).
\]
This supplies the proof omitted at manuscript p.~25 and justifies the earlier use of the lemma in 2.8. It also explains why the category of elements, rather than only the set of objects of a site, controls connectedness.
\end{dcomment}


Ainsi (revenant à la démonstration de d) $\Longrightarrow$ c) dans 2.6.) $S$ étant 0-connexe et localement cofiltrante,

\msp{page 26}\texttt{[elle]} est en fait cofiltrante. Donc pour deux objets $a$, $b$ de $S$, $\exists\ c$ avec $c\longrightarrow a$, $c\longrightarrow b$ (condition PS 1$'$, plus forte que PS 1). Pour prouver que $\X$ est totalement 0-connexe, on est réduit à prouver par 2.3 que les $a\times b$ dans $\X^\wedge$ sont 0-connexes. Par 2.9, il revient au même de prouver qu'ils sont 0-connexes dans $S^\wedge$. Donc on est ramené à prouver le

\textbf{Lemme 2.10.} \textit{Soit $S$ une catégorie cofiltrante. Si $a$, $b$ sont deux objets de $S$, $a\times b$ est un objet} $0$\textit{-connexe de $S^\wedge$.}

On voit de suite qu'il suffit de prouver (et en fait il revient au même, compte tenu que $a\times b$ est localement irréductible) que si on a deux flèches
\[
\begin{tikzcd}
x \arrow[r] & a\times b \\
 & y, \arrow[u]
\end{tikzcd}
\]
alors $x\times_{a\times b}y\neq\emptyset$, i.e. $\exists\ z$ dans $S$ rendant commutatif
\[
\begin{tikzcd}
x \arrow[r] & a\times b \\
z \arrow[u] \arrow[r] & y. \arrow[u]
\end{tikzcd}
\]
Or par PS 1$'$ il existe bien un $z$ s'envoyant dans $x$ et dans $y$, mais les deux flèches composées
\[
z\rightrightarrows a\times b,
\]
i.e. chacune des double flèches
\[
z\rightrightarrows a \tag*{$(*)$}
\]
\[
z\rightrightarrows b \tag*{$(*')$}
\]
peuvent n'être pas égales. Mais par PS 2,

\msp{page 27}on peut trouver
\[
z'\longrightarrow z
\]
\printedpage{17}
% head: A. Grothendieck : Dérivateurs, Ch. II, Édition des Universités de Montpellier II et Paris VII
% foot: Version 8 juillet 2001
égalisant la double flèche $(*)$, donc on peut supposer que celle-ci est déjà triviale. Appliquant encore une fois PS 2, par un $z''\longrightarrow z'$, on égalise de même la double flèche $(*')$, et on gagne. On a donc prouvé, finalement,
\[
\begin{tikzcd}
\mathrm{a} \arrow[r, Leftrightarrow] \arrow[d, Rightarrow] & \mathrm{a}' \arrow[r, Rightarrow] & \mathrm{b} \arrow[r, Leftrightarrow] & \mathrm{d}\ (=\mathrm{PS}\ 1+\mathrm{PS}\ 2) \\
\mathrm{c} \arrow[r, Leftrightarrow] & \mathrm{c}',
\end{tikzcd}
\]
mais pas l'implication inverse. En d'autres termes, si les topos induits $\X/a$, pour $a$ dans $S$, sont totalement 0-connexes, je n'ai pas prouvé qu'il en est de même des $\X/W$, $W$ un objet 0-connexe \textit{quelconque} de $\X^\wedge$. Ainsi, on en est au point suivant

\dgrblock{dgrb96}{Where both equalizations are used}
\begin{dcomment}{exposition}{p27:mais pas l'implication inverse}{Where both equalizations are used}
The completion of 2.6 is a useful instance of the finite-cone argument underlying filteredness. It is not enough to find a common predecessor of two objects: one must also arrange that the resulting maps agree in both factors. This is why the proof of Lemma 2.10 invokes PS 2 twice.

Suppose $S$ is cofiltered. An object of $S/(h_a\times h_b)$ is a triple $(x,r:x\to a,s:x\to b)$. It exists by the common-predecessor axiom. Given two such triples with domains $x$ and $y$, choose $z\to x$ and $z\to y$. Equalize the two resulting maps $z\rightrightarrows a$ by a map $z'\to z$, then equalize the two maps $z'\rightrightarrows b$ by $z''\to z'$. The second operation preserves the first equality. The object with domain $z''$ is a common predecessor in $S/(h_a\times h_b)$. This proves its connectedness, as required in 2.10. In fact it is cofiltered: parallel arrows in it can be equalized after projection to $S$, and the resulting arrow still belongs to the category of elements.

Now let $S$ merely be locally cofiltered, and let $W$ be connected in the topos. The projection $S/W\to S$ is a discrete fibration. PS 1 lifts along a discrete fibration: once a common predecessor is chosen downstairs, lifting one leg determines the object upstairs, and equality of the two composites gives the other leg. PS 2 lifts for the same reason, using uniqueness of lifts. Thus $S/W$ is locally cofiltered. Lemma 2.9 says it is connected, so it is cofiltered. Applying the preceding argument in this generating category and then the generator test of 2.3 shows that $\mathcal E/W$ is totally $0$-connected. This is (d)$\Rightarrow$(a).

The proof establishes the announced equivalence of (a), (a'), (b), and (d). It does not establish (c')$\Rightarrow$PS 2, and the annotation correctly removes that assertion. Restricting the base of a pullback to a representable object has lost information: the object $b\times b$ required for the diagonal argument need not be among the chosen generators. The following group example makes this loss explicit.
\end{dcomment}


\textbf{Perplexité 2.11.} \textit{Soit $X$ dans} Cat. \textit{Si $X$ est localement filtrante à gauche, alors les produits fibrés $a\times_x b$ d'objets dans $X$ sont des objets} $0$\textit{-connexes de $X^\wedge$, et cette condition à son tour implique PS 1, i.e. que $X$ est localement irréductible. Mais en présence de cette hypothèse PS 1, ou $X$ localement irréductible, la condition sur les produits fibrés $a\times_x b$ équivaut à la suivante : pour tout diagramme} \texttt{[en]} \textit{traits pleins}

\msp{page 28}
\[
\begin{tikzcd}
c \arrow[r] & a\times_x b \\
e \arrow[u, dotted] \arrow[r, dotted] & d, \arrow[u]
\end{tikzcd}
\qquad\qquad
\textit{avec $x$, $a$, $b$, $c$, $d$ dans $X$,}
\]
\textit{on peut compléter en un diagramme} commutatif, \textit{avec $e$ dans $X$. En d'autre termes, tout diagramme commutatif à traits pleins ci-dessous se complète en y rajoutant un `objet initial' $e$} % [sic] d'autre termes -> d'autres termes, missing plural s
\[
\begin{tikzcd}[column sep=small, row sep=small]
 & & a \arrow[ddr] & \\
 & c \arrow[ur] \arrow[dddr] & & \\
e \arrow[ur, dotted] \arrow[dr, dotted] & & & x \\
 & d \arrow[uuur] \arrow[dr] & & \\
 & & b \arrow[uur] &
\end{tikzcd}
\]
(\footnote[33]{Il suffit d'exiger pour $e$ la commutativité des deux carrés $(e\ c\ b\ d)$ et $(e\ d\ a\ c)$.}).

Cette condition est donc impliquée par $X$ localement cofiltrante, et elle est satisfaite aussi par ce qui précède plus généralement si et seulement si (\footnote[34]{si et seulement si ?}) les $X/x$ sont cofiltrantes (\footnote[35]{Elle signifie que $X$ satisfait PS 1, et la condition

\textbf{PS$'$ 2} Tout diagramme commutatif
\[
a\rightrightarrows b\longrightarrow x
\]
% footnote continued on printed page 18
s'insère dans un diagramme commutatif
\[
y\longrightarrow a\rightrightarrows b\longrightarrow x.
\]
}). Mais la réciproque est-elle vraie - cette condition de totale 0-connexité des $X/x$ implique-t-elle que $X$ soit totalement 0-connexe ? Mais je n'ai pas construit d'exemple où les $X/x$
\printedpage{18}
% head: A. Grothendieck : Dérivateurs, Ch. II, Édition des Universités de Montpellier II et Paris VII
% foot: Version 8 juillet 2001
sont cofiltrantes sans que $X$ soit localement cofiltrante (ou ce qui revient au même, de $X$ 0-connexe avec les $X/x$ cofiltrantes sans que $X$ en soit cofiltrante)(\footnote[36]{Mais il y a un \textit{contre-exemple} évident, si $X$ est un groupoïde 0-connexe, les $X/x$ sont $\approx e$, donc cofiltrantes, mais $X$ n'est cofiltrant (i.e. \texttt{[mot(s) illisible(s)]}) que si $X$ est 1-connexe. P.ex. si $C$ est un groupoïde $\neq 1$, alors $X=B_G$ est tel que les $X/x$ sont cofiltrantes, mais $X$ non cofiltrant. % [sic] si $C$ est un groupoïde -> si $G$ est un groupoïde, the groupoid is then written $B_G$
}) ; ni d'exemple où les $X/x$ sont totalement 0-connexes, sans que les $X/x$ soient cofiltrantes. Mais notons le

\dgrblock{dgrb97}{A groupoid invisible to its representable slices}
\begin{dcomment}{examples}{p28:Mais la réciproque est-elle vraie}{A groupoid invisible to its representable slices}
The marginal answer to Perplexity 2.11 is decisive, and it also explains the weaker condition that remains. Let $G$ be a nontrivial group and let $BG$ be its one-object category. The objects of $BG/*$ are the elements $g\in G$, viewed as arrows to the unique object. For any $g,h$, there is exactly one arrow from $g$ to $h$, namely the element $k$ satisfying $hk=g$. Hence $BG/*$ is a nonempty indiscrete groupoid, equivalent to the terminal category, and is cofiltered.

The category $BG$ itself is not cofiltered. If $g\ne h$ are parallel arrows, an equalization $gk=hk$ would imply $g=h$ because $k$ is invertible. Its presheaf topos is the category of right $G$-sets. The representable is the regular $G$-set $G$; it is connected. Its square with diagonal action has
\[
 \pi _0(G\times G)\simeq G,
 \qquad [(a,b)]\longmapsto ba^{-1},
\]
after choosing the corresponding convention for the right action. Thus the square is not connected. This verifies directly that the topos is locally irreducible but is not totally $0$-connected, although the slice over the representable is equivalent to sets. The image of the regular object onto the terminal $G$-set is exactly the quotient missed by the generator-only test.

There is, however, no further example in which \emph{every} $X/x$ is totally $0$-connected but some $X/x$ is not cofiltered. Indeed, total $0$-connectedness of all $X/x$ makes all these categories irreducible, so $X$ is locally irreducible. For $y\to x$, the slice $(X/x)/(y\to x)$ is isomorphic to $X/y$, and is therefore irreducible. Hence $X/x$ is both totally $0$-connected and locally irreducible. Corollary 2.7, whose proof has just been completed, makes it cofiltered.

In elementary terms, the surviving weaker condition is PS 1 together with
\[
 a\mathrel{\substack{\xrightarrow{u}\\[-.5ex]\xrightarrow[v]{}}}b
 \xrightarrow{h}x,\quad hu=hv
 \quad\Longrightarrow\quad
 \exists\,t:c\to a\ \text{with }ut=vt.
\]
This is PS$'$ 2 in footnote 35. It equalizes parallel arrows only when they already become equal after a common outgoing arrow. In a groupoid that antecedent forces equality, so the condition is automatic. PS 2 has no such antecedent. The two questions in the paragraph are therefore settled differently: the first has the groupoid counterexample already supplied by Grothendieck; the second has a negative existence answer from his preceding theorem.
\end{dcomment}


\textbf{Corollaire 2.12 :} \textit{Les conditions équivalentes a), a$'$), b), d) (à l'exclusion de c), c$'$), plus faibles apparemment) du théorème 2.6}

\msp{page 29}\textit{équivalent encore à la suivante}
\begin{itemize}[leftmargin=2.5em]
\item[\textit{e)}] \textit{$\forall\ x$ dans $S$, le topos induit $\X/x$ et le topos induit $\X/x\times x$ est irréductible} (\footnote[37]{NB $\X/x$ inutile.}),
\end{itemize}
\textit{où encore à la suivante} % [sic] où encore -> ou encore, conjunction "ou" without accent
\begin{itemize}[leftmargin=2.5em]
\item[\textit{e$'$)}] \textit{$\forall\ x$ dans $S$, $S^\wedge/x$ et $S^\wedge/x{\times}x$ sont irréductibles, i.e. $S/x$ et $S/x{\times}x$ sont irréductibles} (\footnote[38]{$S/x$ inutile, car $x\xrightarrow{\mathrm{diag}}x\times x$ est une immersion ouverte non vide, donc si $x\times x$ est irréductible, $x$ l'est aussi. Mais il revient au même de dire que les $S/x$ sont irréductibles (condition PS$^*$ 1, cf. 2.8, p. 21) \textit{et} que les $x\xrightarrow{\mathrm{diag}}x\times x$ sont dominants (ce que implique en effet que $x\times x$ est irréductible). Ceci redonne immédiatement la forme PS$^*$ 2 de la condition cofiltrante. % [sic] ce que implique -> ce qui implique, relative pronoun subject
}).
\end{itemize}
C'est la forme `topologique' la plus faible en apparence que je peux trouver pour la condition $S$ localement cofiltrante, i.e. $\X$ (ou Top$(S)$, au choix) localement totalement 0-connexe.

\dgrblock{dgrb98}{Why the square of one generator detects equalizers}
\begin{dcomment}{exposition}{p29:C'est la forme `topologique' la plus faible}{Why the square of one generator detects equalizers}
Corollary 2.12 is not merely another presentation of a long list of equivalent conditions. It isolates the diagonal as the part of the product test which forces equalization. We can read its criterion entirely in the topos, without treating the chosen site as a space.

Suppose that $s\times s$ is irreducible for every $s\in S$. The diagonal $s\hookrightarrow s\times s$ is a noninitial subobject. Every noninitial subobject of an irreducible object is irreducible, so each $s$ is irreducible and 2.8 gives PS 1. For two maps $u,v:a\rightrightarrows b$ in $S$, the image of $(u,v):a\to b\times b$ is noninitial. It meets the diagonal, because $b\times b$ is irreducible. Pullback stability of epimorphisms then says that
\[
 \operatorname{Eq}(u,v)=a\times_{b\times b}b
\]
is noninitial. A generator mapping into this equalizer supplies PS 2. Conversely, PS 1 and PS 2 give local cofilteredness; the category of generators over $b\times b$ is cofiltered by the proof of 2.10, so this product is irreducible.

This also proves the comment that the separate hypothesis on $\mathcal E/s$ is redundant. The diagonal is a nonzero open subobject in the language of slice topoi, and irreducibility passes to such an open. Saying instead that $s$ is irreducible and its diagonal is dense gives the same test: density means that the diagonal meets the image of every noninitial object mapping to $s\times s$. The pair $(u,v)$ then gives the equalizer test just used. The relevant ``open immersion'' is the morphism of topoi associated with a monomorphism in $\mathcal E$; it is not an assertion that an arbitrary topological diagonal is open.

Definition 2.13 packages these equivalent conditions as local total $0$-connectedness. The generator category is allowed to vary: once the condition is expressed for all connected slices of the topos, the proof of 2.6 applies to every full generating subcategory of connected objects.
\end{dcomment}


\textbf{Définition 2.13.} On dit qu'un topos $\X$ est \textit{localement totalement} $0$\textit{-connexe} s'il satisfait aux conditions équivalentes a, a$'$, b, d du théorème 2.6, i.e. si $\X$ est localement connexe et si de plus tout topos induit 0-connexe est totalement 0-connexe - le produit fibré de deux objets 0-connexes sur un troisième est 0-connexe. Il revient au même de dire que toute (ou que une) sous-catégorie pleine génératrice $S$ de $\X^\wedge$ formée d'objets 0-connexes est localement cofiltrante (en supposant toujours $\X$ localement connexe, ou ce qui revient au même, l'existence d'un $S$ formé d'objets 0-connexes).

On n'a pas terminé encore la démonstration de 2.6 - il reste à prouver que si $\X$ (supposé localement totalement 0-connexe) est 0-connexe, il est $\infty$-connexe - ou encore que les composantes connexes de $\X$, supposé localement totalement 0-connexe, sont $\infty$-connexes, i.e. $W_\infty$-asphériques.

\msp{page 30}\textbf{Corollaire 2.14.} \textit{Notons quand même que si $\X$ est de la forme Top$(T)$, $T$ un espace topologique, alors $\X$ est localement totalement} $0$\textit{-connexe si et seulement si il est localement irréductible, i.e. si et seulement si $T$ est localement irréductible.}

% ---- from Der-ChII-234_p021-025.tex
\printedpage{19}
% head: A. Grothendieck : Dérivateurs, Ch. II, Édition des Universités de Montpellier II et Paris VII
% foot: Version 8 juillet 2001

On est ramené à prouver que cette dernière condition implique la condition d) de 2.6, et on peut prendre pour $S$ l'ensemble ordonné des ouverts connexes non vides de $S$, cela devient alors évident.

Ceci dit, il faut se rappeler que la théorie cohomologique des espaces topologiques irréductibles n'est nullement triviale - comme nous le montre FAC de \textsc{Serre}\,! D'autre part, c'est \textsc{Serre} qui m'avait fait observer que pour des coefficients \textit{constants,} la cohomologie d'un espace irréductible est triviale.

\dgrblock{dgrb99}{Irreducible spaces still have nonconstant cohomology}
\begin{dcomment}{examples}{p30:la cohomologie d'un espace irréductible est triviale}{Irreducible spaces still have nonconstant cohomology}
The reference to Serre separates two coefficient questions which must remain distinct. On an irreducible nonempty topological space $T$, every nonempty open is connected. For an abelian group $M$, the constant sheaf $\underline M$ consequently has sections $M$ on every nonempty open, with identity restriction maps. Restrictions to the empty open are also surjective. Thus $\underline M$ is flabby and
\[
 H^q(T,\underline M)=0\quad(q>0),\qquad
 H^0(T,\underline M)=M.
\]
The same argument on each irreducible component applies when the space is a disjoint union of irreducible open components, which is the local condition used in 2.14.

For a concrete contrast, take the Zariski space $T=\mathbf P^1_k$, which is irreducible, and the line bundle $\mathcal O(-2)$. On the standard affine charts with coordinate $t$ and $t^{-1}$, its degree-one \v{C}ech quotient is
\[
 \frac{k[t,t^{-1}]}{k[t]+t^{-2}k[t^{-1}]}
   =k\,[t^{-1}].
\]
The first summand removes all nonnegative powers and the second all powers at most $-2$; exactly the coefficient of $t^{-1}$ remains. The affine charts and their intersection are acyclic for this coherent sheaf, so the quotient is $H^1(\mathbf P^1_k,\mathcal O(-2))\simeq k$. The affine-cover comparison used here is \citet[no.~47, Theorem 4]{Serre1955FAC}; the displayed computation gives the particular example in full.

For 2.14 itself, choose the connected nonempty opens as generators. If the space is locally irreducible in Grothendieck's sense, each such open is irreducible. Two nonempty opens inside it meet, giving PS 1; the generator category is a poset, so parallel arrows are already equal and PS 2 imposes nothing further. The converse follows from local total $0$-connectedness implying local irreducibility. This is why the two notions coincide for spaces while the group example separates them for general topoi.
\end{dcomment}


Dans le cas actuel, prouvons déjà, si $\mathcal{X}$ est 0-connexe et localement totalement 0-connexe, que tout revêtement de l'objet final de $\mathcal{X}$ est trivial. Il suffit de voir que si $U$ est 0-connexe, $U \longrightarrow e$ est un isomorphisme. Mais considérons $U \hookrightarrow U \times U$, on sait que c'est un sommande direct, et comme $U \times U$ est connexe, c'est un isomorphisme. Donc $U \longrightarrow e$ \texttt{[est un]} monomorphisme, donc sommande direct, et comme $U \neq \emptyset$, on a $U \longrightarrow e$ isomorphisme.

Reste à prouver que pour tout faisceau \msp{page 31} abélien constant $F$ sur $\mathcal{X}$, on a $H^i(X,F)=0$ pour $i>0$. On procède par récurrence sur $i$, c'est déjà fait pour $i=1$. Supposons que ce soit prouvé pour les $j<i$ ($i$ fixé) pour tout topos \textit{totalement} 0\textit{-connexe.} Soit $\xi \in H^i(\mathcal{X},F)$, soit $U \longrightarrow e$ \texttt{[un]} épimorphisme tel que $\xi|U=0$, et considérons la suite spectrale % [sic] H^i(X,F) -> H^i(\mathcal{X},F), plain X printed where the topos \mathcal{X} is meant elsewhere
\[
E_2^{pq} = H^p(C(U^*,\mathcal{H}^q(F))) \Longrightarrow H^*(\mathcal{X},F).
\]
On a par hypothèse totalement 0-connexe
\[
\pi_0(U^q) \simeq \pi_0(U)^q, \tag*{$(*)$}
\]
et les composantes connexes de $U^q$ sont totalement 0-asphériques,
\begin{center}
\begin{tikzpicture}[scale=0.8]
\draw[->] (0,0) -- (0,5);
\node[left] at (0,5) {$q$};
\draw[->] (0,0) -- (5,0);
\node at (4.3,-0.6) {$p$};
\node[left] at (0,4) {$i$};
\foreach \x in {0,1,2,3,4} {\fill (\x,4) circle (2pt); \fill (\x,0) circle (2pt);}
\end{tikzpicture}
\end{center}
donc par hypothèse de récurrence
\[
E_1^{pq} = 0 \qquad \text{si } 0<q<i,
\]
donc on trouve une suite exacte
\[
0 \longrightarrow H^i(C(U^*,F)) \longrightarrow H^i(\mathcal{X},F) \longrightarrow H^0(\ \underbrace{H^i(U^*,F)}_{\substack{\text{complexe}\\ \text{co-semisimplicial}\\ \text{pour } i\\ \text{variable}}}\ ) = \Ker(H^i(U,F) \rightrightarrows H^i(U^2,F)).
\]

\dgrblock{dgrb100}{The hereditary hypothesis in the cohomological proof}
\begin{dcomment}{exposition}{p31:Supposons que ce soit prouvé pour les}{The hereditary hypothesis in the cohomological proof}
Grothendieck's induction proves the intended assertion after keeping its hypothesis fixed: the topos must be connected and \emph{locally} totally $0$-connected. The word ``localement'' disappears in the sentence introducing the induction at manuscript p.~31, but it is needed when the argument passes to components of powers of a covering object. Those components define connected slices of the original topos, and hence remain in the same class by 2.6. Mere total $0$-connectedness does not provide this inheritance.

First consider a connected locally constant object $L$. Its diagonal in $L\times L$ is complemented, as one checks after a cover trivializing $L$. The product is connected, so this noninitial summand is the whole product. Hence $L\to1$ is monic. Its image is also complemented, by the same local triviality; connectedness of the base and nonemptiness of $L$ give $L\simeq1$. Every covering object is therefore a coproduct of copies of $1$. In the source paragraph, $U$ must still denote a connected component of a \emph{covering}, not an arbitrary connected sheaf.

For the higher cohomology, assume the vanishing in degrees $0<q<n$ for every connected topos in this hereditary class. Given $\xi\in H^n(\mathcal E,\underline M)$, choose an epimorphism $U\to1$ killing $\xi$. Such a cover exists because a positive-degree cohomology class represented in an injective resolution becomes a boundary locally. The augmented \v{C}ech nerve gives the descent spectral sequence
\[
 E_1^{pq}=H^q(\mathcal E/U^{p+1},\underline M)
 \ \Longrightarrow\ H^{p+q}(\mathcal E,\underline M).
\]
Each $U^{p+1}$ is a coproduct of connected objects whose slices remain in the class; thus the rows $0<q<n$ vanish. The edge sequence gives
\[
 0\longrightarrow H^n(C^\bullet)
 \longrightarrow H^n(\mathcal E,\underline M)
 \longrightarrow \ker\bigl(H^n(U,\underline M)
       \rightrightarrows H^n(U^2,\underline M)\bigr).
\]
The last image of $\xi$ is zero by choice of $U$. If $I=\pi _0(U)$, product preservation identifies
\[
 C^p=\operatorname{Map}(I^{p+1},M).
\]
Choose $i_*\in I$, which exists since $U\to1$ is epic. Inserting $i_*$ as first coordinate defines
$(sc)(i_0,\ldots,i_{p-1})=c(i_*,i_0,\ldots,i_{p-1})$.
For the alternating deletion differential, direct cancellation gives $ds+sd=1$ in positive degrees, with the usual augmentation in degree zero. Thus $H^n(C^\bullet)=0$ and $\xi=0$. The ``trivial'' complex in the manuscript is this contractible augmented complex; its individual groups and differentials are not zero.
\end{dcomment}


\printedpage{20}
% head: A. Grothendieck : Dérivateurs, Ch. II, Édition des Universités de Montpellier II et Paris VII
% foot: Version 8 juillet 2001

Par hypothèse sur $U$, l'image de $\xi$ dans le troisième terme est nulle, donc $\xi$ appartient au premier terme. Mais pour $(*)$, le complexe $C(U^*,F)$ est le complexe \textit{trivial} des puissances de $I=\pi_0(U)$, à coefficients dans $M$ (si $F=M_{\mathcal{X}}$). Sa cohomologie \msp{page 32} en degrés $>0$ est nulle, donc $\xi=0$, q.e.d.

\dgrblock{dgrb101}{From connected components to the hypercover homotopy type}
\begin{dcomment}{later}{p32:en degrés $>0$ est nulle}{From connected components to the hypercover homotopy type}
The stronger homotopical conclusion of 2.6 can also be proved directly, without deducing weak contractibility from an unspecified version of a cohomological Whitehead theorem. On a connected locally totally $0$-connected topos, the component functor preserves finite limits, as shown above, and preserves epimorphisms: epimorphisms in a topos are regular, and the left adjoint $\pi _0$ preserves their coequalizer presentations. Let $U_\bullet\to1$ be a hypercover by objects of this topos. Its matching objects are finite limits. Therefore
\[
 \pi _0\bigl(M_n(U_\bullet)\bigr)
       \simeq M_n\bigl(\pi _0(U_\bullet)\bigr),
\]
and the hypercover epimorphisms become surjections of sets. The simplicial set $\pi _0(U_\bullet)$ consequently has a filler for every boundary $\partial\Delta^n\to\pi _0(U_\bullet)$, including the nonemptiness condition in degree zero. Its map to a point is a trivial Kan fibration, and it is weakly contractible.

In the hypercover construction of the homotopy type, whose classical setting is \citet[\S\S8--9]{ArtinMazur1969}, every stage is therefore weakly contractible. This proves the $W_\infty$-asphericity meant here, as well as the corresponding assertion for each connected object by passage to its slice. For a presheaf topos it says concretely that every connected category of elements is weakly contractible.

There is also a short algebraic verification of the constant-coefficient calculation. Since $\pi _0$ preserves finite limits, it induces an exact functor on abelian group objects, left adjoint to the constant-abelian-sheaf functor. The latter therefore preserves injectives. Applying it to an injective resolution of $M$, and then taking global sections, recovers the same resolution of $M$, because the topos is connected. Thus $R\Gamma(\mathcal E,\underline M)\simeq M[0]$. These proofs make explicit why the hereditary condition repairs the false inference after 2.3. They impose no claim that the topos itself is equivalent to sets; the preceding projective-line example shows that arbitrary coefficients retain substantial information.
\end{dcomment}


\noindent\textbf{Proposition 2.15.} (\footnote[39]{C'est déjà dans SGA 4, I 8.}) \textit{Soit $X$ un objet de} Cat\textit{. Conditions équivalentes :}
\begin{enumerate}[label={}, leftmargin=2em, itemsep=0pt]
\item \textit{a) $X$ est filtrant, i.e. $X$ est $0$-connexe et $X^o$ satisfait les conditions équivalentes du théorème 2.6 ($X^{o,\wedge}$ est localement totalement $0$-connexe).}
\item \textit{b) L'objet final de $X^\wedge$ est ind-représentable.}
\end{enumerate}
\noindent\textsc{Démonstration}. On sait qu'on a
\[
e_{X^\wedge} \simeq \varinjlim_{x \in X} x, \tag*{(2.15.1)}
\]
limite inductive prise dans $X^\wedge$. En effet, cela signifie aussi que pour tout objet $a$ dans $X$, on a
\[
\varinjlim_{x \in X} \ \underbrace{x(a)}_{\Hom_X(a,x)} \ \xrightarrow{\ \sim\ } e \qquad \text{(ensemble ponctuel)}
\]
Mais considérons l'objet $X' = a\backslash X \longrightarrow X$ coétale sur $X$, qui est donc Cat-cofibré sur $X$, on sait que l'on a
\[
\pi_0(X') \simeq \varinjlim_{x \in X} \pi_0(X'_x),
\]
où $X'_x = \Hom(a,x)$, donc on trouve
\[
\varinjlim_{x \in X} \Hom(a,x) = \pi_0(a\backslash X) = e
\]
car $a\backslash X$ est 0-connexe.

\dgrblock{dgrb102}{The terminal presheaf as an ind-object}
\begin{dcomment}{exposition}{p32:C'est déjà dans SGA 4}{The terminal presheaf as an ind-object}
The passage to ind-objects replaces the cone conditions by a representability statement. The general form, recorded in 2.18, is the familiar criterion from \citet[Expos\'e I, Theorem 8.3.3]{SGA4I1972}: a presheaf $F$ on a small category $X$ is a filtered colimit of representables if and only if its category of elements $X/F$ is filtered. We give the argument, since it also supplies all the details of 2.15.

For every presheaf there is the density presentation
\[
 F\simeq\mathop{\rm colim}_{(x,s)\in X/F}h_x.
\]
Objectwise, the element represented by a pair $a\to x$, $s\in F(x)$ maps to its pullback in $F(a)$; it has the representative $(a,t)$ with identity map, and the relations identify every other representative with this one. If $X/F$ is filtered, this is the required ind-presentation.

Conversely, write $F=\operatorname{colim}_{i\in I}h_{x_i}$ with $I$ filtered. The category of elements is nonempty, since the image of $1_{x_i}$ supplies an object. Two elements $s\in F(a)$ and $t\in F(b)$ have representatives at stages $i,j$; passage to a common stage $k$ supplies arrows from $(a,s)$ and $(b,t)$ to $(x_k,\overline{1}_{x_k})$. Given parallel arrows $u,v:(a,s)\rightrightarrows(b,t)$, represent $t$ by $b\to x_i$. The two composites $a\rightrightarrows x_i$ have the same image in $F(a)$. Equality in a filtered colimit of sets is witnessed at a common later stage. The resulting map $(b,t)\to(x_j,\overline{1}_{x_j})$ equalizes $u,v$. Thus $X/F$ is filtered.

For $F=1$, its category of elements is $X$ itself, giving 2.15. The formula $1=\operatorname{colim}_{x\in X}h_x$ is valid even when $X$ is not filtered: at $a$ its value is $\pi _0(a\backslash X)=\{*\}$ because $(a,1_a)$ is initial. The condition lies in the possibility of a \emph{filtered} presentation. The reindexing by a directed poset used in the manuscript is Deligne's result, \citet[Expos\'e I, Proposition 8.1.6]{SGA4I1972}; the proof above does not need this extra reindexing.
\end{dcomment}


La formule (2.15.1) montre que si \msp{page 33} $X$ est filtrant, $e_{X^\wedge}$ est limite inductive filtrante de foncteurs représentables. On sait d'ailleurs qu'on peut réindexer par une catégorie \textit{ordonnée} filtrante (cf. SGA 4, I 8.1.6 (c'est dû à P. \textsc{Deligne}\,!)).

On a prouvé a $\Longrightarrow$ b, prouvons b $\Longrightarrow$ a. Supposons donc
\[
e_{X^\wedge} = \varinjlim_{i \in I} x_i,
\]
où $I$ est un (petit) ensemble ordonné filtrant, et $I \longrightarrow X$ un foncteur. Prouvons PS 1$'$ (forme droite), i.e. pour $a, b$ dans $X$, \texttt{[il]} existe $\vcenter{\hbox{\begin{tikzcd}[row sep=0.1em, column sep=1.2em]
a \arrow[dr] & \\
& x \\
b \arrow[ur] &
\end{tikzcd}}}$. En effet, comme $\varinjlim\limits_{I} \Hom(a,x_i) = e$, il s'ensuit que $\Hom(a,x_i) \neq \emptyset$ pour $i$ grand, de même $\Hom(b,x_i) \neq \emptyset$, donc il suffit de prendre $\vcenter{\hbox{\begin{tikzcd}[row sep=0.1em, column sep=1.2em]
a \arrow[dr] & \\
& x_i \\
b \arrow[ur] &
\end{tikzcd}}}$. Prouvons PS 2$_{\mathrm{d}}$, i.e. que pour \texttt{[une]} double flèche $a \overset{u}{\underset{v}{\rightrightarrows}} b$, \texttt{[il]} existe $b \longrightarrow x$ qui égalise.

\printedpage{21}
% head: A. Grothendieck : Dérivateurs, Ch. II, Édition des Universités de Montpellier II et Paris VII
% foot: Version 8 juillet 2001

\msp{page 34}
Soit $\xi_{i_0}$ dans $\Hom(b,x_{i_0})$, ses deux images dans $\Hom(a,x_{i_0})$ ne sont pas nécessairement égales, mais le deviennent dans la limite inductive, réduite à $e$, donc deviennent égales en composant avec $x_{i_0} \longrightarrow x_i$ pour $i$ assez grand, donc le composé $b \longrightarrow x_{i_0} \longrightarrow x_i$
\[
a \rightrightarrows b \longrightarrow x_{i_0} \longrightarrow x_i
\]
égalise la double flèche.

\noindent\textbf{Corollaire 2.16.} \textit{Soit $W$ un localiseur fondamental satisfaisant l'axiome des limites} Loc (8)\textit{. Si $X$ dans} Cat \textit{est une catégorie filtrante ou cofiltrante, alors $X$ est $W$-asphérique.}

Plus généralement et agréablement :

\dgrblock{dgrb103}{Filtered categories and the minimal localizer}
\begin{dcomment}{later}{p34:Plus généralement et agréablement}{Filtered categories and the minimal localizer}
Corollary 2.16 needs only the standard contractibility of a filtered category, together with the comparison between its weak equivalences and the chosen localizer. Here is the finite-data argument behind the first assertion. In a filtered category, any finite collection of objects has maps into one object. Given finitely many arrows between those objects, successive applications of the parallel-arrow axiom modify the common target so that all the required triangles commute. Thus every finite diagram of objects and arrows occurring in a finite simplicial subset of the nerve admits a compatible cone.

Every map from a finite sphere into the realization of the nerve can, after simplicial approximation, be represented in a finite such subset. The cone makes that map nullhomotopic. The nerve is nonempty and connected, so all its homotopy groups vanish. Reversing the category does not change the homotopy type of its nerve, and gives the cofiltered assertion as well.

The manuscript formulates the conclusion for a fundamental localizer satisfying its axiom Loc(8). A later theorem permits a sharper modern statement: the nerve weak equivalences $W_\infty$ form the smallest fundamental localizer. This is \citet[Theorem 2.2.11, p.~136]{Cisinski2004Minimal}. With the standard modern definition of fundamental localizer, every filtered or cofiltered category is consequently $W$-aspherical without an additional filtered-limit axiom. This is a later strengthening of the hypothesis in the manuscript, rather than a reason to alter its historical formulation. The finite-cone argument also underlies the higher categorical characterization cited earlier: filtered colimits of spaces commute with finite limits.
\end{dcomment}


\noindent\textbf{Corollaire 2.17.} \textit{Soit $W$ comme dans 2.16. Alors toute catégorie filtrante est} totalement $W$-asphérique (\footnote[40]{cf. XII, §V pour cette notion, que j'ai beaucoup peiné à dégager et à comprendre\,!}), \textit{i.e. pour toute flèche $f : X \longrightarrow Y$ dans} Cat\textit{, l'objet $f_!(e_X)$ de $e_Y$ (défini aussi comme $y \longmapsto \pi_0(y\backslash X)$, et intervenant dans la factorisation canonique de $f$ en $ip$, avec $p$ étale et les cofibres} \texttt{[?]} \textit{de $p$ $0$-connexes)} \msp{page 35} \textit{est $W$-asphérique. En fait, $Y/f_!(e_{X^\wedge})$ est même une catégorie filtrante.}

En effet, le foncteur $f_!$ commute aux limites inductives, donc $f_! e_{X^\wedge}$ est limite inductive filtrante des $f_!(x_i) = f(x_i)$. Donc $Y/f_!(e_{X^\wedge})$ est limite inductive filtrante des $Y/f(x_i)$. Donc c'est même une catégorie filtrante, a fortiori elle est $W$-asphérique.

Donc finalement, on trouve des exemples non triviaux de catégories totalement $W$-asphériques, i.e. des exemples donc des catégories sans objet final. % [sic] donc -> dont (?), the conjunction "donc" makes no sense here

\dgrblock{dgrb104}{Left Kan extension carries the ind-presentation}
\begin{dcomment}{exposition}{p35:Donc finalement, on trouve des exemples non triviaux}{Left Kan extension carries the ind-presentation}
The reason behind 2.17 is that an ordinary left Kan extension sends a representable to a representable and preserves colimits. If $f:X\to Y$ and $X$ is filtered, then
\[
 f_!1_X\simeq\mathop{\rm colim}_{x\in X}f_!h_x
          \simeq\mathop{\rm colim}_{x\in X}h_{f(x)}.
\]
The result is an ind-representable presheaf on $Y$, so $Y/(f_!1_X)$ is filtered by the criterion just proved. This gives the precise argument for the last assertion of 2.17 and therefore its asphericity conclusion. The proof does not require $f$ to be fully faithful or cofinal.

The canonical factorization mentioned in the statement has the following form:
\[
 X\xrightarrow{p}Y/(f_!1_X)\xrightarrow{i}Y.
\]
Here $i$ is the discrete fibration associated with the presheaf, and
\[
 (f_!1_X)(y)=\pi _0(y\backslash X).
\]
The first map sends $x$ to the component represented by $(x,1_{f(x)})$. For an object $(y,c)$ of the category of elements, the comma category $(y,c)\backslash X$ is exactly the component $c$ of $y\backslash X$; it is nonempty and connected. Thus $p$ is cofinal in the terminology of these pages. The edition assigns ``\'etale'' to $p$, but it is the second map $i$ which is \'etale; the connected cofibres belong to $p$. The phrase ``objet $\cdots$ de $e_Y$'' similarly means an object of $\widehat Y$.

For the simplest example without a terminal object, take $X=Y=\mathbf N$ with its usual order and $f=1_X$. Its terminal presheaf is the increasing colimit of the representables, and its category of elements is again $\mathbf N$. This is filtered and weakly contractible, although no natural number is terminal. It supplies exactly the kind of example for which Grothendieck introduces the ind-object formulation. Corollary 2.18 is the same argument with an arbitrary presheaf in place of the terminal one.
\end{dcomment}


\noindent\textbf{Corollaire 2.18.} (\footnote[41]{déjà dans SGA 4, I 8.}) \textit{Soit $X$ dans} Cat\textit{, $F$ dans $X^\wedge$. Pour que $F$ soit ind-représentable il faut et il suffit que $X/F$ soit filtrante.}

En effet, par (2.15), $X/F$ est filtrante si et seulement si l'objet final $F$ de $X/F$ provient d'un ind-objet de $X/F$, mais cela signifie aussi que $F$ est la limite inductive dans $X^\wedge$ d'un ind-objet de $X$, i.e. que $F$ est ind-représentable.

\msp{page 36}
\noindent\textbf{Question.} \textit{Si $X$ dans} Cat \textit{est totalement $W_\infty$-asphérique, i.e. les $f_!(e_{X^\wedge})$ sont $W_\infty$-asphériques pour toute flèche $f : X \longrightarrow Y$ dans} Cat\textit{, ou encore pour toute telle flèche, avec les $y\backslash Y$ $0$-connexes ($y \in \Ob Y$), $Y$ est $W_\infty$-asphérique - cela implique-t-il que $X$ soit filtrante ?} Ça ne m'étonnerait pas - et même en supposant seulement les $f_!(e_{X^\wedge})$ 1-connexes (\footnote[42]{Ils le sont toujours (cf. XII, p. 101), c'est 2-\textit{connexe} qu'il faut lire ici\,!}) (ils sont 0-connexes, dès que $X$ l'est). Mais je ne veux pas m'attarder à présent sur cette question.

\dgrblock{dgrb105}{The quantifier in total asphericity}
\begin{dcomment}{exposition}{p36:Mais je ne veux pas m'attarder à présent sur cette question.}{The quantifier in total asphericity}
Grothendieck is asking which indexing categories remain homotopically harmless after an ordinary change of index. The adjective ``total'' concerns every possible target of a functor; it does not merely strengthen the contractibility of the nerve of the original category. It is useful to isolate the precise condition before reading the proposed answer. Write $B A=|N A|$, and write $T(A)$ for the assertion that $B E$ is weakly contractible whenever $A\longrightarrow E$ is an ordinary cofinal functor. Here and below ``cofinal'' has the orientation fixed in Section~3: each category $e\backslash A$ must be nonempty and connected.

This formulation is equivalent to the one using $f_!1_A$ on the page. Indeed, for $f:A\longrightarrow Y$, put $P=f_!1_A$, so that
\[
 P(y)=\pi_0(y\backslash A).
\]
Let $Y/P$ be its category of elements. There is a canonical factorization
\[
 A\xrightarrow{q}Y/P\longrightarrow Y,
 \qquad a\longmapsto(fa,[a,1_{fa}]).
\]
For an object $(y,c)$ of $Y/P$, the comma category $(y,c)\backslash A$ is precisely the connected component $c$ of $y\backslash A$. Thus $q$ is cofinal. Conversely, if $f$ itself is cofinal, then $P=1_Y$ and $Y/P=Y$. Consequently $T(A)$ says exactly that the category of elements of every ordinary presheaf $f_!1_A$ is weakly contractible. The occurrence of $y\backslash Y$ in the printed question is a slip: it is $y\backslash A$ that must be connected.

A filtered category satisfies $T$. Its ordinary cofinal images are filtered by Corollary~3.4 below, and their nerves are weakly contractible. The difficult direction is whether this test forces the equalization of parallel arrows. The counterexample inserted on pages~36bis--36ter withdraws the attempted negative answer; it does not prove the converse. The converse holds for monoids and for categories with finitely many objects and fails in general; both are proved in the commentary ``Total asphericity is not filteredness'' below.

There is a terminological precaution for the later literature. In \citet[Definition~1.7.2 and Proposition~1.7.1]{Maltsiniotis2020}, a category is called totally aspherical when it is aspherical and its diagonal is aspherical, equivalently when products of representable presheaves are aspherical. This is a different quantifier from $T(A)$. In particular, the fact that the simplex category is a strict test category does not settle the question on this page. We retain $T(A)$ only inside this commentary, to keep the question itself unchanged.
\end{dcomment}


La réponse à la question précédente est négative, comme on voit par la :

\printedpage{22}
% head: A. Grothendieck : Dérivateurs, Ch. II, Édition des Universités de Montpellier II et Paris VII
% foot: Version 8 juillet 2001

\noindent\textbf{Proposition 2.19.} (\footnote[43]{\textbf{Faux.}}) \textit{Soit $X$ dans} Cat \textit{telle que $\mathit{diag}_X : X \longrightarrow X \times X$ soit cofinal} (\footnote[44]{cf. définition 3.3, p. 45.})\textit{, i.e. telle que $X^o$ soit totalement $0$-connexe. Alors $X$ est totalement $W_\infty$-asphérique (donc aussi totalement $W$-asphérique pour tout localiseur fondamental $W$ tel que $W \supseteq W_\infty$).}

Soit en effet un foncteur cofinal $f : X \longrightarrow Y$, il faut prouver que $Y$ est $W_\infty$-asphérique. Or on sait que totalement 0-connexe implique $W_\infty$-asphérique (\footnote[45]{C'est \textbf{faux.}}) (cf. proposition 2.3, p. 17-18), donc il suffit de voir que $Y^o$ est totalement 0-connexe, i.e. $Y \xrightarrow{d_Y} Y \times Y$ cofinal. Or on a le diagramme commutatif
\[
\begin{tikzcd}[row sep=large, column sep=large]
X \arrow[r, "d_X"] \arrow[d, "f"'] & X \times X \arrow[d, "f \times f"] \\
Y \arrow[r, "d_Y"] & Y \times Y,
\end{tikzcd}
\]
où par hypothèse $f$, donc aussi $f \times f$, et $d_X$ sont cofinaux. Donc $\varphi = (f \times f) \circ p_X$ \texttt{[plutôt} $\varphi = (f \times f) \circ d_X$\texttt{]} est cofinal, et comme $\varphi = d_Y f$ et $f$ est aussi cofinal, on conclut que $d_Y$ l'est, q.e.d.

\noindent\textbf{Corollaire.} $\boldsymbol{\Delta}$ \textit{étant totalement $0$-connexe, $\boldsymbol{\Delta}^o$ est totalement $W_\infty$-asphérique, pourtant elle n'est pas filtrante.}

\noindent\textbf{Question.} \textit{$X$ totalement $W_\infty$-asphérique $\Longrightarrow$ $X^o$ totalement $0$-connexe ? Ou du moins implique PS 1$'$ ?}

\msp{page 36bis}
\noindent\textbf{Contrexemple 2.20} à la proposition fausse 2.19.

Soit $X$ un objet de Cat, rajoutons-lui la catégorie finale $e$ comme fermé, i.e. considérons un objet $X'$ de Cat dans lequel $X$ soit un sous-objet ouvert, $e$ le fermé complémentaire :
\[
X \overset{i}{\hookrightarrow} X' \overset{j}{\hookleftarrow} e.
\]
On sait (cf. XII) qu'une telle opération est déterminée, à isomorphisme près, par la donnée d'un objet
\[
H \in \Ob X^\wedge
\]
en posant
\[
H(x) = \Hom_{X'}(x,e),
\]
la catégorie $X/e$ s'identifiant (à isomorphisme près) à $X/H$. Elle est donc 0-connexe si et seulement si l'objet $H$ de $X^\wedge$ est 0-connexe. Cela signifie aussi que le morphisme $i$ est cocofinal, i.e. que $i^o : X^o \longrightarrow X'^o$ est cofinal. Supposons maintenant $X$ $W_\infty$-asphérique et totalement 0-connexe, i.e. l'actuel Loc (6), cf. XVI. La proposition 2.19 affirme que, sous ces conditions, $X'$ (qui est lui aussi totalement 0-connexe) est même $W$-asphérique. Il s'ensuivrait donc que $i \in W_\infty$. Mais par XII, th. 4, cor. 4 (p. 149), il s'ensuivrait que
\printedpage{23}
% head: A. Grothendieck : Dérivateurs, Ch. II, Édition des Universités de Montpellier II et Paris VII
% foot: Version 8 juillet 2001
\msp{page 36ter}
$X/e$ est $W_\infty$-asphérique, i.e. $H$ $W_\infty$-asphérique. Ainsi, si $X$ est totalement 0-asphérique, tout objet 0-connexe de $X^\wedge$ serait $\infty$-connexe. Mais c'est manifestement faux, p.ex. si on prend pour $X$ une catégorie-test pour $W_\infty$, p.ex. $X = \boldsymbol{\Delta}$.

\dgrblock{dgrb106}{The retracted proposition and a visible sphere}
\begin{dcomment}{examples}{p36:p.ex. si on prend pour $X$ une catégorie-test}{The retracted proposition and a visible sphere}
The useful part of Proposition~2.19 is its elementary calculation with cofinal functors: the ordinary cofinal image of a nonempty sifted category is sifted. The unsupported step is the passage from that ordinary connectedness condition to weak contractibility. The edition marks this step and the proposition as false. The appeal in Counterexample~2.20 to Chapter~XII also needs its hypothesis restored: the reflection of a weak equivalence to the attaching category uses the strong form of $W(7\mathrm{bis})$, explicitly required in \citet[Chapter~XII, Section~V.3.4, Corollary~4, note~104]{Grothendieck1990XII}. That reflection is not valid for arbitrary weak homotopy equivalences. We can make the intended counterexample entirely explicit without using it.

The omitted reflection condition has a concrete obstruction. Choose a finite connected acyclic noncontractible simplicial complex $K$; a finite presentation complex with this homotopy type is given in \citet[Example~2.38 and Theorem~2C.5]{Hatcher2002}. Its suspension is simply connected and acyclic, hence contractible by the homology Whitehead theorem. The face posets of the two cones on $K$ are open subcategories covering the face poset of $\Sigma K$, with intersection the face poset of $K$. The two cones and their union are weakly contractible, but their intersection is not. Thus the weak equivalence from one cone into the union does not reflect to a weak equivalence from the intersection into the other cone. This contradicts the strong biconditional $W(7\mathrm{bis})$, while leaving the direct homotopy-pushout theorem intact.

Let $H=\Delta[1]/\partial\Delta[1]$, the simplicial circle, and form the category $E$ with full subcategory $\Delta$, one additional object $z$, and
\[
 E([n],z)=H_n,\qquad E(z,[n])=\varnothing,
 \qquad E(z,z)=\{1_z\}.
\]
Composition into $z$ is the simplicial action of $H$. Thus $\Delta/z=\Delta/H$, the category of simplices of $H$. Its classifying space has the homotopy type of $|H|=S^1$; this is the usual category-of-simplices realization, or the test-category equivalence for $\Delta$ \citep[Proposition~1.7.14]{Maltsiniotis2020}. Since $\Delta$ has terminal object $[0]$, its classifying space is contractible.

There is a direct description of $B E$. An $n$-simplex of $N E$ that does not lie in $N\Delta$ consists of a chain in $\Delta$ followed by an arrow into $z$, with possible repetitions of $z$. The arrow into $z$ and all its composites are exactly a simplex of $N(\Delta/H)$. These simplices glue to the mapping cone of the forgetful map $B(\Delta/H)\longrightarrow B\Delta$. Hence
\[
 B E\simeq\operatorname{hocolim}\bigl(B\Delta\longleftarrow
 B(\Delta/H)\longrightarrow *\bigr)\simeq\Sigma S^1\simeq S^2.
\]
The cone description can also be read degree by degree: adjoining the last vertex $z$ supplies the cone direction, and deleting that vertex applies the forgetful map. It therefore identifies the attaching maps, not merely the vertices of the two spaces.

The inclusion $i:\Delta\hookrightarrow E$ is ordinary initial: $\Delta/e$ has a terminal object when $e\in\Delta$, while $\Delta/z=\Delta/H$ is nonempty and connected. Thus $i^{\mathrm{op}}:\Delta^{\mathrm{op}}\longrightarrow E^{\mathrm{op}}$ is cofinal. The source $\Delta^{\mathrm{op}}$ is sifted: for $[m],[n]$, its category of common targets is opposite to $\Delta/(\Delta[m]\times\Delta[n])$, which is connected because the product of the two simplices is contractible. The target $E^{\mathrm{op}}$ is consequently sifted as well, whereas $B(E^{\mathrm{op}})\simeq S^2$. This supplies both a counterexample to the implication used in Proposition~2.19 and a direct counterexample to its asserted corollary about $\Delta^{\mathrm{op}}$.

The distinction is the same one encountered earlier in the chapter between products of connected objects and their higher homotopy. Here its failure has a specific witness, a two-dimensional sphere. It is not a failure of the composition or cancellation properties of ordinary cofinal functors.
\end{dcomment}


Cela me fait soupçonner que la conjecture
\[
X \text{ totalement } W_\infty\text{-asphérique} \Longrightarrow X \text{ filtrante}
\]
pourrait être valable.

\dgrblock{dgrb107}{A necessary directedness condition}
\begin{dcomment}{development}{p36:pourrait être valable.}{A necessary directedness condition}
Grothendieck asks, before returning to cofinality, whether his condition at least forces the existence of common successors. This weaker question admits an elementary answer. The argument below also settles his full question for preorders. We did not find this particular implication isolated in the consulted treatments of sifted categories, asphericity, or lax epimorphisms; the following is our continuation of the argument, with no assertion of priority.

\textbf{Proposition (the common-successor obstruction).}\ours
\textit{If $T(A)$ holds, then $A$ is nonempty and any two objects of $A$ admit arrows to a common object. Consequently, for a preorder $A$, the condition $T(A)$ is equivalent to filteredness.}

\textit{Proof.}
The identity of $A$ is cofinal, so $T(A)$ implies that $B A$ is weakly contractible. In particular $A$ is nonempty and connected. Suppose that objects $a,b$ have no common successor. Let $U$ be the full subcategory of objects receiving an arrow from $a$, and let $V$ be the full subcategory of objects receiving an arrow from $b$. These subcategories are disjoint and closed under passage to successors.

Let $P$ be the three-object category with precisely the nonidentity arrows
\[
 0\xrightarrow{p}1,\qquad 0\xrightarrow{q}2.
\]
Define $f:A\longrightarrow P$ by sending $U$ to $1$, $V$ to $2$, and all other objects to $0$. Closure under successors makes this a functor: an arrow cannot leave $U$ or $V$, and there are no arrows between them. The category $0\backslash A$ computed for $f$ is $A$ itself. The other two comma categories are $U$ and $V$. They are connected, since every object of $U$ receives an arrow from $a$, and every object of $V$ receives one from $b$. They need not have initial objects; the existence of those arrows is all that is used. Hence $f$ is cofinal.

We next give $P$ a cofinal image with nonzero second homology. Adjoin an object $z$ and arrows
\[
 r,s:z\rightrightarrows0,
 \qquad v:z\longrightarrow1,
 \qquad w:z\longrightarrow2,
\]
subject to $pr=ps=v$ and $qr=qs=w$, and no other nonidentity arrows. Call the resulting category $Q$. The inclusion $j:P\hookrightarrow Q$ is cofinal. For an object of $P$, the relevant comma category has an initial object. For $z$, the category $z\backslash P$ has four objects $r,s,v,w$ and four arrows
\[
 r\longrightarrow v,\quad r\longrightarrow w,
 \quad s\longrightarrow v,\quad s\longrightarrow w;
\]
it is nonempty and connected.

To see exactly what obstructs contractibility, use normalized simplicial chains of $NQ$. There are four nondegenerate two-simplices, denoted $(r,p),(s,p),(r,q),(s,q)$, and no nondegenerate simplices in higher dimensions. Their boundaries are
\[
\begin{aligned}
 \partial(r,p)&=p-v+r,&\partial(s,p)&=p-v+s,\\
 \partial(r,q)&=q-w+r,&\partial(s,q)&=q-w+s.
\end{aligned}
\]
Therefore
\[
 c=(r,p)-(s,p)-(r,q)+(s,q)
\]
is a nonzero cycle. Since there are no three-simplices, it is not a boundary. More precisely, the coefficients of $p,q,r,s,v,w$ in the boundary show that every two-cycle is an integral multiple of $c$, so $H_2(BQ;\mathbf Z)\cong\mathbf Z$. The composite $jf$ is cofinal, contradicting $T(A)$. This proves the common-successor assertion.

In a preorder, parallel arrows are already equal. Nonemptiness and common successors therefore imply filteredness. The converse follows from preservation of filteredness by ordinary cofinal images, as in Corollary~3.4. This proves the last assertion.\hfill$\square$

For general categories the argument leaves precisely the parallel-arrow equalization problem. The distinction is substantive: a one-object group category satisfies the common-successor condition for every group, while filteredness forces that group to be trivial. The next commentary settles it: the implication $T(A)\Rightarrow A\text{ filtered}$ holds when some object receives an arrow from every object, and fails for the category of finite ordinals and increasing injections.
\end{dcomment}


\dgrblock{dgrb108}{Total asphericity is not filteredness}
\begin{dcomment}{development}{p36:pourrait être valable.}{Total asphericity is not filteredness}
The question whether $T(A)$ forces $A$ to be filtered has a negative answer in general and a positive one for monoids and for categories with finitely many objects. Both rest on a description of $T(A)$ in which no target category appears. For a functor $H:A\to\mathrm{Set}$ write $\operatorname{El}(H)$ for its category of elements: the objects are the pairs $(a,h)$ with $h\in H(a)$, and an arrow $(a,h)\to(a',h')$ is an arrow $v:a\to a'$ with $H(v)h=h'$. The statements below were not found in the sources consulted, except where a reference is given; no priority is claimed.

\textbf{Theorem (a criterion on set-valued functors).}\ours
\textit{For a small category $A$ the following are equivalent: (i) $T(A)$; (ii) $A$ is nonempty and connected, and for every functor $H:A\to\mathrm{Set}$ every connected component of $\operatorname{El}(H)$ has the integral homology of a point; (iii) the same with ``is weakly contractible''. If they hold, every connected component $E$ of every $\operatorname{El}(H)$ satisfies $T(E)$.}

\textit{Proof.}
Let $K:A\to\mathrm{Set}$ have nonempty connected $\operatorname{El}(K)$, and let $A_K$ be $A$ with one object $x$ adjoined, $A_K(x,a)=K(a)$, no arrows into $x$ but its identity, and composition given by the action of $A$ on $K$. For the inclusion $A\to A_K$ the category $x\backslash A$ is $\operatorname{El}(K)$, and $a\backslash A$ has an initial object for $a\in A$; so the inclusion is cofinal. A chain of arrows of $A_K$ not contained in $A$ consists of $x$ (possibly repeated) followed by a chain of $\operatorname{El}(K)$. Hence $N(A_K)$ is $NA$ with the cone on $N\operatorname{El}(K)$ attached along the forgetful map, and
\[
 \widetilde H_{j+1}(BA_K)\longrightarrow\widetilde H_j(B\operatorname{El}(K))\longrightarrow\widetilde H_j(BA)
\]
is exact. Under (i) both $BA$ and $BA_K$ are weakly contractible, so $\operatorname{El}(K)$ has the homology of a point. A component of $\operatorname{El}(H)$ for an arbitrary $H$ is $\operatorname{El}(K)$ for the subfunctor $K$ of the elements lying in it. This proves (i)$\Rightarrow$(ii).

(ii)$\Rightarrow$(iii). Let $E$ be a component of $\operatorname{El}(H)$ and $L:E\to\mathrm{Set}$ a functor. Then $\operatorname{El}(L)\cong\operatorname{El}(H')$ for the functor $H'(a)=\coprod_hL(a,h)$ on $A$, the sum over the $h\in H(a)$ with $(a,h)\in E$. So $E$ again has property (ii). Let $\Pi E$ be the groupoid obtained by inverting all arrows of $E$, fix an object $e_0$, and identify $G=\operatorname{Aut}_{\Pi E}(e_0)$ with $\pi_1(BE,e_0)$. For a subgroup $\Gamma\subseteq G$ the functor $L_\Gamma(e)=\operatorname{Hom}_{\Pi E}(e_0,e)/\Gamma$ has bijective transition maps, and $\operatorname{El}(L_\Gamma)\to E$ is the connected covering with fundamental group $\Gamma$. By (ii) for $E$ its first homology $\Gamma^{\mathrm{ab}}$ vanishes. Taking $\Gamma$ cyclic gives $G=1$. A simply connected CW complex with the homology of a point is weakly contractible.

(iii)$\Rightarrow$(i). Let $f:A\to B$ be cofinal. For $b\in B$ the category $b\backslash A$ is $\operatorname{El}(a\mapsto B(b,f(a)))$; it is nonempty and connected, hence weakly contractible by (iii), and Quillen's Theorem~A makes $Bf$ a weak equivalence. The constant functor with value a point shows that $BA$ is weakly contractible. So is $BB$. The last assertion was shown in the second step.\hfill$\square$

In particular the integral-homology form of $T(A)$ (every cofinal image has the homology of a point) is equivalent to $T(A)$. With coefficients in a nonzero commutative ring $R$ the same proof, with the homology form of Theorem~A, gives: every cofinal image of $A$ has the $R$-homology of a point if and only if $A$ is nonempty and connected and the components of all $\operatorname{El}(H)$ do. For $R=\mathbf Q$ this is weaker: a finite group satisfies it.

\textbf{Proposition (exact colimits suffice).}\ours
\textit{If $A$ is nonempty and connected and $\varinjlim_A$ is exact on functors $A\to\mathrm{Ab}$, then $T(A)$.}

\textit{Proof.}
For $H:A\to\mathrm{Set}$ and $F:\operatorname{El}(H)\to\mathrm{Ab}$ put $(\Sigma_HF)(a)=\bigoplus_{h\in H(a)}F(a,h)$. The functor $\Sigma_H$ is exact and $\varinjlim_{\operatorname{El}(H)}F=\varinjlim_A\Sigma_HF$. So colimits over $\operatorname{El}(H)$ are exact, and so are colimits over a component $E$ (extend by zero). The left derived functors of $\varinjlim_E$ on the constant functor $\mathbf Z$ are the groups $H_j(BE)$; they vanish for $j>0$. Criterion (ii) applies.\hfill$\square$

\textbf{Theorem (a category with $T$ that is not filtered).}\ours
\textit{Let $\Delta'$ be the category of the ordered sets $[n]=\{0,\dots,n\}$, $n\geq0$, and strictly increasing maps. Then $T(\Delta')$ holds and $\Delta'$ is not filtered; it is not sifted either.}

\textit{Proof.}
The two maps $[0]\to[1]$ are not equalized by any injection, so $\Delta'$ is not filtered. A pair of maps $[0]\to[N]\leftarrow[0]$ is two points of $[N]$, and whether the first is smaller than, equal to or larger than the second is preserved by increasing injections; so the category of such pairs has three components and the diagonal of $\Delta'$ is not cofinal.

For exactness we show: \emph{for every $n$ there are $M$ and an integral combination $t=\sum_gt_g\,g$ of maps $g:[n]\to[M]$ of $\Delta'$, with $\sum_gt_g=1$, such that $t\circ h$ is the same combination for all maps $h:[m]\to[n]$ with the same $m$.} For $n=1$, $M=2$, one can take $t=(0,1)-(0,2)+(1,2)$, naming a map by its values. In general, a map $g:[n]\to[M]$ is determined by its gaps $e_0=g(0)$, $e_i=g(i)-g(i-1)-1$, $e_{n+1}=M-g(n)$, that is, by a monomial $z^e$ of degree $M-n$ in $z_0,\dots,z_{n+1}$; a combination $t$ becomes a homogeneous polynomial $P_t$, and $\sum t_g=P_t(1,\dots,1)$. Let $D_i$ be the substitution $z_{i+1}\mapsto z_i$ followed by renumbering to variables $y_0,\dots,y_n$. Omitting the value $g(i)$ merges two gaps and the omitted point, so $P_{t\delta^i}=y_iD_i(P_t)$ for the map $\delta^i:[n-1]\to[n]$ that omits $i$. We construct, for all $n\geq0$ and $k\geq1$, a homogeneous polynomial $C_{n,k}$ with integer coefficients, $C_{n,k}(1,\dots,1)=1$, such that $y_i^kD_i(C_{n,k})$ does not depend on $i$. Put $C_{0,k}=1$. Given $Q=C_{n-1,2k}$ in the variables $y_0,\dots,y_n$, we want $D_i(C_{n,k})=F_i:=Q\prod_{l\neq i}y_l^k$ for all $i$; then $y_i^kD_i(C_{n,k})=Q\prod_ly_l^k$. In the coordinates $x_i=z_i-z_{i+1}$ ($0\leq i\leq n$), $u=z_{n+1}$, the map $D_i$ is reduction modulo $x_i$. The prescriptions agree modulo $(x_i,x_j)$: for $j=i+1$ both give $D_i(Q)\prod_lw_l^k$, and for $j\geq i+2$ they give $w_{j-1}^{2k}D_{j-1}(Q)\prod_{l\neq i,j-1}w_l^k$ and $w_i^{2k}D_i(Q)\prod_{l\neq i,j-1}w_l^k$, which are equal by the property of $Q$. For a nonempty set $S$ of indices let $\Phi_S$ be the common value of the $F_i$, $i\in S$, after setting $x_i=0$ for all $i\in S$ (in homogeneous lifts written in the coordinates $x,u$). Then
\[
 C_{n,k}=\sum_{S\neq\varnothing}(-1)^{|S|+1}\Phi_S
\]
reduces to $F_j$ modulo $x_j$, because the terms $S$ and $S\cup\{j\}$, $j\notin S$, cancel modulo $x_j$ and the term $S=\{j\}$ remains. It is homogeneous, and its value at $(1,\dots,1)$, a point where all $x_i$ vanish, is $F_0(1,\dots,1)=1$. Now let $t$ correspond to $C_{n,1}$. Then $s=t\delta^i$ is independent of $i$, and from $\delta^{b+1}\delta^a=\delta^a\delta^b$ ($a\leq b$) we get $s\delta^a=s\delta^b$; by induction on $n$, $s\circ h'$ depends only on the source of $h'$, and every non-identity $h:[m]\to[n]$ factors through some $\delta^i$. This proves the claim.

Let $F\subseteq F'$ be functors $\Delta'\to\mathrm{Ab}$. Any two objects of $\Delta'$ map to a third, so every element of $\varinjlim F$ is the class of one element $\xi\in F([n])$. Suppose its image in $\varinjlim F'$ is zero. The colimit is the quotient of $\bigoplus_mF'([m])$ by the elements $F'(v)\eta-\eta$, so $\xi=\sum_k(F'(v_k)\eta_k-\eta_k)$ in this direct sum, for finitely many $v_k:[s_k]\to[r_k]$ and $\eta_k\in F'([s_k])$. Choose $N$ and maps $d_m:[m]\to[N]$ for the finitely many $m$ that occur; applying $F'(d_m)$ on each summand gives $F'(d_n)\xi=\sum_k\bigl(F'(d_{r_k}v_k)-F'(d_{s_k})\bigr)\eta_k$ in $F'([N])$. Choose $t$ for $N$ and write $F'(t)=\sum_gt_gF'(g)$. Since $t\circ d_{r_k}v_k=t\circ d_{s_k}$, applying $F'(t)$ gives $F'(\tau)\xi=0$ for $\tau=t\circ d_n$, a combination of maps $[n]\to[M]$ with coefficient sum $1$. As $F\subseteq F'$, also $F(\tau)\xi=0$. The class of $F(g)\xi$ in $\varinjlim F$ is the class of $\xi$ for every $g$, so the class of $F(\tau)\xi$ is that of $\xi$; it is zero. So $\varinjlim_{\Delta'}$ preserves monomorphisms and is exact; $\Delta'$ is nonempty and connected, and the proposition applies.\hfill$\square$

That colimits of abelian groups over this category are exact is attributed to Isbell and Mitchell in \citet[A.5]{AdamekKoubekReiterman1981}, together with the general criterion: $\varinjlim_A$ is exact on abelian groups if and only if each component of the category with the same objects and the integral combinations of arrows of coefficient sum $1$ as morphisms is filtered. The withdrawn corollary on this page claimed $T$ for $\Delta^{\mathrm{op}}$, which is false; the category $\Delta'$ is an example of what that corollary was after.

\textbf{Theorem (monoids; finitely many objects).}\ours
\textit{Let $A$ have an object $q$ that receives an arrow from every object. Then $T(A)$ holds if and only if $A$ is filtered. This applies to every monoid, and to every category with finitely many objects, where both $T(A)$ and filteredness provide such an object $q$.}

\textit{Proof.}
Assume $T(A)$. \emph{Step 1: for parallel arrows $a,b:p\to q'$ there are an object $r$ and an integral combination $t$ of arrows $q'\to r$, with coefficient sum $1$, such that $ta=tb$.} Let $P_c=\mathbf Z[A(c,-)]$, let $d:P_{q'}\to P_p$ be $c\mapsto ca-cb$, with kernel $K$ and image $J$, and let $H$ be the quotient of $A(p,-)$ by the relations $ca\sim cb$, so that $0\to J\to P_p\to\mathbf Z[H]\to0$ is exact. The functors $P_c$ are projective with $\varinjlim P_c=\mathbf Z$, and $\varinjlim d=0$. So $L_1\varinjlim\mathbf Z[H]$ is the kernel of $\varinjlim J\to\varinjlim P_p$; this map is zero, because $\varinjlim P_{q'}\to\varinjlim J$ is surjective and $\varinjlim d=0$. Hence $L_1\varinjlim\mathbf Z[H]=\varinjlim J=\operatorname{coker}(\varinjlim K\to\mathbf Z)$, the image of $\varinjlim K$ being generated by the coefficient sums of the elements of the groups $K(r)$. The left side is $H_1(B\operatorname{El}(H))$, and $\operatorname{El}(H)$ is connected (every object receives an arrow from the class of $\mathrm{id}_p$); by the criterion it vanishes. So $1$ is a sum of coefficient sums of finitely many elements of groups $K(r_i)$; pushing them to a common successor $r$ of the $r_i$ (the common-successor proposition of the preceding commentary) and adding gives $t$.

\emph{Step 2: the monoid $M=A(q,q)$ is filtered.} For $x,m\in M$, Step~1 for the pair $x,xm$, followed by an arrow $r\to q$, gives $t\in\mathbf Z[M]$ with coefficient sum $1$ and $tx=txm$. The element $\nu=tx$ of $\mathbf Z[Mx]$ is invariant under right multiplication $\rho$ by $m$. Let $S$ be its finite support. Each point of $S$ has a $\rho$-preimage in $S$, distinct points have distinct preimages, so $\rho$ permutes $S$ and $\nu$ is constant on its cycles; the coefficient sum $1$ is thus an integral combination of cycle lengths. Among all points of $Mx$ that are $\rho$-periodic choose one, $w$, of least period $\pi$. If $\pi>1$, let $\ell$ be a prime dividing $\pi$ and apply the same to $w$ in place of $x$: some cycle in $Mw$ has length prime to $\ell$, and its points $w'$ satisfy $w'm^\pi=w'$ too, so their period is a proper divisor of $\pi$, a contradiction. Hence there is $y$ with $yxm=yx$. Given $m,m'\in M$, take $y$ with $ym=y$ (case $x=1$), then $y'$ with $y'ym'=y'y$; then $y'ym=y'y=y'ym'$.

\emph{Step 3.} Let $a,b:p\to q'$ be parallel, choose $g:q'\to q$, and by Step~1 an element $t=\sum n_ic_i\in\mathbf Z[M]$ of coefficient sum $1$ with $tga=tgb$. In the filtered monoid $M$ there is $y$ with $yc_i=y$ for all $i$. Then $yga=\sum n_iyc_iga=\sum n_iyc_igb=ygb$. With common successors this shows that $A$ is filtered. The converse is the fact that filtered categories satisfy $T$.\hfill$\square$

So Grothendieck's expectation is right in every finite case and for all monoids, and the smallest counterexamples have infinitely many objects. For the monoid $\{1,e,f\}$ with $uv=v$ for $u,v\in\{e,f\}$, which is cofiltered with contractible classifying space, the theorem says that $T$ fails; a cofinal image of the homotopy type of $S^2$ is obtained by adjoining a source $x$ with two arrows $h,i:x\to*$ on which $e$ and $f$ act by the constant map with value $i$. Whether $T(A)$ implies exactness of $\varinjlim_A$ on abelian groups, the converse of the proposition, is not decided here.
\end{dcomment}


% ---- from Der-ChII-234_p026-030.tex
\printedpage{24}
% head: A. Grothendieck : Dérivateurs, Ch. II, Édition des Universités de Montpellier II et Paris VII
% foot: Version 8 juillet 2001
\msp{page 37}

\section*{3\quad Retour sur la cofinalité}

Il est instructif de relier les questions de cofinalité (pour une flèche $u : I' \longrightarrow I$ dans Cat) au dérivateur $\HOT_{W_0}$, où $W_0 \subseteq \Fl\,\Cat$ est le plus grossier des localiseurs dans Cat satisfaisant l'axiome de connexité Loc 5, i.e.
\[
W_0 = \{ f \in \Fl\,\Cat \mid \pi_0(f) \text{ bijectif}\}. \tag*{(3.1)}
\]
On trouve alors, pour $I$ dans Cat, une équivalence canonique, 2-fonctorielle en $I$,
\[
\HOT_{W_0}(I) \xrightarrow{\;\sim\;} I^{\wedge} \defeq \uHom(I,(\Ens)), \quad \text{\pl}\ \uHom(I^o,\Ens)\,\text{\texttt{]}} \tag*{(3.2)}
\]
en d'autres termes, on a une équivalence
\[
\HOT_{W_0} \simeq \mathbf{D}_{(\Ens)}. \tag*{(3.3)}
\]
Quand $\HOT_{W_0}$ est défini par localisation à partir de la catégorie $\Cat(I) \defeq \uHom(I^o,\Cat)$, la flèche (3.2) est celle déduite, par localisation, de la flèche
\[
\begin{array}{rcl}
\pi_0^I : \uHom(I^o,\Cat) & \longrightarrow & \uHom(I^o,(\Ens)) \\[1ex]
X & \longmapsto & \pi_0 \circ X,
\end{array} \tag*{(3.4)}
\]
où
\[
\pi_0 : \Cat \longrightarrow (\Ens) \tag*{(3.5)}
\]
est le foncteur ‘composantes connexes’. Si on décrit $\HOT_{W_0}(I)$ comme localisée de $\Cat/I$, (3.2) est déduit de la flèche
\[
\left\{
\begin{array}{rcl}
\Cat/I & \longrightarrow & I^{\wedge} = \uHom(I^o,(\Ens)) \\[1ex]
X & \longmapsto & \underline{\pi_0}(X/I) = (i \longmapsto \pi_0(i\backslash X)).
\end{array}
\right. \tag*{(3.6)}
\]

\msp{page 38}
Ainsi, les foncteurs familiers, pour une flèche
\[
u : I' \longrightarrow I
\]
dans Cat
\[
\left\{
\; I'^{\wedge} \quad
\begin{array}{c}
\xrightarrow{\;u_!\;} \\[1.5ex]
\xleftarrow{\;u^*\;} \\[1.5ex]
\xrightarrow{\;u_*\;}
\end{array}
\quad I^{\wedge},
\right. \tag*{(3.7)}
\]
\printedpage{25}
% head: A. Grothendieck : Dérivateurs, Ch. II, Édition des Universités de Montpellier II et Paris VII
% foot: Version 8 juillet 2001
s'interprètent comme les foncteurs de même nom pour le dérivateur $\HOT_{W_0}$, ou plus directement, pour le dérivateur (qu'on pourrait faussement considérer comme canularesque) $\mathbf{D}_{(\Ens)}$.

\dgrblock{dgrb109}{The represented derivator at dimension zero}
\begin{dcomment}{exposition}{p37}{The represented derivator at dimension zero}
Grothendieck now uses sets to separate the genuinely categorical part of cofinality from its higher homotopical refinements. The equivalence $\HOT_{W_0}(I)\simeq\operatorname{Fun}(I^{\mathrm{op}},\mathrm{Set})$ has a short proof that explains why no hidden higher information survives. Let $d:\mathrm{Set}\longrightarrow\mathrm{Cat}$ assign the discrete category, and let $\pi_0$ be connected components. There is an adjunction $\pi_0\dashv d$, with $\pi_0d=1$. For each category $C$, the unit $C\longrightarrow d\pi_0C$ is a $W_0$-equivalence. Applied pointwise to an $I^{\mathrm{op}}$-diagram, these units give a natural objectwise $W_0$-equivalence from that diagram to its discrete component diagram. Thus the two functors become inverse after localization. Since the construction commutes with restriction, it gives the stated equivalence of diagram theories.

The description by $\mathrm{Cat}/I$ requires the relative, rather than the global, component construction. To $X\longrightarrow I$ one associates the presheaf $i\longmapsto\pi_0(i\backslash X)$. Localizing the slice merely at functors that induce a bijection on the components of their total categories would give a different category. The relevant class is detected in all these comma categories, exactly as in the relative localization recalled in the commentary on Chapter~I, manuscript p.~43 (``What the localization axioms do and do not say'').

The large table then has an economical interpretation. The words ``aspherical'' and ``coaspherical'' become ordinary initiality and finality; the proper and smooth conditions say that specified strict lift categories are nonempty and connected. Right and left Kan extensions retain their ordinary meanings. Nothing in this reduction licenses replacing connected lift categories by contractible ones. Conversely, it makes the degree-zero assertions accessible by elementary calculations with connected components, without invoking a construction of general homotopy Kan extensions.
\end{dcomment}


Cela donne un sens à la cohorte des notions cohomologiques associées à un dérivateur, ou plus particulièrement, à un localiseur fondamental $W$ sur Cat (donnant naissance au dérivateur $\HOT_W$). On les désignera par le nom de $W_0$-fourbi (au lieu de $\HOT_{W_0}$-fourbi, ou $\mathbf{D}_{(\Ens)}$-fourbi), notamment pour une flèche $u : I' \longrightarrow I$ de Cat telle la notion
\[
\left\{
\begin{array}{lll}
\text{$W_0$-équivalence} & \Bigl(\iff & \text{$\pi_0(u)$ bijectif.}\Bigr) \\
\text{$W_0$-asphérique} & \Bigl(\iff & \text{les $I'/i$ pour $i \in \Ob I$ sont} \\
 & & \text{$0$-connexes.}\Bigr) \\
\text{$W_0$-coasphérique} & \Bigl(\iff & \text{les $i\backslash I'$ pour $i \in \Ob I$ sont} \\
 & & \text{$0$-connexes.}\Bigr) \\
\text{$W_0$-propre (\footnotemark[46])} & \Bigl(\iff & \text{les inclusions $I'_i \longrightarrow I'/i$ ($i \in \Ob I$)} \\
 & & \text{sont $W_0$-coasphériques, i.e. les fibres} \\
 & & \text{des colocalisées $i'\backslash I' \longrightarrow i\backslash I$} \\
 & & \text{($i' = u(i)$) \pl\ ($i = u(i')$)]} \\
 & & \text{sont $0$-connexes.}\Bigr) \\
\text{\mspin{page 39}} & & \\
\text{$W_0$-lisse (\footnotemark[47])} & \Bigl(\iff & \text{les inclusions $I'_i \longrightarrow i\backslash I'$} \\
 & & \text{sont $W_0$-asphériques, i.e. les fibres} \\
 & & \text{des flèches localisées $I'/i' \longrightarrow I/i$} \\
 & & \text{($i = u(i')$) sont $0$-connexes.}\Bigr) \\
\text{$W_0$-parfait} & \Bigl(\iff & \text{$W_0$-propre et $W_0$-lisse.}\Bigr) \\
\text{$W_0$-équivalence universelle (\footnotemark[48])} & \Bigl(\iff & \text{$\pi_0$ bijectif et le reste par tout} \\
 & & \text{changement de base. Ce sont aussi} \\
 & & \text{les $W_0$-fibrations (ou $W_0$-fibrations} \\
 & & \text{faibles à fibres $W_0$-asphériques,} \\
 & & \text{i.e. $0$-connexes.}\Bigr) \\
\text{$W_0$-fibration} & \Bigl(\iff & \text{$u^* : \Cat/I \longrightarrow \Cat/I'$ transforme} \\
 & & \text{$W_0$-équivalences en $W_0$-équivalences.}\Bigr) \\
\text{$W_0$-fibration faible (\footnotemark[49])} & \Bigl(\iff & \text{le changement de base} \\
 & & \text{$u^* : \Cat/I \longrightarrow \Cat/I'$ transforme} \\
 & & \text{homotopismes $J_1 \longrightarrow J_2$ en} \\
 & & \text{$W_0$-équivalences - et il suffit de} \\
 & & \text{l'exiger pour l'inclusion d'un objet} \\
 & & \text{final (ou au choix initial) et pour} \\
 & & \text{l'inclusion d'un point dans $\boldsymbol{\Delta}^1$.}\Bigr)
\end{array}
\right. \tag*{(3.8)}
\]
\footnotetext[46]{Exemple : les Cat-cofibrations.}
\footnotetext[47]{Exemple : les Cat-fibrations.}
\footnotetext[48]{Exemple : les flèches $W_0$-lisses ou $W_0$-propres, à fibres 0-connexes.}
\footnotetext[49]{Exemple : les flèches $W_0$-lisses telles que $\pi_0(I'/I) = (i \longmapsto \pi_0(i\backslash I')) = u_!(e_{I'^{\wedge}})$ soit localement constant sur $I$, i.e. $\underline{\pi_0}(I'/I)$ soit un \textit{revêtement} étale de $I$, ou si $u^{\mathrm{opp}}$ satisfait cette condition.}

\printedpage{26}
% head: A. Grothendieck : Dérivateurs, Ch. II, Édition des Universités de Montpellier II et Paris VII
% foot: Version 8 juillet 2001
\texttt{[Nous avons une]} implication \textit{stricte} (et comment)
\[
\text{$W_0$-équivalence universelle} \Longrightarrow \text{$W_0$-fibration}
\]
\texttt{[et une]} implication \textit{stricte} (et comment\,!) (\footnote[50]{Mais si les fibres sont 0-connexes, les deux notions de $W_0$-fibrations sont équivalentes et signifient qu'on a une $W_0$-équivalence universelle \dots})
\[
\text{$W_0$-fibration} \Longrightarrow \text{$W_0$-fibration faible.}
\]

Signalons, comme exemple d'utilisation de ces notions, que pour un diagramme cartésien dans Cat
\[
\begin{tikzcd}
I' \arrow[d, "u"'] & J' \arrow[l, "p'"'] \arrow[d, "v"] \\
I & J, \arrow[l, "p"']
\end{tikzcd} \tag*{(3.9)}
\]
la flèche de changement de base pour $u_*$, $v_*$ pour les préfaisceaux est \texttt{[un]} isomorphisme si $u$ est $W_0$-propre, ou $p$ $W_0$-lisse, et la flèche de changement de base pour $u_!$, $v_!$ est \texttt{[un]} isomorphisme si $u$ est $W_0$-lisse ou $p$ $W_0$ propre (\footnote[51]{Sauf erreur, ca donne \textit{deux caractérisations} des morphismes $W_0$-propres, ou des morphismes $W_0$-lisses.}). Je signale à ce propos que ces énoncés, qui a priori se formulent pour les préfaisceaux à valeurs dans (Ens), restent valables % [sic] ca -> ça (in footnote 51), missing cedilla
\msp{page 40}
pour des préfaisceaux à valeurs dans une catégorie $\mathcal{M}$ quelconque (sous réserve d'existence des operations d'images directes cohomologiques ou homologiques envisagées). Cela se voit p.ex. en notant que % [sic] operations -> opérations, missing accent
\[
\text{$W_0$-équivalence} \;\Longrightarrow\; \text{$\mathbf{D}_{\mathcal{M}}$-équivalence} \tag*{(3.10)}
\]
pour toute $\mathcal{U}$-catégorie $\mathcal{M}$, d'où résulte que toute (\footnote[52]{Pas tout à fait exacte pour \textit{toute} ‘$W_0$-notion’, il faut excepter les deux notions de $W_0$-fibration forte, et faible.}) ‘$W_0$-notion’ implique la $\mathbf{D}_{\mathcal{M}}$-notion correspondante, et permet donc d'appliquer à $\mathbf{D}_{\mathcal{M}}$ les résultats le concernant. (Il y a bien sûr un argument direct, moins savant, consistant à se ramener au cas où $\mathcal{M}$ est petite, et à utiliser alors le plongement pleinement fidèle $\mathcal{M} \hookrightarrow \mathcal{M}^{\wedge}$ resp. $\mathcal{M} \hookrightarrow ((\mathcal{M}^o)^{\wedge})^o$.)

Je rappelle pour mémoire une (troisième\,?) caractérisation des flèches $u : I' \longrightarrow I$ $W_0$-propres (resp. $W_0$-lisses) : ce sont celles pour lesquelles le foncteur changement de base $u^* : \Cat/I \longrightarrow \Cat/I'$ transforme flèches $W_0$-coasphériques (resp. $W_0$-asphériques) en itou. Cela gagne en intérêt, du fait que la notion de foncteur $W_0$-coasphérique (resp. $W_0$-asphérique) s'interprète comme celle de foncteur cofinal (resp. co-cofinal), cf. plus bas.

\printedpage{27}
% head: A. Grothendieck : Dérivateurs, Ch. II, Édition des Universités de Montpellier II et Paris VII
% foot: Version 8 juillet 2001
\msp{page 41}
La notion de $W_0$-fibration apparaît comme extrêmement forte - beaucoup plus forte à certains égards que celle (disons) de $W_\infty$-fibration ou de $W_\omega$-fibration. Elle équivaut à ceci (comme on voit aussitôt) : pour toute composante connexe $I'_\alpha$ de $I'$, désignant par $I_\alpha$ la composante connexe correspondante de $I$, et pour tout $J_\alpha$ 0-connexe et au dessus de $I_\alpha$, le produit fibré $I'_\alpha \times_{I_\alpha} J_\alpha$ est encore 0-connexe. (Ça ressemble aux histoires de totale 0-connexité d'un topos\,!) Mais cela signifie aussi que $I'_\alpha \longrightarrow I_\alpha$ est une $W_0$-équivalence universelle. Cela implique que ses fibres sont 0-connexes, et la réciproque est vraie si $u : I' \longrightarrow I$ est $W_0$-propre ou $W_0$-lisse. Ainsi, la condition essentielle pour une $W_0$-fibration est celle que les fibres des flèches induites $I'_\alpha \longrightarrow I_{u(\alpha)}$ \pl\ $I'_\alpha \longrightarrow I_\alpha$\texttt{]} soient 0-connexes. En fait, on a ceci

\dgrblock{dgrb110}{The lift categories behind ordinary base change}
\begin{dcomment}{exposition}{p40}{The lift categories behind ordinary base change}
The base-change assertion is an application of the lift-category criterion already developed in the commentary on Chapter~I, manuscript pp.~68--70 (``The category of lifts behind properness'' and ``Completing the tests in the proper base-change theorem''). Its present specialization is useful because all the relevant tests are tests of connectedness. For $u:A\longrightarrow B$, properness means that, for every $a\in A$ and arrow $\beta:u(a)\longrightarrow b$, the category
\[
 L(a,\beta)=\{a\xrightarrow{h}a'\mid u(a')=b, u(h)=\beta\}
\]
of strict lifts, with arrows over $1_b$, is nonempty and connected. Smoothness is the dual condition on arrows into $u(a)$. An opfibration supplies initial objects in the displayed lift categories; a fibration supplies the dual terminal objects. Connectedness is the exact weakening that suffices for ordinary Kan extensions.

For a strict pullback square as in (3.9), the comparison of the comma categories computing a right Kan extension is final if $u$ is proper or if the horizontal base-change functor is smooth. The ordinary cofinality theorem therefore gives the right base-change isomorphism. Reversing arrows gives the left formula. The same argument works in any coefficient category in which the requisite comma limits exist, so that these Kan extensions are computed pointwise. One can apply $\operatorname{Hom}(m,-)$ to the limits, or dually $\operatorname{Hom}(-,m)$ to the colimits, and reduce to the set-valued cofinality theorem. This avoids imposing unnecessary completeness assumptions on the coefficient category.

The passage from a $W_0$-equivalence to an equivalence measured by constant coefficients has a similarly elementary explanation. A limit of a constant diagram is a product indexed by the connected components of its shape, and a colimit is the corresponding coproduct, whenever those objects exist. A bijection on components preserves both. The printed footnote correctly excludes the two fibration notions from the blanket transfer of terminology: preserving a larger class of equivalences under every base change is a stronger requirement, even when the equivalences themselves have a simpler definition. The next proposition measures this effect exactly.
\end{dcomment}


\noindent\textbf{Proposition 3.1.} \textit{Soit $u : I' \longrightarrow I$ une flèche dans} Cat.
\begin{itemize}
\item[\textit{a)}] \textit{Pour que ce soit une $W_0$-équivalence universelle, il faut et il suffit que pour toute $\boldsymbol{\Delta}^1 \longrightarrow I$, i.e. toute flèche $\alpha$ de $I$, la catégorie $I' \times_I \boldsymbol{\Delta}^1 \defeq I'|\alpha$ soit $0$-connexe, ou encore que les fibres de $I$ soient $0$-connexes, et que pour toute flèche} non identique \textit{$\alpha$ de $I$, $I'|\alpha$ soit $0$-connexe.}
\item[\textit{b)}] \textit{Pour que $u : I' \longrightarrow I$ soit une $W_0$-fibration, il faut et il suffit que pour toute composante connexe $I'_{\alpha'}$ de}
\msp{page 42}
\textit{$I'$, posant $\alpha = u(\alpha')$, la flèche induite $I'_{\alpha'} \longrightarrow I_\alpha$ soit une $W_0$-équivalence universelle (cf. critère a) pour cette notion).}
\item[\textit{c)}] \textit{Pour que ce soit une $W_0$-fibration faible, il} faut \textit{que $\forall\, i \in I$, l'inclusion $I'_i \hookrightarrow i\backslash I'$ soit une $W_0$-équivalence, i.e. $\pi_0(I'_i) \xrightarrow{\;\sim\;} \pi_0(i\backslash I')$ soit bijectif,} et \textit{que $i \longmapsto \pi_0(i\backslash I')$ soit un préfaisceau localement constant sur $I$ (i.e. $\underline{\pi_0}(I'/I)$ soit un} revêtement \textit{étale de $I$). Cette condition est aussi suffisante, si $u$ est $W_0$-lisse. De plus (pour mémoire) $u$ est une $W_0$-fibration faible si et seulement si $u^{\mathrm{opp}}$ l'est.}
\end{itemize}

\dgrblock{dgrb111}{Why arrows detect universal connectedness}
\begin{dcomment}{exposition}{p42:Cette condition est aussi suffisante}{Why arrows detect universal connectedness}
Proposition~3.1 says that universal preservation of connected components can be checked on vertices and single arrows. Its proof is a useful model for reading the much stronger componentwise condition in part~(b). We give the details because the proposition is stated, rather than proved, here.

Suppose first that $u:A\longrightarrow B$ is a universal $W_0$-equivalence. Pullback along any $[1]\longrightarrow B$ gives a connected nonempty category, since $[1]$ is connected. Conversely, assume this arrow test. Applying it to an identity arrow at $b$ shows that $A_b\times[1]$, and hence $A_b$, is connected and nonempty. Let $v:J\longrightarrow B$ be arbitrary. Every object of $J$ has a nonempty fibre in $A\times_BJ$, so the map on components is surjective. If two objects of the pullback project to objects joined by a zigzag in $J$, each edge of that zigzag has a connected pullback by hypothesis. Its endpoints can be joined through that pullback, and the connected vertex fibres join any two choices of endpoints. Concatenating these finite zigzags proves injectivity. Thus
\[
 \pi_0(A\times_BJ)\xrightarrow{\sim}\pi_0(J)
\]
for every $J$, which is the required universal equivalence. Notice that the test concerns strict pullbacks, not comma categories.

Now let $A_\lambda$ be a connected component of $A$, mapping to the component $B_\mu$. If pullback by $u$ preserves every $W_0$-equivalence, apply it to $J\amalg\coprod_{\nu\ne\mu}B_\nu\longrightarrow B$, where $J\longrightarrow B_\mu$ is any connected nonempty category over $B_\mu$. The resulting map on components forces $A_\lambda\times_{B_\mu}J$ to be connected and nonempty. Decomposing an arbitrary category over $B_\mu$ into its components then shows that $A_\lambda\longrightarrow B_\mu$ is a universal $W_0$-equivalence. Conversely, if all these component maps are universal equivalences, then for any $J\longrightarrow B$ there is a natural bijection
\[
 \pi_0(A\times_BJ)\cong
 \pi_0(A)\times_{\pi_0(B)}\pi_0(J).
\]
This formula immediately shows that pullback by $u$ preserves every bijection on connected components. It proves part~(b), including components of $B$ not met by $A$.

The necessary conditions in part~(c) come from the same tests with categorical homotopy equivalences. The inclusion of the initial object $1_b$ into $b\backslash B$ is such an equivalence; its pullback identifies $A_b\longrightarrow b\backslash A$ on components. Testing both endpoints of an arrow in $B$ makes the resulting component presheaf locally constant. When $u$ is smooth, ordinary left base change identifies the component presheaf after every base change with the restriction of this local system. Its colimit depends only on the fundamental groupoid of the indexing category, and categorical homotopy equivalences induce equivalences of those groupoids. This proves the stated sufficiency. The opposite-category assertion follows directly because categorical homotopies and $\pi_0$ are unchanged by reversal of arrows.
\end{dcomment}


La plus grande différence entre la notion de $W_0$-fibration, et celle de $W_0$-fibration faible, me paraît que les morphismes $W_0$-parfaits (i.e. $W_0$-lisses et $W_0$-propres), et a fortiori les $W$-parfaits pour un localiseur fondamental $W \subseteq W_0$, sont forcément des $W$-fibrations faibles, mais (\textit{même si $W = W_\omega$,} le plus fin de tous les localiseurs fondamentaux) ne sont pas des $W$-fibrations. L'exemple le plus frappant à cet égard est sans doute celui des revêtements étales, qui sont $W_\omega$-parfaits (donc
\msp{page 43}
des $W_\omega$-fibrations faibles et a fortiori des $W_0$-fibrations faibles), mais pas des $W_0$-fibrations\,; car dans le cas où $I'$, $I$ sont 0-connexes, il faudrait que les fibres soient 0-connexes (cf. b) ci-dessus), i.e. que $u$ soit un isomorphisme\,!

\dgrblock{dgrb112}{Monodromy explains the strictness}
\begin{dcomment}{examples}{p43:i.e. que $u$ soit un isomorphisme}{Monodromy explains the strictness}
The covering example isolates what the componentwise test excludes. This is the same monodromy issue as in the commentary on Chapter~I, manuscript p.~98 (``Sets already detect the missing monodromy''), now applied to the two precise definitions in (3.8). Let $G$ be a nontrivial group. Its translation groupoid $E G$ has object set $G$ and exactly one arrow between each ordered pair of objects. The functor
\[
 E G\longrightarrow B G
\]
sends $g\longrightarrow h$ to $hg^{-1}$. It is a discrete fibration and a discrete opfibration, so the strict lift categories are singletons and the map is proper and smooth. Both total and base categories are connected, but the strict fibre is the discrete set $G$.

The inclusion of the sole object $*\longrightarrow BG$ is a $W_0$-equivalence. Its pullback along $EG\longrightarrow BG$ is the map $G_{\mathrm{disc}}\longrightarrow EG$, whose map on components is $G\longrightarrow *$. This proves directly that the covering is not a $W_0$-fibration. It is a weak $W_0$-fibration by the smooth criterion of Proposition~3.1(c): the relative component presheaf is the locally constant $G$-set with regular action. Thus weak fibrations permit monodromy between different fibre components, whereas the strong $W_0$ condition treats each connected component of the total category separately.

There is a subscript caution in the preceding printed paragraph. The conclusion established by this example is ``not necessarily a $W_0$-fibration.'' It does not assert that coverings fail to preserve weak homotopy equivalences under pullback. In fact their nerves are simplicial coverings, and they do preserve weak homotopy equivalences. Similarly a $W_0$-perfect map gives a weak $W_0$-fibration; a $W$-perfect map gives the corresponding weak $W$-fibration. These are distinct assertions when $W\subsetneq W_0$.
\end{dcomment}


Une propriété sympathique des morphismes $W_0$-lisses (resp. $W_0$-propres), c'est que ce sont des applications \textit{ouvertes} (resp. fermées) pour les espaces topologiques sous-jacents, ce qui signifie en effet que les fibres locales, i.e. des flèches localisées en haut $I'/i' \longrightarrow I/i$ (resp. colocales, i.e. les flèches colocalisées en bas $i'\backslash I' \longrightarrow i\backslash I$), sont surjectives, i.e. à fibres non-vides (ce qui est impliqué par la condition que celles-ci soient 0-connexes).

\printedpage{28}
% head: A. Grothendieck : Dérivateurs, Ch. II, Édition des Universités de Montpellier II et Paris VII
% foot: Version 8 juillet 2001
Dans tout ça je n'ai encore guère parlé des flèches $W_0$-asphériques resp. co-asphériques, qui pourtant sont ici une motivation principale pour introduire $W_0$\,!

\noindent\textbf{Proposition 3.2.} \textit{Conditions équivalentes sur la flèche $u : I' \longrightarrow I$ dans} Cat :
\begin{itemize}
\item[a)] \textit{Pour toute $\mathcal{U}$-catégorie $\mathcal{M}$ et tout foncteur $F : I \longrightarrow \mathcal{M}$, la flèche canonique (sous réserve que les petites limites inductives soient représentables dans $\mathcal{M}$),}
\[
\varinjlim_{I'} u^*(F) \longrightarrow \varinjlim_{I} F \tag*{(3.2.1)}
\]
\msp{page 44}
\textit{est} \texttt{[un]} \textit{isomorphisme.}
\item[a$'$)] \textit{Pour toute $\mathcal{U}$-catégorie $\mathcal{M}$ où les petites limites projectives sont représentables et tout $F \in \mathcal{M}(I) = \uHom(I^o,\mathcal{M})$, i.e. tout préfaisceau sur $I$ à valeurs dans $\mathcal{M}$, la flèche canonique}
\[
\varprojlim_{I} F \longrightarrow \varprojlim_{I'} u^*(F) \tag*{(3.2.2)}
\]
\textit{est} \texttt{[un]} \textit{isomorphisme.}
\item[a$''$)] (\footnote[53]{C'est une condition sur les morphismes de topos, savoir que si $F$ \texttt{[est]} dans $\mathcal{X}^{\wedge}$, on a
\[
H^0(\mathcal{X},F) \xrightarrow{\;\sim\;} H^0(\mathcal{X}',u^*F).
\]
Il y a \texttt{[une]} notion similaire pour les $H^i$ ($i \leq n$), à dégager. Notion de morphismes $W_0$-coasphériques de topos\,! C'est dans la nature d'une condition de 0-connexité relative, duale de la notion $W_0$-asphérique, qui signifie simplement que pour $U$ 0-connexe dans $U^{\wedge}$ \pl\ $\mathcal{X}^{\wedge}$\texttt{]}, $f^*(U)$ \pl\ $u^*(U)$\texttt{]} est 0-connexe. De même, on a sans doute des notions (‘duales’ des précédentes) de $W_n$-coasphéricité (bien comprise en principe si $\mathcal{X}$ est localement $n$-asphérique) (La $W_n$-coasphéricité, étant celle évoquée plus haut.)}) \textit{Comme a$'$, avec $\mathcal{M} = (\Ens)$, i.e. $F \in \Ob I^{\wedge}$.}~(\footnote[54]{\textbf{NB} On peut introduire
\begin{itemize}
\item[a$'''$)] Comme a), mais avec $\mathcal{M} = (\Ens)$.
\item[a$^{IV}$)] Comme a), avec le foncteur canonique $I' \longrightarrow I'^{\wedge}$.
\end{itemize}}).
\item[b)] \textit{La flèche $u$ est $W_0$-coasphérique, i.e. pour tout $i \in I$, $i\backslash I'$ est $0$-connexe.}
\end{itemize}
L' équivalence de a, a$'$, a$''$ est $\pm$ tautologique en termes des définitions (cf. SGA 4, I, 8\dots). La forme a$''$ s'écrit, posant $\mathbf{D} = \HOT_{W_0}$, que $\forall\, F \in \mathbf{D}(I)$, on veut que
\[
\underbrace{H^{\bullet}_{\mathbf{D}}(I,F)}_{\defeq p_{I,*}(F)} \longrightarrow \underbrace{H^{\bullet}_{\mathbf{D}}(I',u^*F)}_{\defeq p_{I',*}(F)} \qquad \text{\pl}\quad \underbrace{H^{\bullet}_{\mathbf{D}}(I',u^*F)}_{\defeq p_{I',*}(u^*F)}\ \text{\texttt{]}} \tag*{(3.2.3)}
\]
($p_I : I \longrightarrow e$ et $p_{I'} : I' \longrightarrow e$ les foncteurs canoniques) soit \texttt{[un]} isomorphisme. On sait que cela équivaut à $u$ $\mathbf{D}$-coasphérique. La démonstration est un jeu amusant entre $u_!$, $u^*$, $u_*$ à la fois. Il faut exprimer que le Hom d'un objet $M$ de $\mathbf{D}(e)$ ($=$ (Ens) en l'occurence) dans la flèche (3.2.3) est bijective. Ici, il % [sic] occurence -> occurrence, French spelling

% ---- from Der-ChII-234_p031-035.tex
\printedpage{29}
% head: A. Grothendieck : Dérivateurs, Ch. II, Édition des Universités de Montpellier II et Paris VII
% foot: Version 8 juillet 2001

\msp{page 45}
est plus naturel de prendre $M = e$ (ensemble ponctuel) et d'interpréter la flèche (3.2.3) comme
\[
\Hom(e_I, F) \longrightarrow \Hom(e_{I'}, u^*F)
\tag*{(3.2.4)}
\]
et par adjonction de $u_!$, $u^*$ le deuxième membre s'interprète comme $\Hom(u_!(e_{I'}), F)$, et la flèche est déduite de
\[
u_!e_{I'} = u_!u^*(e_I) \xrightarrow{\text{adj.}} e_I
\tag*{(3.2.5)}
\]
en prenant le Hom dans $F$. Donc dire que (3.2.4) est \texttt{[un]} isomorphisme pour tout $F$ signifie que (3.2.5) est \texttt{[un]} isomorphisme. Passant aux valeurs des deux membres en $i$, on trouve $\pi_0(i\bs I') \xrightarrow{\sim} 1$, i.e. que $i\bs I'$ est 0-connexe, q.e.d.

\dgrblock{dgrb113}{What the ordinary cofinality theorem tests}
\begin{dcomment}{exposition}{p45:i.e. que $i\bs I'$ est 0-connexe, q.e.d.}{What the ordinary cofinality theorem tests}
Grothendieck has now proved the criterion announced on the opening page, with the sufficient directedness condition replaced by the exact connectedness condition. The proof is the representability calculation already explained in ``The presheaf represented by a change of index'': the functor $F\longmapsto\Gamma(I',u^*F)$ is represented by $u_!1_{I'}$, and
\[
 (u_!1_{I'})(i)=\pi_0(i\backslash I').
\]
Consequently the comparison on all set-valued presheaves is an isomorphism precisely when $u_!1_{I'}\longrightarrow1_I$ is an isomorphism. Applying ordinary Hom functors reduces arbitrary existing limits and colimits to this case. This is the classical theorem of SGA~4, Expos\'e~I, Section~8.1, in the variance used throughout these pages \citep{SGA4I1972}.

Two small examples distinguish its quantifiers. The inclusion $\{1\}\hookrightarrow[1]$ is cofinal, because each $i\backslash\{1\}$ is a point. For a presheaf $F(1)\longrightarrow F(0)$, its limit is therefore $F(1)$. But the restriction functor forgets $F(0)$ and is not fully faithful. On the other hand, for any nonempty connected category $C$, the projection $C\longrightarrow *$ is cofinal even when $BC$ has nontrivial homotopy groups. Ordinary cofinality is therefore stronger than a bijection on components of arbitrary source and target, weaker than full faithfulness of restriction, and weaker than homotopy cofinality in a different direction. These comparisons refer to different properties of a functor, not to alternative conventions for the same word.

Footnote~53 distinguishes two degree-zero conditions on a morphism of topoi $f:\mathcal E'\longrightarrow\mathcal E$. Its coasphericity condition preserves global sections of every object. Its asphericity condition instead requires $f^*U$ to be noninitial and connected whenever $U$ is. For locally connected topoi, the latter is equivalent to preservation of component sets: decompose $U$ into its noninitial connected components and use preservation of coproducts by $f^*$. For the morphism induced by $u:I'\longrightarrow I$, the restricted representable $u^*h_i$ has category of elements $I'/i$. Thus this asphericity condition is ordinary initiality of $u$, by the cofinality theorem with arrows reversed; the section condition is finality. The two endpoints of $[1]$ distinguish them: inclusion of $0$ is initial, while inclusion of $1$ is final.

The supplementary Yoneda test in footnote~54 should use the functor $I\longrightarrow I^{\wedge}$. Its colimit is $1_I$, whereas its restriction along $u$ has colimit $u_!1_{I'}$. The printed $I'\longrightarrow I'^{\wedge}$ has the wrong domain for the test as stated. The further suggestion in footnote~53 about higher cohomology should be read with coefficients specified: preservation of all set-valued sections, preservation of abelian cohomology through a fixed degree, and equivalence on all homotopy-valued diagrams detect different truncations. The earlier cohomological cofinality block supplies the derived comparison; the ordinary theorem here establishes its degree-zero model.
\end{dcomment}


\noindent\textbf{Définition 3.3.}
\begin{itemize}[label={}]
\item[a)] Quand les conditions équivalentes de 3.2 sont satisfaites, donc quand les $i\bs I'$ ($i \in \Ob I$) sont 0-connexes, on dit que $u : I' \longrightarrow I$ est un foncteur \textit{cofinal.} On dit que $u$ est \textit{co-cofinal} si $u^o : I'^o \longrightarrow I^o$ est cofinal, i.e. si les localisées $I'/i$ de $I'$ sur $I$ ($i \in \Ob I$) sont 0-connexes.
\item[b)] Si $I'$ est une \textit{sous-catégorie} de $I$, on dit que $I'$ \textit{est cofinale dans} $I$ si le foncteur d'inclusion est cofinal, \textit{co-cofinale} si ce foncteur est co-cofinal.
\end{itemize}

\noindent\textbf{Corollaire 3.4.} (\footnote[55]{\textbf{NB} Pour un énoncé complémentaire de cofinalité, cf. XVIII prop. 4.9.6.3 (p. 224-225, 228-242). Il faudrait un énoncé qui coiffe les deux.})
\begin{itemize}[label={}]
\item[\textit{a)}] \textit{Pour que $u$ soit cofinal, il} faut \textit{que les catégories colocalisées $i\bs I'$ soient non-vides, i.e. qu'on ait}

\msp{page 46}
\begin{itemize}[label={}]
\item[\textit{F 0)}] $\pi_0(f)$ \texttt{[plutôt $\pi_0(u)$]} \textit{est bijectif.}
\item[\textit{F 1)}] \textit{Pour tout} $i \in \Ob I$, \texttt{[il]} \textit{existe} \texttt{[un]} $i' \in I'$ \textit{et une flèche} $i \longrightarrow u(i')$ \textit{dans} $I$.
\end{itemize}
(\footnote[56]{\textbf{NB} F 1 implique déjà \texttt{[la]} surjectivité de $\pi_0$, donc sa bijectivité si $I'$ est 0-connexe - dans ce cas F 0 résulte de F 1.}).
\item[\textit{b)}] \textit{Si $I'$ est faiblement filtrante, pour que $u$ soit cofinale, il} suffit \textit{qu'elle satisfasse F 1, ainsi que la condition :}
\begin{itemize}[label={}]
\item[\textit{F 2)}] \textit{Pour toute double flèche} $i \dbl{\beta}{\gamma} u(i')$ \textit{(avec} $i \in \Ob I$, $i' \in \Ob I'$\textit{),} \texttt{[il]} \textit{existe une flèche} $i' \xrightarrow{\alpha'} j'$ \textit{dans} $I'$\textit{, de telle façon que} $u(\alpha')$ \textit{égalise} $(\beta, \gamma)$\textit{, i.e.} $u(\alpha')\beta = u(\alpha')\gamma$.
\end{itemize}
\textit{et ces conditions F 1, F 2 impliquent que les $i\bs I'$ sont faiblement filtrantes, et $I$ aussi} (\footnote[57]{Si $I'$ est localement faiblement filtrante (= préfiltrante), alors pour la cofinalité de $u$, il suffit qu'on ait F 0, F 1, F 2, et alors les $i\bs I'$ sont faiblement filtrantes, et $I$ est préfiltrante tout comme $I'$.}).

\printedpage{30}
% head: A. Grothendieck : Dérivateurs, Ch. II, Édition des Universités de Montpellier II et Paris VII
% foot: Version 8 juillet 2001
\item[\textit{c)}] \textit{Si $I'$ est filtrante, alors $u$ est cofinal} si et seulement s'\textit{il satisfait les deux conditions F1, F2, ou encore si et seulement si les $i\bs I'$ sont} filtrantes \textit{(et non seulement} 0-\textit{connexes). Et cela implique que $I'$ est également filtrante} (\footnote[58]{Si $I'$ est localement filtrante, $u$ est cofinal si et seulement si elle satisfait F 0, F 1, F 2, i.e. si et seulement si les $i\bs I'$ sont filtrantes, et alors $I$ est localement filtrante.}).
\item[\textit{d)}] \textit{Si $I$ est préfiltrante (= localement faiblement filtrante, i.e. PS 1), et $u$} pleinement fidèle, \textit{pour que $u$ soit cofinal, il} faut \textit{et il} suffit \textit{que $u$ satisfasse la condition F 1 plus haut. Et alors $I'$ est également préfiltrante, et les $i\bs I'$ sont faiblement filtrantes. Si de plus $I$ est filtrante, les $i\bs I'$ et $I'$ le sont également. (Comparer SGA 4, I 8.1.3.)}
\end{itemize}

\msp{page 47}
\textsc{Démonstration.} a) F 1 est évident, car 0-connexe $\Longrightarrow$ non vide. F 0 car $W_0$-coasphérique $\Longrightarrow$ $W_0$-équivalence, i.e. $\pi_0(u)$ bijectif.

Pour \textit{prouver} b), il suffit de prouver que F 1, F 2 (si $I'$ est faiblement filtrante) impliquent que les $i\bs I'$ sont faiblement filtrantes (a fortiori 0-connexes) et de plus que $I$ l'est également. On sait déjà qu'ils sont non vides par F 1, reste à voir que si $(i', \alpha)$, $(j', \beta)$ sont dans $i\bs I'$, alors ils sont `majorés' par un troisième objet de $i\bs I'$. Comme $I'$ est faiblement filtrante, on trouve
\begin{tikzcd}[row sep=0pt, column sep=small, baseline=(k.base)]
i' \arrow[dr, "r"] & \\
& |[alias=k]| k' \\
j' \arrow[ur, "s"'] &
\end{tikzcd}, d'où $u(r)$, $u(s)$ dans $I$ (cf. diagramme)
\[
\begin{tikzcd}[row sep=small, column sep=small]
& i' \arrow[rr, "r"] \arrow[dd, dash] & & k' \arrow[r, dotted, "t"] \arrow[dd, dash] & l' \arrow[dd, dash] \\
& & j' \arrow[ur, "s"] \arrow[dd, dash] & & \\
& u(i') \arrow[rr, "u(r)"] & & u(k') \arrow[r, dotted, "u(t)"] & u(l') \\
i \arrow[ur, "\alpha"] \arrow[rr, "\beta"'] & & u(j'), \arrow[ur, "u(s)"'] & &
\end{tikzcd}
\]
appliquant F 2 à la double flèche $u(r)\alpha, u(s)\beta : i \rightrightarrows u(k')$, on trouve $k' \xrightarrow{t} l'$ (en pointillés sur \texttt{[le]} diagramme) telle que $(u(t)u(r))\alpha = (u(t)u(s))\beta$, soit $\gamma$, ce qui montre que l'objet $(l', \gamma : i \longrightarrow u(l'))$ majore $(i', \alpha)$ et $(j', \beta)$.

La démonstration que si $I'$ est faiblement filtrante et $u$ satisfait F 1, alors $I$ est faiblement filtrante, est immédiate.

\textit{Prouvons} c). Il reste à prouver que si $I'$ est \textit{filtrante} (non seulement faiblement filtrante), alors si $u$ \texttt{[est]} cofinal, alors F 2 est vérifiée, et de plus les $i\bs I'$ et $I$ sont filtrantes (pas seulement faiblement filtrantes). Mais notons tout de suite que $i\bs I' \longrightarrow I'$ est coétale, donc $I'$ filtrante, donc localement filtrante implique que $i\bs I'$ est aussi localement filtrante\,; comme par hypothèse les $i\bs I'$ sont 0-connexes, ils sont donc filtrantes. La condition F 2 sort facilement à présent :

\printedpage{31}
% head: A. Grothendieck : Dérivateurs, Ch. II, Édition des Universités de Montpellier II et Paris VII
% foot: Version 8 juillet 2001
\msp{page 48}
\[
\begin{tikzcd}[row sep=small, column sep=small]
& i' \arrow[rr, "r"] \arrow[dd, dash] & & j' \arrow[r, dotted, "t"] \arrow[dd, dash] & k' \arrow[dd, dash] \\
& & i' \arrow[ur, "s"] \arrow[dd, dash] & & \\
& u(i') \arrow[rr, "u(r)"] & & u(j') \arrow[r, dotted, "u(t)"] & u(k') \\
i \arrow[ur, "\alpha"] \arrow[rr, "\beta"'] & & u(i'). \arrow[ur, "u(s)"'] & &
\end{tikzcd}
\]
Comme $i\bs I'$ est filtrante, donc faiblement filtrante (PS 1), on peut trouver $i' \dbl{r}{s} j'$ telles que
\[
u(r)\alpha = u(s)\beta,
\]
et comme $I'$ est filtrante, on peut trouver $t : j' \longrightarrow k'$ qui égalise la double flèche $(r, s)$ :
\[
tr = ts(\defeq \tau),
\]
alors
\[
\underbrace{u(t)u(r)}_{u(tr)}\alpha = \underbrace{u(t)u(s)}_{u(ts)}\beta \qquad\quad \text{i.e.} \qquad\quad u(\tau)\alpha = u(\tau)\beta,
\]
i.e. on a vérifié F 2.

Reste à prouver que $I$ est filtrante (si $I'$ l'est et \texttt{[si]} $u : I' \longrightarrow I$ est cofinal). On sait déjà que $I$ est faiblement filtrante par b), reste à vérifier la condition de la double flèche
\[
i \dbl{\alpha}{\beta} j
\]
(PS 2). Soit $i'$ dans $I'$ et $j \longrightarrow u(i')$, et considérons
\[
i \dbl{\gamma\alpha}{\gamma\beta} u(i'),
\]
par hypothèse F 2 $\exists\, r : i' \longrightarrow j'$ dans $I'$ telle que
\msp{page 49}
$u(r)(\gamma\alpha) = u(r)(\gamma\beta)$, donc $u(r)\gamma$ égalise $(\alpha, \beta)$, q.e.d.

La vérification de d) par AQT \texttt{[abréviation dont le sens échappe aux éditeurs]}.

\dgrblock{dgrb114}{The filtered criterion and eventual lifting}
\begin{dcomment}{development}{p49:La vérification de d) par AQT}{The filtered criterion and eventual lifting}
Corollary~3.4 converts connectedness of comma categories into the two concrete operations needed for filtered colimits: reaching the image of $u$, and equalizing parallel arrows after reaching it. The proof of parts~(b) and~(c) is present on the page. There is one harmless slip in the statement of~(c): its last assertion concerns filteredness of $I$, not the already assumed filteredness of $I'$. Part~(d) is left under the abbreviation AQT. Its proof uses full faithfulness in two different places.

One variance qualification is essential. Definition~2.2 called the completion of a cospan by a common predecessor PS~1 and explicitly called its right-hand dual ``colocalement faiblement filtrante.'' Here part~(d) and footnote~57 instead use ``pr\'efiltrante'' or ``localement faiblement filtrante'' for the right-hand condition: a span must admit a commuting square with a common successor. The proof requires that interpretation. Under the earlier left-hand reading, part~(d) would be false: the poset $I=(0\longrightarrow1,\,0\longrightarrow2)$ satisfies left PS~1, and the full inclusion $\{1,2\}\hookrightarrow I$ satisfies F~1, but its comma category under $0$ is a discrete two-point category. Thus this is a necessary clarification of the printed terminology, not a variance-neutral use of PS~1.

Assume that $I$ satisfies the right-hand Ore condition PS~1, that $u:I'\longrightarrow I$ is fully faithful, and that F~1 holds. Given objects $(a,\alpha:i\longrightarrow u a)$ and $(b,\beta:i\longrightarrow u b)$ of $i\backslash I'$, PS~1 supplies arrows $r:u a\longrightarrow j$ and $s:u b\longrightarrow j$ with $r\alpha=s\beta$. Choose $j\longrightarrow u c$ by F~1. Fullness lifts the two resulting arrows $u a\longrightarrow u c$ and $u b\longrightarrow u c$ to arrows in $I'$, and they give a common successor of the original two objects in $i\backslash I'$. F~1 also makes that comma category nonempty. It is therefore weakly filtered, hence connected, and $u$ is cofinal.

To transfer PS~1 itself to $I'$, start with a span $a\longrightarrow b$, $a\longrightarrow c$ in $I'$. Complete its image in $I$ to a commuting square, reach an object in the image of $u$ by F~1, and lift the two new arrows by fullness. Faithfulness ensures that the lifted square commutes. Thus $I'$ satisfies PS~1. If $I$ is filtered, a parallel pair $r,s:a\rightrightarrows b$ in $I'$ can likewise be equalized in $I$, followed by an arrow to $u c$ and lifted to an equalizer arrow $b\longrightarrow c$ in $I'$. Faithfulness again supplies the equality. Nonemptiness follows from F~1, so $I'$ is filtered. The same reasoning, or the discrete-opfibration projection $i\backslash I'\longrightarrow I'$, makes each comma category filtered.

For locally filtered categories, the additional condition F~0 in footnotes~57--58 ensures that objects whose images lie in the same component already lie in the same source component. The arguments can then be carried out in that component. Without this component condition, F~1 and F~2 do not prevent several disconnected copies of one indexing category from being mapped onto it. This is why the global filtered statement and its local variant have different hypotheses.

Footnote~55 asks for a statement encompassing this corollary and Chapter~XVIII. The exact later result is Lemma~4.9.6.3, part~1, manuscript p.~224, whose proof is in \S4.9.6.7. Its additional condition is eventual lifting:
\[
 \text{F~3:}\quad
 \theta:u(a)\longrightarrow u(b)
 \ \Longrightarrow\
 \exists\,v:b\longrightarrow c,\ w:a\longrightarrow c,
 \qquad u(v)\theta=u(w).
\]
For a faithful functor into a filtered category, the lemma identifies F~1 and F~3 with filteredness of the source together with cofinality \citep[Lemma~4.9.6.3(1), pp.~81--82]{Grothendieck1990XVIII}. The following finite formulation combines the two arguments, retaining only the right Ore hypothesis where that suffices. It is our proof-level synthesis of the cited passages, without a claim of priority.

\textbf{A finite comparison with eventual lifting.}\ours
\textit{If the target satisfies right PS~1, F~1 and F~3 make $u$ cofinal. If $u$ is faithful, its source also satisfies right PS~1 and is filtered whenever the target is filtered.}

Indeed, given $(a,\alpha)$ and $(b,\beta)$ under $i$, right PS~1 followed by F~1 produces $\rho_a:u(a)\to u(c)$ and $\rho_b:u(b)\to u(c)$ with $\rho_a\alpha=\rho_b\beta$. Apply F~3 first to $\rho_a$, obtaining $v:c\to d$, $w:a\to d$, and then to $u(v)\rho_b$, obtaining $t:d\to e$, $z:b\to e$. The equality
\[
 u(tw)\alpha=u(t)u(v)\rho_a\alpha
              =u(t)u(v)\rho_b\beta=u(z)\beta
\]
gives a common successor in the comma category. F~1 gives nonemptiness, hence cofinality. For a span in the source, this construction makes its two composites equal after applying $u$; faithfulness makes them equal before applying it. If the target is filtered, equalize the images of any parallel pair, reach another image object by F~1, and lift the resulting arrow eventually by F~3. Faithfulness then gives an equalizing arrow in the source. Cofinality supplies connectedness and nonemptiness, completing filteredness. Conversely, when the source is filtered and $u$ cofinal, the already proved filteredness of $u(a)\backslash I'$ gives F~3 by taking a common successor of $(a,1)$ and $(b,\theta)$.

Full faithfulness makes F~3 automatic, recovering part~(d). The source-filtered case is the F~1--F~2 criterion in part~(c). Thus the finite comparison requested in footnote~55 is explicit. Part~2 of the later lemma strengthens F~3 to simultaneous lifting for families of bounded cardinality and imposes the corresponding larger filtering hypotheses; that cardinal refinement is a separate assertion, not a consequence of the finite conditions alone.
\end{dcomment}


\printedpage{32}
% head: A. Grothendieck : Dérivateurs, Ch. II, Édition des Universités de Montpellier II et Paris VII
% foot: Version 8 juillet 2001
\section*{4\quad Foncteurs biasphériques}

Soit encore
\[
u : I' \longrightarrow I
\]
une flèche de Cat, soit $\cM$ une catégorie, d'où
\[
u^* : \cM(I) \longrightarrow \cM(I');
\tag*{(4.1)}
\]
nous cherchons des conditions sur $u$, $\cM$ pour que ce foncteur soit pleinement fidèle. Si le foncteur adjoint à gauche $u_!$ (resp. adjoint à droite $u_*$) existe, on sait qu'il revient au même de dire que le morphisme d'adjonction
\[
u_!u^*F \longrightarrow F
\tag*{(4.2)}
\]
resp.
\[
u_*u^*F \longleftarrow F
\tag*{(4.2 bis)}
\]
soit \texttt{[un]} isomorphisme pour tout $F$ dans $\cM(I)$. Disons alors que $u$ est $\cM$-connexe, ou $\cM$-biasphérique, cette terminologie sous-entendant que c'est là une propriété qui se réfère au dérivateur $\mathbf{D}_{\cM} = (S \longmapsto \cM(S))$ sur Cat défini par $\cM$. Notons que
\[
\begin{aligned}
&\text{Si $u$ est $\cM$-connexe, alors il est aussi $\cM(S)$-connexe, pour tout $S$ dans Cat,}\\
&\text{puisque $\cM(S)(I) \simeq \cM(I)(S) \defeq \uHom(S^o, \cM(I))$,}
\end{aligned}
\tag*{(4.3)}
\]

\msp{page 50}
et
\[
\text{Si $\cM' \hookrightarrow \cM$ est pleinement fidèle, alors $\cM$-connexe implique $\cM'$-connexe,}
\tag*{(4.4)}
\]
enfin, inversement,
\[
\begin{aligned}
&\text{pour que $u$ soit $\cM$-connexe, il faut et il suffit que pour toute \textit{petite}}\\
&\text{sous-catégorie pleine $\cM'$ de $\cM$, $u$ soit $\cM'$-connexe.}
\end{aligned}
\tag*{(4.5)}
\]
Comme toute petite $\cM'$ se plonge dans $\cM'^\wedge = \operatorname{Ens}(\cM')$, il résulte de (4.3), (4.4), (4.5) que
\[
\text{Si $u$ est Ens-connexe, il est $\cM$-connexe pour toute catégorie $\cM$,}
\tag*{(4.6)}
\]
car c'est le cas si $\cM$ est petite par (4.3), donc pour tout $\cM$ par (4.5).

Disons qu'un foncteur $u$ dans Cat est $\cM$-\textit{cofinal} ou $\cM$-coasphérique si l'homomorphisme fonctoriel
\[
\varprojlim_{I} F \longrightarrow \varprojlim_{I'} u^*(F)
\tag*{(4.7)}
\]
(flèche dans $\uHom(\cM(I), \cM)$ si les $\varprojlim$ de types $I$, $I'$ existent dans $\cM$, sinon flèche dans $\uHom(\cM(I), \cM^\wedge)$) est un isomorphisme. En se ramenant au cas de petites catégories $\cM$, chacune plongée dans $\cM^\wedge$, on voit par le critère (4.2 bis) que
\[
\text{$u$ est $\cM$-connexe si pour tout $i \in \Ob I$,}
\tag*{(4.8)}
\]

\printedpage{33}
% head: A. Grothendieck : Dérivateurs, Ch. II, Édition des Universités de Montpellier II et Paris VII
% foot: Version 8 juillet 2001
\msp{page 51}
le foncteur induit $I'/i \longrightarrow I/i$ est $\cM$-cofinal. La réciproque n'est pas claire, car $u$ est $\cM$-connexe si et seulement si pour tout $i \in \Ob I$ et tout $\cM$-préfaisceau $F_i$ sur $I/i$ \textit{induit par un préfaisceau $F$ sur $I$,} l'homomorphisme $\varprojlim\limits_{I/i} F_i \longrightarrow \varprojlim\limits_{I'/i} (u/i)^*(F/i)$ est \texttt{[un]} isomorphisme. On ne peut donc affirmer d'emblée que cela implique itou pour \textit{tout} $F_i$ sur $I/i$, que si on sait que $(u/i)^* : \cM(I) \longrightarrow \cM(I/i)$ est essentiellement surjectif. Il en est ainsi pour tout $i$ quand $I$ est \textit{ordonné,} car alors $I/i \longrightarrow I$ est pleinement fidèle, donc $\cM(I/i) \longrightarrow \cM(I)$ \texttt{[est]} pleinement fidèle, donc $\cM(I) \longrightarrow \cM(I/i)$ \texttt{[est]} un foncteur de localisation, en particulier essentiellement surjectif.

\msp{page 52}
Mais quand $\cM = (\text{Ens})$, cette perplexité disparaît. On trouve la

\dgrblock{dgrb115}{The missing pointwise limit hypothesis}
\begin{dcomment}{exposition}{p51}{The missing pointwise limit hypothesis}
Grothendieck isolates a real difference between testing every coefficient on a slice and testing only coefficients restricted from the ambient category. There is also a second issue here: full faithfulness for coefficients in $\mathcal M$ does not imply full faithfulness after replacing $\mathcal M$ by its presheaf completion. The sufficient implication (4.8) is valid. The converse discussed on page~51, including its asserted ordered case, requires hypotheses ensuring that the relevant pointwise Kan-extension limits exist in $\mathcal M$.

Take $I=[1]=(0\longrightarrow1)$, $I'=\{1\}$, and let $\mathcal M$ be the discrete category on two objects $m,n$. Every $\mathcal M$-valued presheaf on $I$ is constant, so restriction $\mathcal M(I)\longrightarrow\mathcal M(I')$ is an equivalence. Nevertheless, at $i=0$ the induced functor on slices is $\varnothing\longrightarrow *$. In $\mathcal M^{\wedge}\simeq\mathrm{Set}\times\mathrm{Set}$, the Yoneda image of $m$ is $(1,\varnothing)$, whereas the limit of the empty diagram is $(1,1)$. Thus the comparison
\[
 \varprojlim_{I/0}y(m)=(1,\varnothing)
 \longrightarrow\varprojlim_{I'/0}y(m)=(1,1)
\]
is not invertible. This coefficient on $I/0$ even extends to $I$. Hence the intermediate assertion that $\mathcal M$-connectedness is equivalent to the slice comparison for restricted coefficients fails without pointwise limit hypotheses.

Existence of an abstract adjoint alone is insufficient: in this example restriction is an equivalence and has both adjoints. What fails is the formula that would compute its right adjoint by an empty limit inside $\mathcal M$. Completeness of $\mathcal M$, or existence of all the particular comma limits involved, supplies the needed repair.

The extension observation for posets is independently correct. Identify $I/i$ with $A=\{x\mid x\leq i\}$. The map $r:I\longrightarrow A$ sending $x$ to itself when $x\leq i$ and to $i$ otherwise is order preserving: no arrow can enter $A$ from outside, and an arrow leaving $A$ is sent to $x\leq i$. With $j:A\hookrightarrow I$, one has $rj=1_A$, so every coefficient on $A$ extends by $r^{\mathrm{op}}$. Under the pointwise limit hypotheses this proves the ordered-case converse. It does not remove those hypotheses. Proposition~4.1 subsequently quantifies over all coefficient categories and can use sets from the start; its intended theorem is unaffected by this coefficient-specific error.
\end{dcomment}


\noindent\textbf{Proposition 4.1.} \textit{Soit $u : I' \longrightarrow I$ une flèche dans} Cat. \textit{Conditions équivalentes :}
\begin{itemize}[label={}]
\item[\textit{a)}] \textit{Pour toute catégorie $\cM$, $u^*_{\cM} : \cM(I) \longrightarrow \cM(I')$ est pleinement fidèle (i.e. $u$ est $\cM$-connexe).}
\item[\textit{a$'$)}] \textit{Pour toute catégorie $\cM$ où les limites inductives de type $I$, $I'$ existent, la flèche d'adjonction pour $F$ dans $\cM(I)$} % Transcription corrected: official printed33 reads a prime.
\[
u_!u^*F \xrightarrow{\sim} F
\]
\textit{est} \texttt{[un]} \textit{isomorphisme.}
\item[\textit{a$''$)}] \textit{Dual de a$'$), avec} \texttt{[la]} \textit{flèche d'adjonction}
\[
u_*u^* \xleftarrow{\sim} F. % [sic] u_*u^* -> u_*u^*F, the object F is missing on the left side (compare a') and (4.2 bis))
\]
\item[\textit{b) b$'$) b$''$)}] \textit{Comme a) a$'$) a$''$), avec $\cM =$} (Ens).
\item[\textit{c)}] \textit{Comme b$'$), avec $F$ représentable, i.e. de la forme $j \longmapsto$} $\Hom_I(j, i)$.
\item[\textit{d)}] \textit{Pour tout $I$ dans $I$, $I'/i \longrightarrow I/i$ est cofinal.} % [sic] Pour tout I dans I -> Pour tout i dans I, the variable of the statement is i (compare d'))
\item[\textit{d$'$)}] \textit{Pour tout $i$ dans $I$, $i\bs I' \longrightarrow i\bs I$ est cofinal.}
\item[\textit{d$''$)}] \textit{Pour toute flèche $i_0 \xrightarrow{\alpha} i_1$ dans $I$, la catégorie} $i_0\bs \underset{\alpha}{I'}/i_1 \defeq I' \times_I \underbrace{i_0\bs \underset{\alpha}{I}/i_1}_{\substack{\text{catégorie biïnduite}\\ (i_0, \alpha)\bs(I/i_1) \simeq (i_0\bs I)/(i_1, \alpha)}}$ \qquad \textit{est}

0-\textit{connexe.}
\end{itemize}

\msp{page 53}
\textsc{Démonstration.} On sait déjà, par les préliminaires, l'équivalence de a, a$'$, a$''$, b, b$'$, b$''$, soit A cette condition - nous dirons que $u$ est $W_0$-\textit{connexe} ou $W_0$-\textit{biasphérique,} si ces conditions équivalentes sont satisfaites. On a de plus l'autodualité
\[
\text{$u$ est $W_0$-connexe si et seulement si $u^o$ l'est,}
\tag*{(4.8)} % [sic] (4.8) -> another number, the label (4.8) is already used for a display on p. 32
\]
car on a évidemment
\[
\text{$u$ est $\cM$-connexe si et seulement si $u^o$ est $\cM^o$-connexe.}
\tag*{(4.9)}
\]


\dgrblock{dgrb116}{Biasphericity, lax epimorphisms, and factorization categories}
\begin{dcomment}{later}{p53}{Biasphericity, lax epimorphisms, and factorization categories}
The new question in Section~4 is whether a diagram and all its natural transformations can be recovered after restriction. The answer is governed by factorizations of arrows, just as ordinary cofinality was governed by objects of comma categories. This is the ordinary version of the factorization test encountered in the commentary on Chapter~I, manuscript p.~49 (``The factorization category controls base change''). The proof of Proposition~4.1 continues beyond these pages; we give the complete argument for its corrected statement here.

For $u:A\longrightarrow B$ and $\gamma:b_0\longrightarrow b_1$, let $\mathcal F_u(\gamma)$ have objects
\[
 (a,p,q),\qquad b_0\xrightarrow{p}u a\xrightarrow{q}b_1,
 \quad qp=\gamma,
\]
and arrows $h:a\longrightarrow a'$ satisfying $u(h)p=p'$ and $q'u(h)=q$. Then the following conditions are equivalent: restriction on presheaves is fully faithful for every coefficient category; it is fully faithful for sets; and every $\mathcal F_u(\gamma)$ is nonempty and connected.

To prove this, first work with sets, where both Kan extensions exist. Full faithfulness of $u^*$ is equivalent to the counit $u_!u^*F\longrightarrow F$ being an isomorphism. Both sides preserve colimits in $F$, and every presheaf is a colimit of representables, so it suffices to take $F=h_{b_1}=B(-,b_1)$. At $b_0$ the pointwise colimit formula gives
\[
 (u_!u^*h_{b_1})(b_0)
 \cong\coprod_{\gamma\in B(b_0,b_1)}
             \pi_0\mathcal F_u(\gamma).
\]
Indeed, the generators are pairs $b_0\longrightarrow u a\longrightarrow b_1$, and the colimit relations are exactly the zigzags induced by arrows of $\mathcal F_u(\gamma)$. Composition labels every such relation by the same $\gamma$. The counit sends a connected component to that composite. It is bijective if and only if each indicated category is nonempty and connected.

The set-valued criterion implies the version for arbitrary coefficients. It first implies full faithfulness for coefficients in every presheaf category, by applying the set-valued statement in each variable. A small category $\mathcal M$ embeds fully faithfully into its presheaf category by Yoneda; the full-faithfulness property descends to this full subcategory. For a possibly large locally small $\mathcal M$, any two diagrams on the small category $B$ take values in a small full subcategory. Applying the preceding argument there descends every natural transformation and proves uniqueness. This establishes the reductions (4.3)--(4.6), with the usual universe enlargement when necessary.

For fixed $b_1$, the comma category of the induced functor $A/b_1\longrightarrow B/b_1$ under $(b_0,\gamma)$ is $\mathcal F_u(\gamma)$. Thus condition~(d) is equivalent to the factorization criterion. Its dual says that
\[
 b_0\backslash A\longrightarrow b_0\backslash B
 \quad\text{is \emph{co-cofinal}, that is, initial.}
\]
The edition instead prints ``cofinal'' in~(d$'$). That cannot be retained as a mathematical equivalence. The inclusion $u:\{1\}\hookrightarrow[1]$ satisfies the printed condition: for $i=0$ it induces the final-vertex inclusion into $[1]$, and for $i=1$ it induces the identity of a point. Yet restriction on presheaves is not faithful: take $F(1)=\varnothing$, $F(0)=\{0,1\}$; its identity and the automorphism transposing $0,1$ become the same map after restriction. The corrected dual condition has comma categories opposite to $\mathcal F_u(\gamma)$, so it gives exactly the required connectedness. The two adjunction formulations and self-duality now follow, proving all the intended equivalences.

In later terminology these functors are \emph{lax epimorphisms}, also called absolutely dense functors. \citet{AdamekElBashirSobralVelebil2001} identify them by full faithfulness of set-valued restriction and by connected categories of factorizations; the latter characterization is restated explicitly in \citet[Theorem~4.1]{NunesSousa2022}. For derivators, \citet[Theorem~6.10]{GrothStovicek2016Trees} replace connectedness by weak contractibility of the same factorization categories and obtain \emph{homotopical epimorphisms}. Thus the arrowwise obstruction persists; the homotopy theory changes the required strength of its vanishing.

A quotient group displays the difference concretely. For a surjective homomorphism $\rho:H\twoheadrightarrow G$, restriction of $G$-sets to $H$-sets is fully faithful, because equivariance under $H$ is equivalent to equivariance under its image $G$. For any $g\in G$, the category $\mathcal F_{B\rho}(g)$ is the transitive action groupoid with object set $G$, on which $H$ acts through $\rho$; its stabilizer is $\ker\rho$. It is therefore equivalent to $B\ker\rho$. The functor $B\rho$ is an ordinary lax epimorphism, but it is a homotopical epimorphism only when $\ker\rho$ is trivial. With $\rho:\mathbf Z\twoheadrightarrow\mathbf Z/n\mathbf Z$, the obstruction has the homotopy type of a circle. The full-faithfulness statement about ordinary diagrams is correct; it simply does not recover that circle.
\end{dcomment}
\begin{otherlanguage}{english}
\dgroverview{II-a}{Chapter II \textperiodcentered\ \textup{II-a} \textperiodcentered\ Overview}
\begingroup
\selectlanguage{english}
\noindent\textbf{Cofinality, connected objects, and the recovery of diagrams.}

Chapter~II begins by asking when a change of indexing category leaves sections unchanged \citep{Grothendieck1990II}. It ends this first part with the stronger question of when restriction recovers a whole diagram and every natural transformation between diagrams. Between these questions Grothendieck develops a topos-theoretic account of filtering, examines the homotopy information lost by ordinary connectedness, and specializes the language of derivators to sets. The movement is deliberate. Ordinary cofinality is not merely an elementary preliminary to the derived theory: its proof identifies the object which must be derived, and its limitations identify the higher information which a derivator must retain.

The conventions are essential. A coefficient on $C$ is a contravariant functor, and the category denoted $c\backslash C'$ has objects $c\to\alpha(c')$. Thus a cofinal functor in these pages is final in the usual covariant terminology; its opposite is initial for the diagrams whose limits are being taken. Under this convention the first question has a representability answer. The presheaf
\[
 K_\alpha(c)=\pi_0(c\backslash C')= (\alpha_!1_{C'})(c)
\]
represents $F\mapsto\Gamma(C',\alpha^*F)$. The comparison with $\Gamma(C,F)$ is induced by $K_\alpha\to1_C$. It is an isomorphism for all $F$ exactly when every indicated comma category is nonempty and connected. This is the classical cofinality theorem of SGA~4, Expos\'e~I, \S8.1 \citep{SGA4I1972}. The initial weak cofiltering hypothesis is sufficient, and Grothendieck replaces it by the exact connectedness criterion as the argument develops.

The same proof suggests its own derived form. For a commutative ring $k$, replace the component presheaf by the derived weight $\mathbf L\alpha_!k_{C'}$. Its value at $c$ is represented by chains on $N(c\backslash C')$. Derived adjunction then gives
\[
 \mathbf R\Gamma(C',\alpha^*F)
 \simeq\mathbf R\operatorname{Hom}_{C;k}
       (\mathbf L\alpha_!k_{C'},F).
\]
Consequently weak contractibility of all the comma nerves implies cohomological cofinality, answering the question on manuscript page~5 affirmatively. For these fixed linear coefficients, their $k$-acyclicity already suffices. For spaces the weight is the corresponding presheaf of classifying spaces, and the mapping-space adjunction gives the homotopical statement. Modern higher category theory places this argument within the general criterion for cofinality \citep[Theorem~4.1.3.1]{Lurie2009HTT}. The commentary constructs the linear comparison by a bar resolution, so the answer does not depend on the internal-Hom formula subsequently proposed in the manuscript.

That formula is a genuine point of failure. The global derived Hom belongs to the category of coefficients at a point; a sheaf-valued internal Hom belongs to the diagram category. Applying global sections can identify two such objects without identifying the objects themselves. For the terminal-vertex inclusion $i:\{1\}\hookrightarrow[1]$, a presheaf is an arrow $F_1\to F_0$, and
\[
 i_!A=(A\xrightarrow{1}A),\qquad i_*A=(A\to *).
\]
The comma categories for cofinality are points, and global sections are $F_1$. Nevertheless $i_*i^*F$ need not be $F$, while $\underline{\operatorname{Hom}}(i_!1,F)=F$. This disproves the unrestricted internal identification in sets, hence already before derived complications arise. The correct internal comparison is equivalent to an additional projection-formula property, whose canonical arrow has the opposite direction from the one displayed in the manuscript. Our point-inclusion calculation identifies its obstruction exactly: for a fixed presheaf, every restriction along an arrow into the chosen object must be invertible. The formula holds for all coefficients precisely when all those incoming arrows are isomorphisms.

The turn to toposes is motivated by a different problem: how to recognize the equalization part of filtering without inspecting individual arrows. An irreducible presheaf topos records common predecessors of objects, a condition visible in the underlying preorder. Local irreducibility records common predecessors compatible with a prescribed pair of arrows, namely PS~1. It still does not supply PS~2, the equalization of arbitrary parallel arrows. The parallel-arrow category $0\rightrightarrows1$ separates the first two conditions. A nontrivial group category $BG$ separates local irreducibility from local cofilteredness: every representable slice is equivalent to a point, but cancellation prevents the equalization of distinct group elements. Several preliminary identifications in the manuscript are explicitly withdrawn by its annotations. The audit preserves this sequence of correction and distinguishes it from errors of transcription.

Total $0$-connectedness adds preservation of products by the component functor. For $\widehat X=\operatorname{Fun}(X^{\mathrm{op}},\mathrm{Set})$, it says that $X^{\mathrm{op}}$ is an ordinary sifted category. Sifted colimits preserve finite products; filtered colimits preserve finite limits, including equalizers \citep{AdamekRosickyVitale2010}. The simplex category supplies the instructive distinction. The products of its representables have connected categories of elements, but the two vertices of $[1]$ admit no common predecessor over $[1]$. Thus its presheaf topos is totally $0$-connected without being locally irreducible. Nothing about this ordinary product test alone makes all the relevant nerves weakly contractible.

The surviving theorem in 2.6 makes the product condition hereditary under connected slices. Its equivalent formulations are (a), (a'), (b), and (d); the generator-slice conditions (c) and (c') are strictly weaker. The diagonal $b\to b\times b$ explains the distinction. Once $b\times b$ is connected and the topos is locally irreducible, its pullback against $(u,v):a\to b\times b$ is noninitial and yields an equalizing arrow. Conversely, PS~1 and PS~2 lift to categories of elements, giving connected products in every connected slice. The commentary supplies the omitted generator-component lemma and the successive equalizations in this argument. It also restores connectedness of the base where the printed formulation of 2.8 omits it.

On a connected topos satisfying this hereditary condition, $\pi_0$ preserves finite limits. This makes the later cohomological argument precise. The induction must run over connected, locally totally $0$-connected toposes, since it repeatedly passes to connected components of slice toposes. The degree-zero row of the covering spectral sequence is the augmented complex of functions $I^{p+1}\to M$, where $I$ is the component set of a covering object. Inserting a chosen element of $I$ contracts that complex. The stronger homotopical conclusion follows from hypercovers: applying the exact component functor to their matching-object diagrams gives boundary fillers in every degree, so the resulting simplicial sets are weakly contractible. This connects the argument to the classical hypercover construction of homotopy types \citep[\S\S8--9]{ArtinMazur1969}. It concerns constant coefficients and shape. It does not trivialize arbitrary sheaf cohomology, as the worked calculation $H^1(\mathbf P^1_k,\mathcal O(-2))\cong k$ illustrates.

The ind-object reformulation in 2.15--2.18 returns to a classical theorem: a presheaf is a filtered colimit of representables exactly when its category of elements is filtered \citep[Expos\'e~I, Theorem~8.3.3]{SGA4I1972}. Taking the terminal presheaf recovers filteredness of the indexing category. Taking its ordinary left Kan extension along an arbitrary functor shows that the resulting category of elements is again filtered. The modern minimal-localizer theorem supplies the asphericity conclusion for every fundamental localizer, without the additional limit axiom used in the historical formulation \citep[Theorem~2.2.11]{Cisinski2004Minimal}.

The converse question on page~36 has a stronger quantifier than later uses of the phrase ``totally aspherical''. In the commentary we denote by $T(A)$ the condition that every ordinary cofinal target of $A$ have weakly contractible nerve. Filtered categories satisfy it. The attempted proof that ordinary siftedness suffices is retracted on pages~36bis--36ter; our explicit completion of the counterexample uses the simplicial circle and a one-point collage whose classifying space is $S^2$. We also prove a necessary common-successor condition for $T(A)$. Its obstruction is a four-object category with a cofinal three-object subcategory and a nonzero integral two-cycle. This settles $T(A)\Leftrightarrow A\text{ filtered}$ for preorders. For general categories the remaining issue is equalization of parallel arrows. We did not locate a resolution under this exact universal quantifier in the sources checked; this is a report of that bounded search, not a declaration of a recognized open problem. Results developed here are marked as the commentators' work and carry no claim of priority.

Section~3 shows how much of the earlier derivator terminology can be tested with sets. Pointwise localization of category-valued diagrams at component equivalences yields the represented derivator of sets. Properness and smoothness become connectedness of strict lift categories. A universal $W_0$-equivalence can be tested over vertices and arrows; a strong $W_0$-fibration is characterized component by component by those universal equivalences. Coverings explain why the weak condition is weaker: monodromy can connect different fibre components without making each fibre connected. The comparison with the commentary on Chapter~I is made at the precise lifting and monodromy arguments, rather than repeating its whole discussion.

The final transition is from preservation of global sections to full faithfulness of restriction. Proposition~4.1 characterizes the latter by nonempty connected categories of factorizations
\[
 b_0\longrightarrow u(a)\longrightarrow b_1
\]
of every arrow $b_0\to b_1$. These are the lax epimorphisms, or absolutely dense functors, of the later literature \citep{AdamekElBashirSobralVelebil2001}. The homotopical version requires weakly contractible factorization categories \citep[Theorem~6.10]{GrothStovicek2016Trees}. A quotient of groups gives a particularly concrete distinction: the factorization category has the homotopy type of the classifying space of the kernel. Two qualifications belong to the audit. The printed coslice condition in 4.1(d') must be initial, not final. Also the coefficient-specific discussion on page~51 invokes pointwise Kan-extension formulas without the necessary comma limits; a discrete coefficient category on two objects gives a counterexample even when the indexing category is ordered. The corrected theorem quantified over all coefficient categories is established by its set-valued case and Yoneda embeddings.

The historical setting is the program of coherent diagram categories described in \citet[\S69]{Grothendieck1983}, alongside the systematic homotopy Kan extensions of \citet{Heller1988}. The commentary keeps the successive formulations, counterexamples, and marginal revisions visible. It distinguishes a proof supplied here from a proof actually written in these pages, and distinguishes a settled theorem from an unsuccessful search for a later treatment.

\endgroup
\end{otherlanguage}
\dgrchunkend
